\documentclass[12pt, a4paper, oneside]{book}

% ── Encoding & language (xelatex) ─────────────────────────────────────────────
\usepackage{fontspec}
\setmainfont{Latin Modern Roman}
\usepackage{polyglossia}
\setmainlanguage{turkish}

% ── Math ───────────────────────────────────────────────────────────────────────
\usepackage{amsmath, amsthm, amssymb}

% ── Theorems ───────────────────────────────────────────────────────────────────
\newtheorem{theorem}{Teorem}[chapter]
\newtheorem{lemma}[theorem]{Lema}
\newtheorem{corollary}[theorem]{Sonuç}
\theoremstyle{definition}
\newtheorem{definition}[theorem]{Tanım}
\newtheorem{example}[theorem]{Örnek}
\theoremstyle{remark}
\newtheorem*{remark}{Not}

% ── Algorithms ─────────────────────────────────────────────────────────────────
\usepackage{algorithm}
\usepackage{algpseudocode}

% ── Code listings ──────────────────────────────────────────────────────────────
\usepackage{listings}
\usepackage{xcolor}
\definecolor{codebg}{rgb}{0.97,0.97,0.97}
\definecolor{codegreen}{rgb}{0.0,0.5,0.0}
\definecolor{codepurple}{rgb}{0.5,0.0,0.5}
\lstset{
  language=Python,
  basicstyle=\ttfamily\small,
  backgroundcolor=\color{codebg},
  commentstyle=\color{codegreen},
  keywordstyle=\color{codepurple}\bfseries,
  stringstyle=\color{orange},
  breaklines=true,
  frame=single,
  framerule=0.4pt,
  numbers=left,
  numberstyle=\tiny\color{gray},
  tabsize=4,
  showstringspaces=false,
}

% ── Tables ─────────────────────────────────────────────────────────────────────
\usepackage{booktabs}
\usepackage{array}

% ── Lists ──────────────────────────────────────────────────────────────────────
\usepackage{enumitem}

% ── Hyperlinks ─────────────────────────────────────────────────────────────────
\usepackage[colorlinks=true, linkcolor=blue, urlcolor=blue, citecolor=blue]{hyperref}

% ── Layout ─────────────────────────────────────────────────────────────────────
\usepackage{geometry}
\geometry{margin=2.5cm}

% ── Title page data ────────────────────────────────────────────────────────────
\title{%
  \Huge\bfseries pytop Kullanım Kılavuzu\\[0.5em]
  \Large Nokta-Küme Topolojisi ve Metrik Uzaylar
}
\author{Kadirhan Polat}
\date{2026}

% ══════════════════════════════════════════════════════════════════════════════
\begin{document}

\maketitle

\frontmatter

\tableofcontents

\chapter*{Ön Söz}
\addcontentsline{toc}{chapter}{Ön Söz}

Bu kılavuz \texttt{pytop} paketinin nokta-küme topolojisi ve metrik uzaylar
modüllerini kapsamlı biçimde tanıtmaktadır.
Her bölüm aynı yapıyı izler:
\begin{enumerate}
  \item Teorik giriş ve formal tanımlar
  \item Temel teoremler (kanıt iskeletleriyle)
  \item Algoritmalar (sözde kod ve karmaşıklık)
  \item \texttt{pytop} API
  \item Çalışan örnekler
  \item Alıştırmalar
\end{enumerate}

Örnekler aynı zamanda \texttt{docs/user\_guide/python/} altındaki Python
script'lerinde ve \texttt{docs/user\_guide/notebook/} altındaki Jupyter
defterlerinde çalıştırılabilir biçimde mevcuttur.

\mainmatter

% ── Ön Koşullar ───────────────────────────────────────────────────────────────
\part{Ön Koşullar}

\chapter{pytop'a Hızlı Giriş}

pytop kurulumu, temel uzay nesneleri, \texttt{Result} tipi ve tag sistemi —
kılavuzun geri kalanını okumadan önce bu bölümü çalıştırın.

%-------------------------------------------------------------------
\section{Kurulum ve Import}

\begin{lstlisting}[language=Python]
pip install -e .   # git kokundan
\end{lstlisting}

Her bolumde ihtiyac duyulan semboller dogrudan \texttt{pytop}'tan import edilir.
Bu bolum bir \emph{tur} oldugundan, ilk hucrede sonraki tum hucrelerin
kullanacagi sembolleri tek seferde ice aktariyoruz:

\begin{lstlisting}[language=Python]
import pytop
from pytop import (
    discrete_topology, indiscrete_topology, sierpinski_space, make_topology,
    is_compact, is_connected, is_t0, is_t1, is_t2,
    count_topologies_on_n_points,
    simplicial_complex, simplicial_homology, betti_numbers,
    FiniteMetricSpace, persistent_homology, persistence_betti_numbers,
)
from pytop.experimental.spaces import finite_circle, pi1_space
print("pytop surumu:", pytop.__version__)
\end{lstlisting}

\begin{verbatim}
pytop surumu: 1.7.0
\end{verbatim}

\begin{sezgi}
pytop'u iki katlı bir bina gibi düşünün. Alt kat \textbf{betimsel}
(descriptive) katmandır: ünlü uzayların ve teoremlerin kayıtlı, kaynaklı
bilgisini tutar --- bilir ama hesaplamaz. Üst kat \textbf{yapıcı}
(constructive) katmandır: ham girdiden invaryantları \emph{hesaplar} ---
homoloji, kalıcı homoloji, düğüm polinomları. Bu bölüm her iki kata da
kısa bir tur attırır.
\end{sezgi}

\begin{figure}[ht]
\centering
\begin{tikzpicture}[font=\small, node distance=6mm,
    kat/.style={draw, rounded corners, minimum width=78mm, minimum height=15mm,
                align=center, inner sep=5pt},
    kut/.style={draw, rounded corners, fill=white, inner sep=3pt,
                align=center, font=\scriptsize}]
  \node[kat, fill=green!8] (desc) {};
  \node[anchor=north west, font=\bfseries] at (desc.north west) {Betimsel katman};
  \node[anchor=south east, font=\scriptsize\itshape] at (desc.south east) {bilir, hesaplamaz};
  \node[kut, below=2mm of desc.north, yshift=-4mm] (d1) {\texttt{*Profile} kayitlari};
  \node[kut, right=4mm of d1] (d2) {\texttt{get\_*\_profiles()}};
  \node[kut, left=4mm of d1] (d0) {unlu uzay/teorem};

  \node[kat, fill=blue!8, below=14mm of desc] (cons) {};
  \node[anchor=north west, font=\bfseries] at (cons.north west) {Yapici katman};
  \node[anchor=south east, font=\scriptsize\itshape] at (cons.south east) {ham girdiden hesaplar};
  \node[kut, below=2mm of cons.north, yshift=-4mm] (c1) {\texttt{simplicial\_homology}};
  \node[kut, right=4mm of c1] (c2) {\texttt{persistent\_homology}};
  \node[kut, left=4mm of c1] (c0) {\texttt{jones\_polynomial}};

  \draw[-{Stealth[length=2.6mm]}, thick]
    (cons.north) -- (desc.south)
    node[midway, right=1mm, font=\scriptsize, align=left]
    {hesaplanan sonuc\\ kayitlari dogrular};
\end{tikzpicture}

\caption{pytop'un iki katmanı: betimsel kayıtlar ile yapıcı motorlar.}
\end{figure}

%-------------------------------------------------------------------
\section{Uzay Nesneleri}

\begin{table}[h]
\centering
\begin{tabular}{lll}
\toprule
\textbf{Kurucu} & \textbf{Topoloji} & \textbf{Sonuç} \\
\midrule
\texttt{discrete\_topology(*pts)} & Her altküme açık & Ayrık \\
\texttt{indiscrete\_topology(*pts)} & Yalnız $\emptyset$ ve $X$ açık & İndiscrete \\
\texttt{sierpinski\_space()} & $\{\emptyset, \{1\}, \{0,1\}\}$ & Sierpiński \\
\texttt{make\_topology(carrier, *open\_sets)} & Kullanıcı tanımlı & Özel \\
\bottomrule
\end{tabular}
\caption{Temel uzay kurucuları}
\end{table}

\begin{lstlisting}[language=Python]
d = discrete_topology(0, 1, 2)
print("carrier:", sorted(d.carrier))
print("topoloji boyutu:", len(d.topology))
print("etiketler:", sorted(d.tags))
\end{lstlisting}

\begin{verbatim}
carrier: [0, 1, 2]
topoloji boyutu: 8
etiketler: ['discrete', 'finite', 'hausdorff', 'metrizable', 'normal', 'regular']
\end{verbatim}

%-------------------------------------------------------------------
\section{Result Tipi}

pytop'un çoğu yüklemi bir \texttt{Result} nesnesi döner.

\begin{lstlisting}[language=Python]
s = sierpinski_space()
r = is_t2(s)
print(".status:", r.status)        # 'true' | 'false' | 'unknown'
print(".value:", r.value)
print(".mode:", r.mode)
print(".justification:", r.justification)
\end{lstlisting}

\begin{verbatim}
.status: false
.value: hausdorff
.mode: exact
.justification: ['The explicit finite topology fails hausdorff.']
\end{verbatim}

\texttt{.status} her zaman \texttt{'true'}, \texttt{'false'} veya \texttt{'unknown'}
dizesidir --- Python bool'u değil.

%-------------------------------------------------------------------
\section{Temel Uzaylar Turu}

\begin{lstlisting}[language=Python]
spaces = {
    "Ayrik D(0,1,2)":     discrete_topology(0, 1, 2),
    "Indiscrete I(0,1,2)": indiscrete_topology(0, 1, 2),
    "Sierpinski S":        sierpinski_space(),
}
for name, sp in spaces.items():
    print(name,
          is_compact(sp).status,
          is_connected(sp).status,
          is_t2(sp).status)
\end{lstlisting}

\begin{verbatim}
Ayrik D(0,1,2)      true  false  true
Indiscrete I(0,1,2) true  true   false
Sierpinski S        true  true   false
\end{verbatim}

%-------------------------------------------------------------------
\section{Özel Topoloji: make\_topology}

\begin{lstlisting}[language=Python]
X = [1, 2, 3, 4]
tau = [set(), {1,2}, {3,4}, {1,2,3,4}]
fts = make_topology(X, *tau)
print("topoloji boyutu:", len(fts.topology))
print("kompakt:", is_compact(fts).status)
print("T2:", is_t2(fts).status)
\end{lstlisting}

\begin{verbatim}
topoloji boyutu: 4
kompakt: true
T2: false
\end{verbatim}

%-------------------------------------------------------------------
\section{Tag Sistemi}

\begin{lstlisting}[language=Python]
print("Ayrik:    ", sorted(discrete_topology(0,1,2).tags))
print("Indiscrete:", sorted(indiscrete_topology(0,1,2).tags))
print("Sierpinski:", sorted(sierpinski_space().tags))
\end{lstlisting}

\begin{verbatim}
Ayrik:      ['discrete', 'finite', 'hausdorff', 'metrizable', 'normal', 'regular']
Indiscrete: ['compact', 'connected', 'finite', 'indiscrete']
Sierpinski: ['compact', 'connected', 'finite', 't0']
\end{verbatim}

Etiketler yüklem fonksiyonlarının karar mekanizmasını hızlandırır:
\texttt{is\_t2(d)} açık küme taraması yerine \texttt{'hausdorff' in d.tags}
kontrolüyle sonuç verir.

%-------------------------------------------------------------------
\section{n Nokta Üzerindeki Topoloji Sayısı}

\begin{lstlisting}[language=Python]
for n in range(1, 6):
    print(f"n={n}: {count_topologies_on_n_points(n)}")
\end{lstlisting}

\begin{verbatim}
n=1: 1
n=2: 4
n=3: 29
n=4: 355
n=5: 6942
\end{verbatim}

Topoloji sayısı hızla büyür; sonlu uzayların tüm topolojileri üzerinde kaba
kuvvet taraması yalnızca küçük $n$ için pratiktir.

%-------------------------------------------------------------------
\section{pytop Ne Hesaplar? — Genel Harita}

Önceki bölümler \textbf{betimsel/sonlu} katmanı gezdirdi. pytop'un asıl gücü
\textbf{yapıcı çekirdektir}: ham bir kompleks ya da nokta bulutundan cebirsel
invaryantları gerçekten \emph{hesaplar}.

\begin{figure}[ht]
\centering
\begin{tikzpicture}[font=\small, node distance=10mm,
    gir/.style={draw, rounded corners, fill=orange!12, align=center, inner sep=5pt},
    inv/.style={draw, rounded corners, fill=blue!8, align=center, inner sep=5pt}]
  \node[gir] (pts) {nokta bulutu};
  \node[gir, below=8mm of pts] (cpx) {simpleks kompleks};
  \node[gir, below=8mm of cpx] (knot) {dugum diyagrami};

  \node[inv, right=34mm of pts] (ph) {kalici homoloji\\ \scriptsize barkod, $H_*$};
  \node[inv, right=34mm of cpx] (hom) {homoloji\\ \scriptsize Betti, torsiyon};
  \node[inv, right=34mm of knot] (jp) {Jones / Alexander\\ \scriptsize polinom};

  \draw[-{Stealth[length=2.6mm]}, thick] (pts)  -- (ph)
        node[midway, above, font=\scriptsize] {Vietoris--Rips};
  \draw[-{Stealth[length=2.6mm]}, thick] (cpx)  -- (hom)
        node[midway, above, font=\scriptsize] {sinir matrisi $\to$ SNF};
  \draw[-{Stealth[length=2.6mm]}, thick] (knot) -- (jp)
        node[midway, above, font=\scriptsize] {Kauffman / Burau};

  \draw[-{Stealth[length=2.2mm]}, thick, gray] (cpx.north) to[bend left=18] (ph.south west);
\end{tikzpicture}

\caption{pytop hesaplama haritası: girdiden invaryanta giden yollar.}
\end{figure}

%-------------------------------------------------------------------
\section{Hesaplamalı Çekirdek: İlk Bakış}

\subsection*{Örnek 9.1 — Sonlu Bir Çember Modeli ve $\pi_1 = \mathbb{Z}$}

\begin{lstlisting}[language=Python]
fc = finite_circle()
r = pi1_space(fc)
print("uzay adi   :", r.space_name)
print("grup tipi  :", r.group_type)
print("uretec say :", len(r.presentation.generators))
print("baginti say:", len(r.presentation.relators))
print("serbest mi :", r.is_free())
print("H1 betti   :", r.abelianization_betti)
\end{lstlisting}

\begin{verbatim}
uzay adi   : S^1
grup tipi  : infinite_cyclic
uretec say : 1
baginti say: 0
serbest mi : True
H1 betti   : 1
\end{verbatim}

Yalnız 4 nokta, hiç bağıntısı olmayan tek üreteç: $\pi_1 = \langle a \rangle
= \mathbb{Z}$. Sonlu bir uzay, sürekli bir çemberle aynı temel grubu taşır.

\subsection*{Örnek 9.2 — Üçgen Kenarından Çemberin Homolojisi}

\begin{lstlisting}[language=Python]
circle = simplicial_complex([[0], [1], [2], [0, 1], [1, 2], [0, 2]])
for k in (0, 1):
    h = simplicial_homology(circle, k)
    print(f"H{k}: betti={h.betti} torsion={h.torsion}")
print("betti_numbers:", betti_numbers(circle))
\end{lstlisting}

\begin{verbatim}
H0: betti=1 torsion=()
H1: betti=1 torsion=()
betti_numbers: (1, 1)
\end{verbatim}

$H_0 = \mathbb{Z}$ (tek bileşen), $H_1 = \mathbb{Z}$ (bir delik) --- çemberin
klasik Betti imzası $(1, 1)$, ham simpleks listesinden SNF yoluyla hesaplandı.

\begin{dikkat}
\texttt{simplicial\_complex} yüz-kapalı girdi bekler: kenarları
\texttt{[0,1]} verip köşeleri \texttt{[0]} eklemeyi unutursanız hata alırsınız.
Üçgeni \texttt{[0,1,2]} olarak eklerseniz içi dolar ve $H_1 = 0$ olur.
\end{dikkat}

\subsection*{Örnek 9.3 — Nokta Bulutundan Kalıcı Homoloji}

\begin{lstlisting}[language=Python]
import math

pts = tuple(
    (round(math.cos(2 * math.pi * k / 8), 4),
     round(math.sin(2 * math.pi * k / 8), 4))
    for k in range(8)
)
space = FiniteMetricSpace(carrier=pts, distance=math.dist)
pairs = persistent_homology(space, max_dimension=1, max_scale=1.0)

essential = [p for p in pairs if p.is_essential]
print("toplam cift :", len(pairs))
print("kalici H0   :", sum(1 for p in essential if p.dimension == 0))
print("kalici H1   :", sum(1 for p in essential if p.dimension == 1))
print("betti       :", persistence_betti_numbers(pairs))
\end{lstlisting}

\begin{verbatim}
toplam cift : 9
kalici H0   : 1
kalici H1   : 1
betti       : {0: 1, 1: 1}
\end{verbatim}

Kalıcı sınıflar yine $(1, 1)$'i verir: bir bileşen ve bir delik. Sadece sıralı
noktalardan, hiçbir bağlantı bilgisi vermeden, pytop çemberin şeklini geri
kurtarır.

%-------------------------------------------------------------------
\section{Alıştırmalar}

\subsection*{Kodlama}

\begin{enumerate}[label=K\arabic*.]
\item \texttt{simplicial\_complex} ile içi \textbf{dolu} üçgeni
      (\texttt{[0,1,2]} facet'i, tüm yüzleriyle) kurun ve \texttt{betti\_numbers}
      çıktısını boş üçgeninkiyle karşılaştırın. $H_1$ neden kayboldu?
\item Örnek 9.3'teki nokta sayısını 8'den 4'e düşürün. Kalıcı $H_1$ hâlâ
      $1$ mi? Sonucu çemberin örnekleme yoğunluğuyla ilişkilendirin.
\end{enumerate}

\subsection*{Teori}

\begin{enumerate}[label=T\arabic*.]
\item Sonlu bir uzayın (Örnek 9.1) sürekli bir çemberle aynı $\pi_1$'e sahip
      olması nasıl mümkün? McCord teoreminin ne söylediğini bir cümleyle açıklayın.
\end{enumerate}

\chapter{Önerme Mantığı}
\label{ch:propositional_logic}

Matematiksel önerme (proposition), doğru ya da yanlış olabilen bir ifadedir.
Bu bölüm \texttt{pytop.logic} modülünün sağladığı \texttt{Proposition},
bağlaçlar ve sonlu niceleyicileri tanıtır; ardından ayrılma aksiyomlarını
bu araçlarla nasıl ifade edebileceğimizi gösterir.

% -----------------------------------------------------------------------
\section{Önerme ve Doğruluk Değeri}
\label{sec:proposition}

\begin{definition}[Önerme]
Bir \textbf{önerme} (proposition), kesin bir doğruluk değerine sahip
(doğru veya yanlış) matematiksel bir ifadedir.
\end{definition}

\begin{lstlisting}[language=Python, caption={Temel önermeler}]
from pytop import (
    Proposition,
    negate, conjunction, disjunction, implies, iff,
    for_all, there_exists, unique_exists,
)

p = Proposition("p", True)
q = Proposition("q", False)

print(f"p : {p.truth_value}")
print(f"q : {q.truth_value}")
\end{lstlisting}

\begin{lstlisting}[style=output]
p : True
q : False
\end{lstlisting}

% -----------------------------------------------------------------------
\section{Mantıksal Bağlaçlar}
\label{sec:connectives}

\begin{table}[h]
\centering
\begin{tabular}{lll}
\toprule
\textbf{Fonksiyon} & \textbf{Sembol} & \textbf{Anlam} \\
\midrule
\texttt{negate(p)}        & $\neg p$       & p değil \\
\texttt{conjunction(p,q)} & $p \wedge q$   & p ve q \\
\texttt{disjunction(p,q)} & $p \vee q$     & p veya q \\
\texttt{implies(p,q)}     & $p \to q$      & p ise q \\
\texttt{iff(p,q)}         & $p \leftrightarrow q$ & p ancak ve ancak q \\
\bottomrule
\end{tabular}
\caption{pytop.logic bağlaç fonksiyonları}
\end{table}

\begin{lstlisting}[language=Python]
print("neg(p)  :", negate(p).truth_value)
print("p and q :", conjunction(p, q).truth_value)
print("p or q  :", disjunction(p, q).truth_value)
print("p -> q  :", implies(p, q).truth_value)
print("p <-> q :", iff(p, q).truth_value)
print("p <-> p :", iff(p, p).truth_value)
\end{lstlisting}

\begin{lstlisting}[style=output]
neg(p)  : False
p and q : False
p or q  : True
p -> q  : False
p <-> q : False
p <-> p : True
\end{lstlisting}

\texttt{implies(p, q)} yalnızca $p = \text{True},\ q = \text{False}$
durumunda yanlıştır; diğer üç kombinasyonda boş doğruluk (vacuous truth)
nedeniyle doğrudur.

Aşağıdaki görsel beş bağlacın dört satırlık doğruluk değerlerini bir arada
özetler: yeşil hücre \emph{doğru} (D), kırmızı hücre \emph{yanlış} (Y).

\begin{figure}[ht]
\centering
\begin{tikzpicture}[font=\small, x=1.05cm, y=0.72cm,
    D/.style={fill=green!16, draw=green!55!black, rounded corners=2pt,
              minimum width=15pt, minimum height=13pt},
    Y/.style={fill=red!12,  draw=red!55!black, rounded corners=2pt,
              minimum width=15pt, minimum height=13pt}]
  % --- Sutun basliklari ---
  \foreach \i/\lbl in {0/$p$, 1/$q$, 2/$\neg p$, 3/$p\wedge q$, 4/$p\vee q$, 5/$p\to q$, 6/$p\leftrightarrow q$} {
    \node[anchor=center] at (\i, 4) {\lbl};
  }
  % --- Baslik / sol blok ayrac cizgileri ---
  \draw[thick] (-0.55, 3.5) -- (6.55, 3.5);
  \draw[thick] (1.55, 4.6) -- (1.55, -0.5);

  % --- Satir 0: p=D q=D ---
  \node[D] at (0,3){\scriptsize D}; \node[D] at (1,3){\scriptsize D};
  \node[Y] at (2,3){\scriptsize Y}; \node[D] at (3,3){\scriptsize D};
  \node[D] at (4,3){\scriptsize D}; \node[D] at (5,3){\scriptsize D};
  \node[D] at (6,3){\scriptsize D};
  % --- Satir 1: p=D q=Y ---
  \node[D] at (0,2){\scriptsize D}; \node[Y] at (1,2){\scriptsize Y};
  \node[Y] at (2,2){\scriptsize Y}; \node[Y] at (3,2){\scriptsize Y};
  \node[D] at (4,2){\scriptsize D}; \node[Y] at (5,2){\scriptsize Y};
  \node[Y] at (6,2){\scriptsize Y};
  % --- Satir 2: p=Y q=D ---
  \node[Y] at (0,1){\scriptsize Y}; \node[D] at (1,1){\scriptsize D};
  \node[D] at (2,1){\scriptsize D}; \node[Y] at (3,1){\scriptsize Y};
  \node[D] at (4,1){\scriptsize D}; \node[D] at (5,1){\scriptsize D};
  \node[Y] at (6,1){\scriptsize Y};
  % --- Satir 3: p=Y q=Y ---
  \node[Y] at (0,0){\scriptsize Y}; \node[Y] at (1,0){\scriptsize Y};
  \node[D] at (2,0){\scriptsize D}; \node[Y] at (3,0){\scriptsize Y};
  \node[Y] at (4,0){\scriptsize Y}; \node[D] at (5,0){\scriptsize D};
  \node[D] at (6,0){\scriptsize D};

  \node[anchor=west, font=\scriptsize, align=left] at (7.0, 3.0)
    {D = doğru\\ Y = yanlış};
  \node[anchor=west, font=\scriptsize, align=left, red!60!black] at (7.0, 1.4)
    {$p\to q$ yalnızca\\ $p$=D,\,$q$=Y'de\\ yanlıştır};
\end{tikzpicture}

\caption{Beş temel bağlacın doğruluk tablosu: $\neg p$, $p\wedge q$,
$p\vee q$, $p\to q$, $p\leftrightarrow q$. $p\to q$ yalnızca tek satırda
($p$=D, $q$=Y) yanlıştır.}
\label{fig:ch02_baglac_dogruluk}
\end{figure}

\begin{sezgi}
Bağlaçları ``kapı'' olarak düşünün. $\wedge$ (ve) her iki girdi de doğruyken
akım geçirir; $\vee$ (veya) en az biri doğru olunca geçirir; $\to$ (ise)
yalnızca ``doğru bir öncülden yanlış bir sonuç'' çıkarılırsa kapanır. İşte bu
yüzden $p\to q$ sütununda yalnızca tek bir satır ($p$=D, $q$=Y) yanlıştır —
diğer her şey, mantığın ``verdiğin sözü bozmadın'' kuralıdır.
\end{sezgi}

\begin{karsiornek}
$p\to q$ ile $p\leftrightarrow q$ aynı şey değildir. $p$=Y, $q$=D alındığında
$p\to q$ \textbf{doğru} (D), ama $p\leftrightarrow q$ \textbf{yanlış} (Y)
döner: içerme tek yönlüdür, çift-koşul iki yönü birden ister.
\end{karsiornek}

% -----------------------------------------------------------------------
\section{Algoritmalar}
\label{sec:logic_algorithms}

\subsection{Doğruluk Tablosu Üretimi}

\begin{algorithm}
\caption{Doğruluk Tablosu (İki Değişkenli)}
\label{alg:truth_table}
\begin{algorithmic}[1]
\Require Bağlaç fonksiyonu $f$
\Ensure Dört satırlı doğruluk tablosu
\For{$(t_p, t_q) \in \{(T,T),(T,F),(F,T),(F,F)\}$}
  \State $p \leftarrow \text{Proposition}("p", t_p)$
  \State $q \leftarrow \text{Proposition}("q", t_q)$
  \State Yazdır $(t_p, t_q, f(p, q).\text{truth\_value})$
\EndFor
\end{algorithmic}
\end{algorithm}

\subsection{Tautoloji Denetimi}

Bir $\phi(p_1,\ldots,p_n)$ formülü tautoloji iff tüm $2^n$ atama için
doğrudur — $O(2^n \cdot |\phi|)$.

% -----------------------------------------------------------------------
\section{pytop API}
\label{sec:logic_api}

\begin{lstlisting}[language=Python, caption={logic modülü importu}]
from pytop import (
    Proposition,
    negate, conjunction, disjunction, implies, iff,
    for_all, there_exists, unique_exists,
)
\end{lstlisting}

Niceleyiciler:
\begin{lstlisting}[language=Python]
for_all(carrier, predicate)       # tum x: P(x)
there_exists(carrier, predicate)  # bazi x: P(x)
unique_exists(carrier, predicate) # tam bir x: P(x)
\end{lstlisting}

% -----------------------------------------------------------------------
\section{Örnekler}
\label{sec:logic_examples}

\subsection{Doğruluk Tablosu}

\begin{lstlisting}[language=Python]
pairs = [(True, True), (True, False), (False, True), (False, False)]
print(f"{'p':>5}  {'q':>5}  {'neg_p':>5}  {'p&q':>5}  {'p|q':>5}  {'p->q':>6}  {'p<->q':>6}")
for tv_p, tv_q in pairs:
    pp = Proposition("p", tv_p)
    qq = Proposition("q", tv_q)
    row = (pp.truth_value, qq.truth_value,
           negate(pp).truth_value, conjunction(pp, qq).truth_value,
           disjunction(pp, qq).truth_value,
           implies(pp, qq).truth_value, iff(pp, qq).truth_value)
    print("  ".join(f"{str(v):>5}" for v in row))
\end{lstlisting}

\begin{lstlisting}[style=output]
    p      q  neg_p    p&q    p|q   p->q  p<->q
----------------------------------------------------
 True   True  False   True   True   True   True
 True  False  False  False   True  False  False
False   True   True  False   True   True  False
False  False   True  False  False   True   True
\end{lstlisting}

\subsection{Tautolojiler — De Morgan ve Karşıt}

\begin{theorem}[De Morgan Yasaları]
\label{thm:de_morgan}
\begin{align*}
\neg(p \wedge q) &\leftrightarrow \neg p \vee \neg q \\
\neg(p \vee q) &\leftrightarrow \neg p \wedge \neg q
\end{align*}
\end{theorem}

\begin{theorem}[Karşıt (Contrapositive)]
\label{thm:contrapositive}
$(p \to q) \leftrightarrow (\neg q \to \neg p)$
\end{theorem}

En sık yapılan mantık hatası, bir içermeyi \textbf{tersiyle (converse)}
karıştırmaktır. $p\to q$'nin doğru biçimde denk olduğu ifade
\emph{kontrapozitiftir}: $\neg q\to\neg p$. Buna karşılık $q\to p$ (tersi)
genellikle denk \emph{değildir}.

\begin{figure}[ht]
\centering
\begin{tikzpicture}[font=\small, node distance=11mm,
    box/.style={draw, rounded corners, fill=blue!7, inner sep=6pt, minimum width=20mm},
    denk/.style={draw=green!55!black, rounded corners, fill=green!8, inner sep=6pt, minimum width=20mm}]

  % --- Ust satir: dogru denklik p->q  <=>  ¬q->¬p ---
  \node[box]  (pq)  {$p \to q$};
  \node[denk, right=28mm of pq] (cp) {$\neg q \to \neg p$};
  \draw[{Stealth[length=2.6mm]}-{Stealth[length=2.6mm]}, thick, green!45!black]
    (pq) -- node[above=1pt, font=\scriptsize, green!45!black] {denk} (cp);
  \node[font=\scriptsize, green!45!black, above=9mm] at ($(pq)!0.5!(cp)$) {kontrapozitif};

  % --- Alt satir: yanlis denklik p->q  ⇏  q->p ---
  \node[box, below=20mm of pq]  (pq2)  {$p \to q$};
  \node[box, right=28mm of pq2] (qp)   {$q \to p$};
  \draw[{Stealth[length=2.6mm]}-{Stealth[length=2.6mm]}, thick, red!65!black, dashed]
    (pq2) -- node[below=1pt, font=\scriptsize, red!65!black] {denk \emph{değil}} (qp);
  \node[font=\scriptsize, red!65!black, below=9mm] at ($(pq2)!0.5!(qp)$)
    {tersi (converse) — yaygın hata};

  % --- Karsi-ornek aciklamasi ---
  \node[anchor=west, font=\scriptsize, align=left, draw=red!40, fill=red!4,
        rounded corners, inner sep=5pt] at (6.7, -1.45)
    {Karşı-örnek: $p$=``kare'', $q$=``dikdörtgen''.\\
     $p\to q$ doğru, ama $q\to p$ yanlış\\
     (her dikdörtgen kare değildir).};
\end{tikzpicture}

\caption{Kontrapozitif denkliği $p\to q \Leftrightarrow \neg q\to\neg p$
(yeşil) ile denk \emph{olmayan} tersi $q\to p$ (kırmızı kesikli).}
\label{fig:ch02_kontrapozitif}
\end{figure}

\begin{karsiornek}
$p$ = ``$x$ bir karedir'', $q$ = ``$x$ bir dikdörtgendir'' olsun. $p\to q$
doğrudur (her kare dikdörtgendir), ama tersi $q\to p$ yanlıştır (her dikdörtgen
kare değildir). Kontrapozitif $\neg q\to\neg p$ = ``dikdörtgen değilse kare de
değildir'' ise yine doğrudur — çünkü o $p\to q$ ile denktir.
\end{karsiornek}

\paragraph{İspat eskizi (kontrapozitif).}
$p\to q$ ile $\neg q\to\neg p$ aynı doğruluk sütununu üretir; dört satırda da
çakışırlar, dolayısıyla $(p\to q)\leftrightarrow(\neg q\to\neg p)$ tautolojidir.
\[
\begin{array}{cc|cc}
p & q & p\to q & \neg q\to\neg p \\ \hline
\text{D} & \text{D} & \text{D} & \text{D} \\
\text{D} & \text{Y} & \text{Y} & \text{Y} \\
\text{Y} & \text{D} & \text{D} & \text{D} \\
\text{Y} & \text{Y} & \text{D} & \text{D}
\end{array}
\]

\paragraph{İspat eskizi (De Morgan).}
$p\wedge q$ yalnızca her ikisi doğruyken doğrudur; öyleyse $\neg(p\wedge q)$
``en az biri yanlış'' demektir. $\neg p\vee\neg q$ de tam olarak bunu söyler.
İki sütun her atamada çakıştığından $\neg(p\wedge q)\leftrightarrow\neg p\vee\neg q$
tautolojidir; ikincil yasa $\neg(p\vee q)\leftrightarrow\neg p\wedge\neg q$
$\wedge$ ile $\vee$ rolleri değiştirilerek aynı şekilde gösterilir.
\[
\begin{array}{cc|ccc}
p & q & p\wedge q & \neg(p\wedge q) & \neg p\vee\neg q \\ \hline
\text{D} & \text{D} & \text{D} & \text{Y} & \text{Y} \\
\text{D} & \text{Y} & \text{Y} & \text{D} & \text{D} \\
\text{Y} & \text{D} & \text{Y} & \text{D} & \text{D} \\
\text{Y} & \text{Y} & \text{Y} & \text{D} & \text{D}
\end{array}
\]

\begin{lstlisting}[language=Python]
def check_tautology(name, func):
    ok = all(func(Proposition("p", tp), Proposition("q", tq)).truth_value
             for tp, tq in [(True,True),(True,False),(False,True),(False,False)])
    print(f"{name}: {'tautoloji' if ok else 'DEGIL'}")

check_tautology("neg(p&q) <-> neg_p|neg_q",
    lambda p, q: iff(negate(conjunction(p, q)), disjunction(negate(p), negate(q))))
check_tautology("(p->q) <-> (neg_q->neg_p)",
    lambda p, q: iff(implies(p, q), implies(negate(q), negate(p))))
\end{lstlisting}

\begin{lstlisting}[style=output]
neg(p&q) <-> neg_p|neg_q: tautoloji
(p->q) <-> (neg_q->neg_p): tautoloji
\end{lstlisting}

\subsection{Niceleyiciler}

\begin{lstlisting}[language=Python]
X = [1, 2, 3, 4, 5]
print("for_all(X, x>0)       :", for_all(X, lambda x: x > 0))
print("for_all(X, x>2)       :", for_all(X, lambda x: x > 2))
print("there_exists(X, x>4)  :", there_exists(X, lambda x: x > 4))
print("unique_exists(X, x==3):", unique_exists(X, lambda x: x == 3))
\end{lstlisting}

\begin{lstlisting}[style=output]
for_all(X, x>0)       : True
for_all(X, x>2)       : False
there_exists(X, x>4)  : True
unique_exists(X, x==3): True
\end{lstlisting}

\subsection{Topoloji Uygulaması — T0 ve T1}

\begin{definition}[T0 — Kolmogorov]
$\forall x \neq y \in X,\ \exists$ açık $U : (x \in U,\ y \notin U)$ veya
$(y \in U,\ x \notin U)$.
\end{definition}

\begin{definition}[T1 — Fréchet]
$\forall x \neq y \in X,\ \exists$ açık $U : x \in U,\ y \notin U$.
\end{definition}

\begin{lstlisting}[language=Python]
from pytop import sierpinski_space, discrete_topology, indiscrete_topology

def is_t0_logic(space):
    pts   = list(space.carrier)
    opens = list(space.topology)
    return for_all(
        [(x, y) for x in pts for y in pts if x != y],
        lambda pair: there_exists(
            opens, lambda U: (pair[0] in U) != (pair[1] in U)
        )
    )

def is_t1_logic(space):
    pts   = list(space.carrier)
    opens = list(space.topology)
    return for_all(
        [(x, y) for x in pts for y in pts if x != y],
        lambda pair: there_exists(
            opens, lambda U: pair[0] in U and pair[1] not in U
        )
    )

for name, sp in [("Sierpinski", sierpinski_space()),
                 ("Discrete  ", discrete_topology(0, 1)),
                 ("Indiscrete", indiscrete_topology(0, 1))]:
    print(f"{name}  T0={is_t0_logic(sp)}  T1={is_t1_logic(sp)}")
\end{lstlisting}

\begin{lstlisting}[style=output]
Sierpinski  T0=True   T1=False
Discrete    T0=True   T1=True
Indiscrete  T0=False  T1=False
\end{lstlisting}

% -----------------------------------------------------------------------
\section{Alıştırmalar}
\label{sec:logic_exercises}

\begin{enumerate}
  \item Beş önerme değeriyle tam bir doğruluk tablosu oluşturun; \texttt{implies}
        kaç \texttt{(p, q)} çiftinde \texttt{False} döner?

  \item \texttt{unique\_exists([0,1,2,3,4], predicate)} değerini \texttt{True}
        yapan üç farklı koşul yazın.

  \item T2 (Hausdorff) aksiyomunu \texttt{for\_all} ve \texttt{there\_exists}
        kullanarak kodlayın:
        $\forall x \neq y,\ \exists$ açık $U,V : x \in U,\ y \in V,\ U \cap V = \emptyset$.
        Sierpiński, discrete ve indiscrete için test edin.

  \item \textit{(Teori)} \texttt{implies(p, q)} neden yalnızca $p = T, q = F$
        durumunda yanlıştır? Boş doğruluk (vacuous truth) kavramını açıklayın.

  \item \textit{(Teori)} De Morgan yasalarını önerme mantığı için ispatlayın:
        $\neg(p \wedge q) \leftrightarrow \neg p \vee \neg q$.

  \item \texttt{check\_tautology} yardımcı fonksiyonunu kullanarak
        \textbf{dağılma yasasını} doğrulayın:
        $p \wedge (q \vee r) \leftrightarrow (p \wedge q) \vee (p \wedge r)$.
        Üç değişken olduğundan sekiz $(p, q, r)$ atamasının tümünü gezmeniz
        gerekir; \texttt{iff}'in her atamada \texttt{True} döndüğünü gösterin.

  \item \textit{(Teori)} Bir öğrenci ``$p \to q$ doğruysa $q \to p$ de
        doğrudur'' diyor. Bu iddianın neden yanlış olduğunu, kontrapozitif
        $\neg q \to \neg p$ ile tersi $q \to p$ arasındaki farkı vurgulayan
        somut bir karşı-örnekle açıklayın. \ipucu{$p$ = ``tam sayı 4'e bölünür'',
        $q$ = ``tam sayı çifttir''.}
\end{enumerate}

\chapter{Küme Teorisi ve Fonksiyon Temelleri}

pytop'un tüm modülleri küme teorisi ve fonksiyon kavramları üzerine kuruludur.
Bu bölüm kılavuzun geri kalanında varsayılan ön koşulları pytop API'siyle pekiştirir.

%-------------------------------------------------------------------
\section{Konu}

\begin{definition}[Altküme ve Güç Kümesi]
$A \subseteq B \Longleftrightarrow \forall a \in A,\, a \in B$.
\textbf{Güç kümesi:} $\mathcal{P}(A) = \{S : S \subseteq A\}$,\quad
$|\mathcal{P}(A)| = 2^{|A|}$.
\end{definition}

İki kümenin işlemleri Venn şemasıyla görselleştirilir: $A\cup B$ üç bölgenin
tamamı, $A\cap B$ ortadaki örtüşen bölge, $A\setminus B$ ise $A$'nın $B$'den
ayrık kalan kısmıdır.

\begin{figure}[ht]
\centering
\begin{tikzpicture}[scale=1.0, font=\small]
  % evren kutusu
  \draw[gray!60, thick] (-2.7,-2.0) rectangle (2.7,2.0);
  \node[gray!60] at (2.35,1.7) {$X$};
  % A ve B daireleri (kesisim taranir)
  \begin{scope}
    \clip (-0.9,0) circle (1.5);
    \fill[violet!18] (0.9,0) circle (1.5);
  \end{scope}
  \draw[blue!70!black, thick] (-0.9,0) circle (1.5);
  \draw[violet!80!black, thick] (0.9,0) circle (1.5);
  \node[blue!70!black] at (-1.9,1.25) {$A$};
  \node[violet!80!black] at (1.9,1.25) {$B$};
  % bolgeler
  \node at (-1.15,0) {$A\setminus B$};
  \node at (0,0) {$A\cap B$};
  \node at (1.15,0) {$B\setminus A$};
  \node[align=center] at (0,-2.55)
    {$A\cup B$ = her uc bolge\\ $A\cap B$ = orta bolge \quad $A\setminus B$ = sol bolge};
\end{tikzpicture}

\caption{İki küme için $A\cup B$, $A\cap B$ ve $A\setminus B$ bölgeleri.}
\end{figure}

\begin{sezgi}
Tümleyen her zaman bir \textbf{evrene} ($X$) göre tanımlıdır; ``$A^c$'' tek
başına anlamsızdır. pytop bunu \texttt{complement(subset, universe)} imzasıyla
zorunlu kılar.
\end{sezgi}

\begin{karsiornek}
``$A\cup B$ her zaman $A$'dan büyüktür'' yanlıştır: $B\subseteq A$ ise
$A\cup B = A$. Örneğin $A=\{1,2,3\}$, $B=\{2\}$ için $|A\cup B|=|A|=3$.
\end{karsiornek}

\begin{definition}[Kartezyen Çarpım]
$A \times B = \{(a,b) : a\in A,\ b\in B\}$ ve $|A\times B| = |A|\cdot|B|$.
$A$ üzerinde bir bağıntı $R \subseteq A\times A$.
\end{definition}

\begin{figure}[ht]
\centering
\begin{tikzpicture}[scale=1.2, font=\small]
  % 2 (yatay, A) x 3 (dikey, B) izgara
  \draw[gray!40] (0,0) grid (2,3);
  % A ekseni (yatay): a, b
  \node[blue!70!black] at (0.5,-0.4) {$a$};
  \node[blue!70!black] at (1.5,-0.4) {$b$};
  \node[blue!70!black] at (1.0,-0.95) {$A=\{a,b\}$};
  % B ekseni (dikey): 1, 2, 3
  \node[violet!80!black] at (-0.4,0.5) {$1$};
  \node[violet!80!black] at (-0.4,1.5) {$2$};
  \node[violet!80!black] at (-0.4,2.5) {$3$};
  \node[violet!80!black, rotate=90] at (-1.0,1.5) {$B=\{1,2,3\}$};
  % carpim noktalari (a,1)...(b,3)
  \foreach \x/\xl in {0.5/a, 1.5/b}
    \foreach \y/\yl in {0.5/1, 1.5/2, 2.5/3} {
      \fill[red!70!black] (\x,\y) circle (2.2pt);
      \node[anchor=south west, font=\scriptsize] at (\x+0.07,\y+0.05) {$(\xl,\yl)$};
    }
  \node at (1.0,-1.6) {$|A\times B| = |A|\cdot|B| = 2\cdot 3 = 6$};
\end{tikzpicture}

\caption{$A\times B$ kartezyen çarpım ızgarası ($|A\times B|=6$).}
\end{figure}

\begin{karsiornek}
Kartezyen çarpım değişmeli değildir: $A\times B \ne B\times A$ (genel olarak).
$(a,1)\in A\times B$ ama $(a,1)\notin B\times A$, çünkü $B\times A$'nın ikilileri
$(1,a)$ biçimindedir.
\end{karsiornek}

\begin{definition}[Denklik Bağıntısı]
$A$ üzerinde $R \subseteq A \times A$ bağıntısı \textbf{denklik bağıntısıdır} iff:
yansımalı ($\forall a,(a,a)\in R$), simetrik ($(a,b)\in R\Rightarrow(b,a)\in R$),
geçişli ($(a,b),(b,c)\in R\Rightarrow(a,c)\in R$).
Denklik sınıfı $[a] = \{b \in A : (a,b) \in R\}$.
\end{definition}

\begin{definition}[Fonksiyon Türleri]
$f: A \to B$ için:
\begin{itemize}
  \item \textbf{Enjeksiyon:} $f(a_1)=f(a_2) \Rightarrow a_1=a_2$
  \item \textbf{Sürjeksiyon:} $\forall b\in B,\, \exists a\in A: f(a)=b$
  \item \textbf{Bijeksiyon:} Enjeksiyon $+$ Sürjeksiyon
\end{itemize}
\end{definition}

%-------------------------------------------------------------------
\section{Teoremler}

\begin{theorem}[De Morgan]
$(A \cup B)^c = A^c \cap B^c$ ve $(A \cap B)^c = A^c \cup B^c$.
\end{theorem}

\begin{proof}[İspat eskizi]
Çift kapsama ile: $x \in (A\cup B)^c \Leftrightarrow x \notin A\cup B
\Leftrightarrow x\notin A \text{ ve } x\notin B \Leftrightarrow x\in A^c\cap B^c$.
Her adım çift yönlü olduğundan iki küme aynı elemanlara sahiptir. İkinci eşitlik
``veya''nın değillemesiyle simetrik biçimde elde edilir.
\end{proof}

\begin{theorem}[Güç Kümesi Boyutu]
$|A| = n \Rightarrow |\mathcal{P}(A)| = 2^n$.
\end{theorem}

\begin{proof}[İspat eskizi]
Her $S\subseteq A$, $A$'nın $n$ elemanının her biri için ``$S$ içinde mi?''
sorusuna evet/hayır cevabı veren bir ikili dizgiyle birebir eşleşir. $n$
bağımsız ikili seçim olduğundan toplam $2^n$. Tümevarımla: $|A|=0$ için
$|\mathcal{P}|=1=2^0$; bir eleman eklendiğinde her eski alt küme hem eleman-sız
hem eleman-lı sayılır, miktar ikiye katlanır.
\end{proof}

\begin{theorem}[Denklik Sınıfları Bölüntü Oluşturur]
$\sim$ denklik bağıntısı ise $\{[a] : a \in A\}$ $A$'nın bir bölüntüsüdür:
sınıflar örtüşmez ve birleşimleri $A$'dır.
\end{theorem}

\begin{theorem}[Ön-Görüntü Özellikleri]
$f: A \to B$ ve $S_1, S_2 \subseteq B$ için:
$f^{-1}(S_1 \cup S_2) = f^{-1}(S_1) \cup f^{-1}(S_2)$ ve
$f^{-1}(S_1 \cap S_2) = f^{-1}(S_1) \cap f^{-1}(S_2)$.
\end{theorem}

%-------------------------------------------------------------------
\section{Algoritmalar}

\subsection{Güç Kümesi — $O(2^{|A|})$}

\begin{algorithm}
\caption{PowerSet$(A)$}
\begin{algorithmic}[1]
\State result $\leftarrow \{\emptyset\}$
\For{each $a \in A$}
    \State result $\leftarrow$ result $\cup\ \{S \cup \{a\} : S \in \text{result}\}$
\EndFor
\State \Return result
\end{algorithmic}
\end{algorithm}

\subsection{Denklik Bağıntısı Doğrulama — $O(|R| \cdot |A|)$}

\begin{algorithm}
\caption{IsEquivalence$(A, R)$}
\begin{algorithmic}[1]
\For{each $a \in A$}
    \If{$(a,a) \notin R$} \Return \textbf{False} \EndIf
\EndFor
\For{each $(a,b) \in R$}
    \If{$(b,a) \notin R$} \Return \textbf{False} \EndIf
\EndFor
\For{each $(a,b) \in R$, $(b,c) \in R$}
    \If{$(a,c) \notin R$} \Return \textbf{False} \EndIf
\EndFor
\State \Return \textbf{True}
\end{algorithmic}
\end{algorithm}

%-------------------------------------------------------------------
\section{pytop API}

\begin{lstlisting}[language=Python]
from pytop import (
    make_set, power_set, empty_set,
    set_union, set_intersection, set_difference, complement,
    cartesian_product, indexed_union, indexed_intersection,
    is_subset, is_proper_subset, equal_sets,
    make_relation, is_equivalence_relation,
    identity_relation, inverse_relation, compose_relations,
    relation_profile,
    equivalence_class, partition_from_equivalence,
    total_order_from_list, is_partial_order, is_total_order,
    normalize_finite_map_data,
    is_injective_finite_map, is_surjective_finite_map, is_bijective_finite_map,
    image_of_subset_finite, preimage_of_subset_finite,
)
\end{lstlisting}

\texttt{make\_set(*elements)} $\to$ \texttt{frozenset}

\texttt{power\_set(values)} $\to$ \texttt{set[frozenset]}

\texttt{complement(subset, universe)} $\to$ \texttt{set} — evrene göre $A^c$

\texttt{cartesian\_product(left, right)} $\to$ \texttt{set[tuple]} — $A\times B$

\texttt{indexed\_union(family)} / \texttt{indexed\_intersection(family)} $\to$ \texttt{set}

\texttt{make\_relation(carrier, *pairs)} $\to$ \texttt{set[tuple]}

\texttt{relation\_profile(carrier, relation)} $\to$ \texttt{dict}

\texttt{normalize\_finite\_map\_data(domain, codomain, mapping)} $\to$ \texttt{FiniteMapData}

%-------------------------------------------------------------------
\section{Örnekler}

\subsection*{Örnek 5.1 — Küme İşlemleri}

\begin{lstlisting}[language=Python]
A = make_set(1, 2, 3)
B = make_set(2, 3, 4)
print("A union B:", sorted(set_union(A, B)))
print("A inter B:", sorted(set_intersection(A, B)))
print("A diff B:",  sorted(set_difference(A, B)))
print("{2,3} subset A:", is_subset({2, 3}, A))
\end{lstlisting}

\begin{verbatim}
A union B: [1, 2, 3, 4]
A inter B: [2, 3]
A diff B: [1]
{2,3} subset A: True
\end{verbatim}

\subsection*{Örnek 5.2 — Güç Kümesi}

\begin{lstlisting}[language=Python]
P = power_set([1, 2, 3])
print("boyut:", len(P))  # 8 = 2^3
\end{lstlisting}

\begin{verbatim}
boyut: 8
\end{verbatim}

\subsection*{Örnek 5.3 — Denklik Bağıntısı}

\begin{lstlisting}[language=Python]
carrier = [1, 2, 3, 4]
rel_eq = make_relation(carrier,
    (1,1),(2,2),(3,3),(4,4),(1,3),(3,1),(2,4),(4,2))
print("denklik:", is_equivalence_relation(carrier, rel_eq))
prof = relation_profile(carrier, rel_eq)
print("yansimalı:", prof['is_reflexive'])
print("simetrik:", prof['is_symmetric'])
print("gecisli:", prof['is_transitive'])

# Gecissiz (1,2),(2,3) var ama (1,3) yok
rel_bad = make_relation(carrier,
    (1,1),(2,2),(3,3),(4,4),(1,2),(2,1),(2,3),(3,2))
print("gecissiz denklik mi:", is_equivalence_relation(carrier, rel_bad))
\end{lstlisting}

\begin{verbatim}
denklik: True
yansimalı: True
simetrik: True
gecisli: True
gecissiz denklik mi: False
\end{verbatim}

\subsection*{Örnek 5.4 — Fonksiyon Türleri}

\begin{lstlisting}[language=Python]
f = normalize_finite_map_data(
    [1,2,3], ['a','b','c'], {1:'a', 2:'b', 3:'c'})
print("bijeksiyon:", is_bijective_finite_map(f))

g = normalize_finite_map_data(
    [1,2,3], ['x','y'], {1:'x', 2:'x', 3:'y'})
print("g enjeksiyon:", is_injective_finite_map(g))
print("g surjeksiyon:", is_surjective_finite_map(g))
\end{lstlisting}

\begin{verbatim}
bijeksiyon: True
g enjeksiyon: False
g surjeksiyon: True
\end{verbatim}

\subsection*{Örnek 5.5 — Görüntü ve Ön-Görüntü}

\begin{lstlisting}[language=Python]
f = normalize_finite_map_data(
    [1,2,3,4], ['a','b','c','d'], {1:'a',2:'b',3:'c',4:'d'})
print("f({1,2}):", sorted(image_of_subset_finite(f, [1, 2])))
print("f^-1({b,c}):", sorted(preimage_of_subset_finite(f, ['b','c'])))
\end{lstlisting}

\begin{verbatim}
f({1,2}): ['a', 'b']
f^-1({b,c}): [2, 3]
\end{verbatim}

\subsection*{Örnek 5.6 — De Morgan Yasaları (Tümleyenle)}

\begin{lstlisting}[language=Python]
X = make_set(1, 2, 3, 4, 5, 6)
A = make_set(1, 2, 3)
B = make_set(3, 4, 5)
lhs1 = complement(set_union(A, B), X)
rhs1 = set_intersection(complement(A, X), complement(B, X))
print("(A u B)^c:", sorted(lhs1), "==", sorted(rhs1))
lhs2 = complement(set_intersection(A, B), X)
rhs2 = set_union(complement(A, X), complement(B, X))
print("(A n B)^c:", sorted(lhs2), "==", sorted(rhs2))
\end{lstlisting}

\begin{verbatim}
(A u B)^c: [6] == [6]
(A n B)^c: [1, 2, 4, 5, 6] == [1, 2, 4, 5, 6]
\end{verbatim}

\subsection*{Örnek 5.7 — Kartezyen Çarpım}

\begin{lstlisting}[language=Python]
P = make_set('a', 'b')
Q = make_set(1, 2, 3)
prod = cartesian_product(P, Q)
print("|PxQ| =", len(prod))  # 2*3 = 6
print(sorted(prod, key=lambda t: (t[0], t[1])))
\end{lstlisting}

\begin{verbatim}
|PxQ| = 6
[('a', 1), ('a', 2), ('a', 3), ('b', 1), ('b', 2), ('b', 3)]
\end{verbatim}

\subsection*{Örnek 5.8 — İndeksli Aile}

\begin{lstlisting}[language=Python]
fam = {0: [1, 2, 3], 1: [2, 3, 4], 2: [3, 4, 5]}
print("birlesim:", sorted(indexed_union(fam)))
print("kesisim :", sorted(indexed_intersection(fam)))
\end{lstlisting}

\begin{verbatim}
birlesim: [1, 2, 3, 4, 5]
kesisim : [3]
\end{verbatim}

%-------------------------------------------------------------------
\section{Alıştırmalar}

\subsection*{Kodlama}

\begin{enumerate}[label=K\arabic*.]
\item $A = \{1,2,3,4,5\}$, $B = \{3,4,5,6,7\}$ için $A\cup B$, $A\cap B$,
      $A\setminus B$'yi hesaplayın.
\item $\{1,2,3,4\}$ için güç kümesini üretin; $|\mathcal{P}|=16$ olduğunu doğrulayın.
\item $R = \{(i,j): 0\le i,j \le 2\}$ (tam çarpım) üzerinde \texttt{relation\_profile}
      çalıştırın ve denklik olduğunu gösterin.
\item $X=\{1,\dots,8\}$, $A=\{1,2,3,4\}$, $B=\{3,4,5,6\}$ için \texttt{complement}
      ile her iki De Morgan eşitliğini doğrulayın.
\item $P=\{x,y,z\}$, $Q=\{0,1\}$ için \texttt{cartesian\_product} ile $P\times Q$
      ve $Q\times P$ üretin; boyutların eşit ama kümelerin farklı olduğunu
      \texttt{equal\_sets} ile gösterin.
\end{enumerate}

\subsection*{Teori}

\begin{enumerate}[label=T\arabic*.]
\item $A \subseteq B \Leftrightarrow A \cap B = A$ olduğunu ispatlayın.
\item $f: A \to B$ ve $g: B \to C$ enjeksiyon ise $g \circ f$ de enjeksiyondur.
\end{enumerate}


% ── Part I: Nokta-Küme Topolojisi ─────────────────────────────────────────────
\part{Nokta-Küme Topolojisi}

\chapter{Topolojik Uzaylar}
\label{ch:topological_spaces}

%% ============================================================
\section{Konu}
%% ============================================================

\subsection{Sezgisel Giriş}

Topoloji, bir kümedeki noktaların birbirine ``yakın'' olup olmadığını, aralarındaki
sürekliliği ve şekilsel özellikleri, mesafe kavramına gerek duymadan inceleyen matematiksel
bir disiplindir. Bunu mümkün kılan temel araç \emph{açık küme} kavramıdır.

Düşünelim: gerçek doğruda $(a, b)$ aralığı ``açık''tır; her noktasının çevresinde küçük bir
açık aralık daha bulunur. Bu sezgiyi soyutlayarak \emph{herhangi} bir küme üzerinde topoloji
tanımlamak mümkündür.

\begin{sezgi}
Bir şehir haritasında ``yakınlık''ı sokak mesafesiyle ölçebilirsiniz; ama
``hangi mahalleler birbirine komşu?'' sorusu mesafe olmadan da anlamlıdır.
Topoloji tam bunu yapar: mesafeyi unutur, yalnızca ``hangi kümeler bir
noktanın çevresini oluşturur?'' bilgisini tutar. Matematiksel karşılığı:
$\tau$ ailesi her noktanın ``çevre'' kavramını açık kümeler aracılığıyla
kodlar; süreklilik ve yakınsama gibi kavramlar yalnız $\tau$'ya
başvurularak tanımlanır.
\end{sezgi}

\subsection{Formal Tanım}

\begin{definition}[Topolojik Uzay]
\label{def:topological_space}
Bir $X$ kümesi ve $\tau \subseteq \mathcal{P}(X)$ ailesi verilsin.
$(X, \tau)$ bir \textbf{topolojik uzay}, $\tau$ ise $X$ üzerinde bir \textbf{topoloji}
olarak adlandırılır, eğer aşağıdaki üç aksiyom sağlanıyorsa:
\begin{enumerate}
  \item[(T1)] $\emptyset \in \tau$ ve $X \in \tau$
  \item[(T2)] $U, V \in \tau \Rightarrow U \cap V \in \tau$
  \item[(T3)] $\{U_\alpha\}_{\alpha \in I} \subseteq \tau \Rightarrow
              \displaystyle\bigcup_{\alpha \in I} U_\alpha \in \tau$
\end{enumerate}
$\tau$'nun elemanları \textbf{açık küme} olarak adlandırılır.
\end{definition}

Aksiyomların her biri ayrı bir iş görür: (T1) uzayın tamamının ve boş
kümenin her zaman ``gözlemlenebilir'' olmasını garanti eder; (T2) iki
gözlemin kesişiminin de gözlem olmasını sağlar --- sonlu kesişimle sınırlı
kalması bilinçli bir tercihtir (bkz.\ Teorem~\ref{thm:axiom_sufficiency});
(T3) keyfî birleşime izin vererek küçük çevrelerden büyük çevre kurma
işlemini serbest bırakır.

\begin{karsiornek}
$X=\{1,2,3\}$ üzerinde $\sigma=\{\emptyset,\{1\},\{2\},X\}$ ailesi bir
topoloji \emph{değildir}: $\{1\}\cup\{2\}=\{1,2\}\notin\sigma$
olduğundan (T3) ihlal edilir. (T1) ve (T2) sağlansa bile tek bir eksik
birleşim aileyi topoloji olmaktan çıkarır. Bu ailenin
\texttt{make\_topology}'ye verilince ne olduğunu Örnekler bölümündeki
``Dikkat'' kutusunda göreceğiz.
\end{karsiornek}

\subsection{Temel Örnekler}

\begin{table}[h]
\centering
\begin{tabular}{lll}
\toprule
\textbf{Topoloji} & \textbf{Tanım} & \textbf{Açık kümeler} \\
\midrule
Ayrık (discrete) & Her alt küme açık & $\tau = \mathcal{P}(X)$ \\
İndirgenmiş (indiscrete) & Yalnızca $\emptyset$ ve $X$ açık & $\tau = \{\emptyset, X\}$ \\
Kosonlu (cofinite) & $U$ açık $\iff$ $U=\emptyset$ veya $X\setminus U$ sonlu & — \\
Sierpiński & $X=\{0,1\}$, $\tau = \{\emptyset, \{1\}, X\}$ & 3 açık küme \\
\bottomrule
\end{tabular}
\caption{Standart topoloji örnekleri}
\end{table}

\subsection{Baz ve Alt-Baz}

\begin{definition}[Baz]
\label{def:basis}
$\mathcal{B} \subseteq \tau$ ailesi, her $U \in \tau$ ve her $x \in U$ için
$x \in B \subseteq U$ sağlayan bir $B \in \mathcal{B}$ varsa $\tau$'ya \textbf{baz} denir.
\end{definition}

Şekil~\ref{fig:ch04:baz} koşulu resmeder: açık kümenin her noktası,
tamamen o kümenin içinde kalan bir baz elemanıyla sarılır.

\begin{figure}[ht]
  \centering
  \begin{tikzpicture}[scale=1.0, font=\small]
  \draw[blue!70!black, thick, fill=blue!8]
    plot[smooth cycle, tension=0.8] coordinates
    {(0,0) (2.2,-0.4) (4.2,0.2) (4.6,2.0) (2.6,3.0) (0.3,2.4)};
  \node[blue!70!black] at (4.35,2.65) {$U$};
  \draw[violet!80!black, thick, fill=violet!12]
    (2.2,1.3) ellipse [x radius=1.15, y radius=0.7];
  \node[violet!80!black] at (3.05,1.95) {$B$};
  \fill (2.2,1.25) circle (1.8pt) node[below right=0.5pt] {$x$};
  \node[align=left, anchor=west] at (5.2,1.5)
    {Her $x\in U$ için\\ $x\in B\subseteq U$ olan\\ bir $B\in\mathcal{B}$ vardır.};
\end{tikzpicture}

  \caption{Baz koşulunun geometrisi: $U$ açık kümesinin her $x$ noktası
    için $x\in B\subseteq U$ sağlayan bir $B\in\mathcal{B}$ vardır.
    Topolojinin tüm açıkları bu tür ``yapı taşları''nın birleşimleridir.}
  \label{fig:ch04:baz}
\end{figure}

\begin{nedenonemli}
Baz, büyük bir topolojiyi küçük bir çekirdekle temsil etme aracıdır.
Gerçek doğrunun standart topolojisi sayılamaz çoklukta açık küme içerir;
ama tümü, sayılabilir $\{(a,b): a<b,\ a,b\in\mathbb{Q}\}$ bazından
üretilir. \texttt{pytop} tarafında \texttt{topology\_from\_basis} tam bu
ilkeyle çalışır: yalnız baz elemanlarını verirsiniz, kapanışı kütüphane
hesaplar (Algoritma~\ref{alg:basis_to_topology}).
\end{nedenonemli}

\noindent\textbf{Baz Koşulları:}
\begin{enumerate}
  \item[(B1)] $\bigcup \mathcal{B} = X$
  \item[(B2)] $B_1, B_2 \in \mathcal{B}$ ve $x \in B_1 \cap B_2$ ise bir $B_3 \in \mathcal{B}$
              vardır: $x \in B_3 \subseteq B_1 \cap B_2$
\end{enumerate}

\begin{definition}[Alt-Baz]
$\mathcal{S} \subseteq \mathcal{P}(X)$ ailesi; sonlu kesişimleri baz, onların birleşimleri
ise $\tau$'yu oluşturur.
\end{definition}

\subsection{Sonlu ve Sonsuz Uzaylar}

\texttt{pytop}'ta iki temel sınıf bulunur:
\begin{itemize}
  \item \textbf{\texttt{FiniteTopologicalSpace}:} Taşıyıcı $(X)$ sonlu; topoloji açıkça
        \texttt{frozenset}'ler koleksiyonu olarak saklanır. Tüm hesaplamalar tam ve kesindir.
  \item \textbf{Sonsuz türevler:} Taşıyıcı simgesel (\texttt{'R'}, \texttt{'N'});
        \texttt{topology=None}; etiket tabanlı sembolik çıkarım yapılır.
\end{itemize}

%% ============================================================
\section{Teoremler}
%% ============================================================

\begin{theorem}[Aksiyomların Yeterliliği]
\label{thm:axiom_sufficiency}
$(X, \tau)$ bir topolojik uzay olsun. T1--T3 aksiyomları, sonlu herhangi bir açık küme
ailesinin kesişiminin $\tau$'da yer aldığını garanti eder:
\[
  U_1, \ldots, U_n \in \tau \Rightarrow \bigcap_{i=1}^{n} U_i \in \tau.
\]
\textit{Not:} Sonsuz kesişim gerekmez. Gerçek doğruda
$\bigcap_{n=1}^{\infty}\bigl(-\tfrac{1}{n}, \tfrac{1}{n}\bigr) = \{0\}$
açık değildir.
\end{theorem}

\begin{proof}[Rehberli Kanıt]
\textbf{Strateji:} $n$ üzerinde tümevarım; tek araç (T2).
\begin{enumerate}
  \item $n=1$: $U_1\in\tau$ zaten verili.
  \item $n=k$ için doğru varsayalım: $V=\bigcap_{i=1}^{k}U_i\in\tau$.
  \item $n=k+1$: $\bigcap_{i=1}^{k+1}U_i = V\cap U_{k+1}$ olup (T2)
        gereği $\tau$'dadır.
\end{enumerate}
Sonsuz kesişimde 3.~adım çalışmaz: tümevarım yalnız sonlu $n$'lere
ulaşır. Şekil~\ref{fig:ch04:kesisim} bunun gerçek bir başarısızlık
olduğunu gösterir: her sonlu kesişim açıkken limit $\{0\}$ açık değildir.
\end{proof}

\begin{figure}[ht]
  \centering
  \begin{tikzpicture}[font=\small, xscale=2.8]
  \draw[-{Stealth}] (-1.35,0) -- (1.5,0) node[right] {$\mathbb{R}$};
  \draw (-1,0.05) -- (-1,-0.05) node[below, font=\scriptsize] {$-1$};
  \draw (0,0.05) -- (0,-0.05);
  \draw (1,0.05) -- (1,-0.05) node[below, font=\scriptsize] {$1$};
  \draw[blue!75, very thick] (-1,0.85) -- (1,0.85);
  \node[blue!75, font=\scriptsize, right] at (1.04,0.85) {$(-1,1)$};
  \draw[blue!55, very thick] (-0.5,0.62) -- (0.5,0.62);
  \node[blue!55, font=\scriptsize, right] at (1.04,0.62) {$(-\tfrac12,\tfrac12)$};
  \draw[blue!40, very thick] (-0.25,0.39) -- (0.25,0.39);
  \node[blue!40, font=\scriptsize, right] at (1.04,0.39) {$(-\tfrac14,\tfrac14)$};
  \node at (0,0.22) {$\vdots$};
  \fill[red!70!black] (0,0) circle (1.2pt);
  \node[red!70!black, below=7pt, font=\scriptsize] at (0,-0.02)
    {$\bigcap_n \bigl(-\tfrac1n,\tfrac1n\bigr)=\{0\}$ --- açık değil};
\end{tikzpicture}

  \caption{(T2)'nin sonlu kesişimle sınırlı olmasının nedeni: iç içe
    $(-\tfrac1n,\tfrac1n)$ açık aralıklarının sonsuz kesişimi tek
    noktaya çöker ve $\{0\}$ standart topolojide açık değildir.}
  \label{fig:ch04:kesisim}
\end{figure}

\begin{theorem}[Baz Teoremi]
\label{thm:basis}
$\mathcal{B} \subseteq \mathcal{P}(X)$ ailesi \emph{(B1)--(B2)} koşullarını sağlıyorsa,
\[
  \tau_{\mathcal{B}}
  = \bigl\{U \subseteq X : \forall x \in U,\, \exists B \in \mathcal{B},\,
    x \in B \subseteq U\bigr\}
\]
bir topoloji olur ve $\mathcal{B}$ bu topolojiye baz oluşturur.
\end{theorem}

\begin{proof}[Rehberli Kanıt]
\textbf{Strateji:} $\tau_{\mathcal{B}}$'nin üç aksiyomu sağladığını, her
adımda yalnız (B1)--(B2)'yi kullanarak doğrularız.
\begin{enumerate}
  \item \textbf{(T1):} $\emptyset\in\tau_{\mathcal B}$, çünkü ``her
        $x\in\emptyset$ için\ldots'' koşulu boş yere doğrudur.
        $X\in\tau_{\mathcal B}$: $x\in X$ verilsin; (B1) gereği
        $x\in\bigcup\mathcal B$, yani $x\in B\subseteq X$ olan bir
        $B\in\mathcal B$ vardır.
  \item \textbf{(T2):} $U,V\in\tau_{\mathcal B}$ ve $x\in U\cap V$
        olsun. Tanım gereği $x\in B_U\subseteq U$ ve
        $x\in B_V\subseteq V$ bulunur. (B2), $x\in B_3\subseteq
        B_U\cap B_V\subseteq U\cap V$ veren bir $B_3$ sağlar; $x$
        keyfî olduğundan $U\cap V\in\tau_{\mathcal B}$.
  \item \textbf{(T3):} $x\in\bigcup_\alpha U_\alpha$ ise
        $x\in U_{\alpha_0}$ olan bir indis vardır; $U_{\alpha_0}$'ın
        tanımından gelen $B$, $x\in B\subseteq
        U_{\alpha_0}\subseteq\bigcup_\alpha U_\alpha$ koşulunu da sağlar.
  \item $\mathcal B$'nin $\tau_{\mathcal B}$'ye baz olması,
        $\tau_{\mathcal B}$'nin tanımının Tanım~\ref{def:basis}'teki
        koşulu birebir içermesindendir.
\end{enumerate}
\end{proof}

\noindent Bu teorem, \texttt{topology\_from\_basis}'in döndürdüğü
ailenin gerçekten topoloji olduğunun matematiksel güvencesidir;
fonksiyon ayrıca (B1)--(B2)'yi sağlamayan girdiyi
\texttt{BasisConstructionError} ile reddeder.

\begin{theorem}[Karşılaştırma]
\label{thm:comparison}
$X$ üzerinde iki topoloji $\tau_1 \subseteq \tau_2$ ise $\tau_1$ \textbf{kaba (coarser)},
$\tau_2$ \textbf{ince (finer)} topoloji olarak adlandırılır. Ayrık topoloji en ince,
indirgenmiş topoloji en kaba topolojidir.
\end{theorem}

\begin{proof}[Rehberli Kanıt]
Herhangi bir $\tau$ için $\tau\subseteq\mathcal{P}(X)=\tau_{\mathrm{disc}}$
topoloji tanımının kendisinden gelir (aksiyom gerekmez); (T1) aksiyomu da
$\{\emptyset,X\}=\tau_{\mathrm{ind}}\subseteq\tau$ verir.
\end{proof}

\begin{figure}[ht]
  \centering
  \begin{tikzpicture}[font=\small,
    kutu/.style={draw, rounded corners, align=center, inner sep=8pt}]
  \node[kutu, fill=orange!6] (kaba) at (0,0)
    {$\tau_{\mathrm{ind}}$ (kaba)\\[3pt] $\emptyset,\quad \{a,b\}$};
  \node[kutu, fill=blue!6] (ince) at (6.4,0)
    {$\tau_{\mathrm{disc}}$ (ince)\\[3pt] $\emptyset,\ \{a\},\ \{b\},\ \{a,b\}$};
  \draw[-{Stealth[length=3mm]}, thick]
    (kaba) -- node[above] {$\subseteq$} node[below] {açık küme eklenir} (ince);
\end{tikzpicture}

  \caption{Aynı $X=\{a,b\}$ üzerinde iki uç: indirgenmiş topoloji (kaba)
    ile ayrık topoloji (ince). Açık küme eklemek topolojiyi inceltir;
    tüm topolojiler bu iki uç arasında yaşar.}
  \label{fig:ch04:kabaince}
\end{figure}

\noindent\texttt{pytop} bu uçları \texttt{tags} ile işaretler:
\texttt{indiscrete\_topology} nesnesi \texttt{indiscrete},
\texttt{discrete\_topology} nesnesi \texttt{discrete} etiketi taşır.

%% ============================================================
\section{Algoritmalar}
%% ============================================================

\subsection{Bazdan Topoloji Üretimi}

\textbf{Problem:} $X$ ve $\mathcal{B} \subseteq \mathcal{P}(X)$ verilsin;
$\tau_{\mathcal{B}}$'yi hesapla.

\begin{algorithm}
\caption{Bazdan Topoloji Üretimi}
\label{alg:basis_to_topology}
\begin{algorithmic}[1]
\Require $X$ (taşıyıcı), $\mathcal{B}$ (baz ailesi)
\Ensure $\tau$ (kapalı topoloji)
\State $\tau \leftarrow \{\emptyset, X\}$
\For{each non-empty $\mathcal{S} \subseteq \mathcal{B}$}
    \State $\tau.\text{add}\bigl(\bigcup_{B \in \mathcal{S}} B\bigr)$
\EndFor
\State \Call{CloseUnderIntersection}{$\tau$}
\Return $\tau$
\end{algorithmic}
\end{algorithm}

\begin{algorithm}
\caption{Kesişim Altında Kapatma}
\begin{algorithmic}[1]
\State $\textit{changed} \leftarrow \textbf{true}$
\While{$\textit{changed}$}
    \State $\textit{changed} \leftarrow \textbf{false}$
    \For{each $U, V \in \tau$}
        \If{$U \cap V \notin \tau$}
            \State $\tau.\text{add}(U \cap V)$;\quad $\textit{changed} \leftarrow \textbf{true}$
        \EndIf
    \EndFor
\EndWhile
\end{algorithmic}
\end{algorithm}

\noindent\textbf{Karmaşıklık:}
\begin{itemize}
  \item En kötü durum: $O(|\mathcal{B}| \cdot 2^{|\mathcal{B}|})$
        (tüm alt-ailelerin birleşimi)
  \item Pratik (çakışmayan baz elemanları): $O(|\tau| \cdot |\mathcal{B}|)$
  \item Uzay: $O(|\tau|)$ açık küme saklar
\end{itemize}

\subsection{İz Sürme: Küçük Bir Girdiyle Adım Adım}

$X=\{1,2,3\}$, $\mathcal{B}=\{\{1\},\{2,3\}\}$ girdisiyle
Algoritma~\ref{alg:basis_to_topology}:

\begin{table}[ht]
\centering
\begin{tabular}{clll}
\toprule
\textbf{Adım} & $\mathcal{S}$ (alt-aile) & \textbf{Eklenen birleşim} & $\tau$ \textbf{(o ana dek)} \\
\midrule
0 & --- & --- & $\{\emptyset, X\}$ \\
1 & $\{\{1\}\}$ & $\{1\}$ & $\{\emptyset, X, \{1\}\}$ \\
2 & $\{\{2,3\}\}$ & $\{2,3\}$ & $\{\emptyset, X, \{1\}, \{2,3\}\}$ \\
3 & $\{\{1\},\{2,3\}\}$ & $\{1\}\cup\{2,3\}=X$ & değişmez \\
4 & kesişim kapanışı & $\{1\}\cap\{2,3\}=\emptyset$ & değişmez \\
\bottomrule
\end{tabular}
\caption{İz sürme: $|\tau|=4$; döngü yeni küme üretmeyince durur.}
\end{table}

\subsection{Alt-Bazdan Topoloji Üretimi}

$\mathcal{S}$ verilsin; önce $\mathcal{S}$'nin sonlu kesişimlerinden baz oluştur,
ardından Algoritma~\ref{alg:basis_to_topology}'yi uygula.

\noindent\textbf{Karmaşıklık:} $O(|\mathcal{S}|^k \cdot 2^{|\mathcal{S}|^k})$ —
$k$: kesişim derinliği.

%% ============================================================
\section{pytop API}
%% ============================================================

\subsection{Sınıflar}

\begin{lstlisting}[language=Python, caption={Temel sınıf importları}]
from pytop import TopologicalSpace, FiniteTopologicalSpace
\end{lstlisting}

\begin{table}[h]
\centering
\begin{tabular}{lll}
\toprule
\textbf{Sınıf} & \textbf{Taşıyıcı} & \textbf{Topoloji} \\
\midrule
\texttt{TopologicalSpace} & \texttt{Any} & \texttt{Any} (None olabilir) \\
\texttt{FiniteTopologicalSpace} & sonlu \texttt{frozenset} & \texttt{frozenset[frozenset]} \\
\bottomrule
\end{tabular}
\caption{Temel uzay sınıfları}
\end{table}

\noindent\textbf{Alanlar:} \texttt{carrier} (altta yatan küme), \texttt{topology}
(açık kümeler), \texttt{tags} (özellik etiketleri), \texttt{metadata} (ek bilgiler).

\subsection{Oluşturucular}

\begin{lstlisting}[language=Python, caption={Topoloji oluşturucu importları}]
from pytop import (
    make_topology, discrete_topology, indiscrete_topology,
    cofinite_topology, sierpinski_space,
    topology_from_basis, topology_from_subbasis,
)
\end{lstlisting}

\begin{table}[h]
\centering
\begin{tabular}{lll}
\toprule
\textbf{Fonksiyon} & \textbf{İmza} & \textbf{Açıklama} \\
\midrule
\texttt{make\_topology} & \texttt{(carrier, *open\_sets)} & Açık küme listesiyle inşa \\
\texttt{discrete\_topology} & \texttt{(*elements)} & Ayrık topoloji \\
\texttt{indiscrete\_topology} & \texttt{(*elements)} & İndirgenmiş topoloji \\
\texttt{cofinite\_topology} & \texttt{(*elements)} & Kosonlu topoloji \\
\texttt{sierpinski\_space} & \texttt{()} & Sierpiński uzayı \\
\texttt{topology\_from\_basis} & \texttt{(carrier, basis)} & Bazdan topoloji üret \\
\texttt{topology\_from\_subbasis} & \texttt{(carrier, subbasis)} & Alt-bazdan topoloji üret \\
\bottomrule
\end{tabular}
\caption{Topoloji oluşturucu fonksiyonlar}
\end{table}

\subsection{Tag Sistemi}

\texttt{pytop} nesneleri \texttt{tags} kümesi aracılığıyla topolojik özellikler taşır.
Örnek etiketler: \texttt{"finite"}, \texttt{"discrete"}, \texttt{"t0"}, \texttt{"t1"},
\texttt{"hausdorff"}, \texttt{"compact"}, \texttt{"connected"}, \texttt{"metric"}.

\subsection{Tag Gerekçeleri}

Etiketler, kurucunun inşa anında \emph{garanti ettiği} gerçeklerdir:

\begin{table}[ht]
\centering
\begin{tabular}{lp{4.2cm}p{5.6cm}}
\toprule
\textbf{Kurucu} & \textbf{Etiketler} & \textbf{Gerekçe} \\
\midrule
\texttt{sierpinski\_space()} & compact, connected, finite, t0 &
  Sonlu $\Rightarrow$ kompakt; $\{1\}$ tek yönlü ayırır $\Rightarrow$ T0
  (T1 değil); ayrık parçalanış yok $\Rightarrow$ bağlantılı \\
\texttt{discrete\_topology(1,2,3)} & discrete, finite, hausdorff,
  metrizable, normal, regular &
  Her tekil açık $\Rightarrow$ tüm ayrılma aksiyomları; 0--1 metriği
  ayrık topolojiyi üretir \\
\texttt{indiscrete\_topology(1,2,3)} & compact, connected, finite,
  indiscrete & Açık yalnız $\emptyset$ ve $X$: parçalanamaz \\
\texttt{cofinite\_topology('a','b','c')} & cofinite, compact, finite,
  t1 & Tekiller kapalı $\Rightarrow$ T1 \\
\texttt{make\_topology(\ldots)},
  \texttt{finite\_chain\_space(n)} & finite &
  Özellik çıkarımı yapılmaz; yalnız sonluluk işaretlenir \\
\bottomrule
\end{tabular}
\caption{Pilot bölümdeki kurucuların etiket gerekçeleri}
\end{table}

\noindent Etiketler \emph{eksiksiz değildir}:
\texttt{cofinite\_topology('a','b','c')} aslında ayrıktır ve
Hausdorff'tur ($2^3=8$ açık), ama \texttt{discrete}/\texttt{hausdorff}
etiketi taşımaz --- \texttt{is\_t2} yüklemi yine \texttt{true} döner.
Kesin sorgu için Bölüm~\ref{ch:separation}'daki \texttt{is\_*}
yüklemlerini kullanın; etiket, yüklemin yerine geçmez.

%% ============================================================
\section{Örnekler}
%% ============================================================

\subsection{Örnek 1: Sierpiński Uzayı}

\begin{lstlisting}[language=Python, caption={Sierpiński uzayı}]
from pytop import sierpinski_space

s = sierpinski_space()
print("Tasiyici:", s.carrier)
print("Topoloji:", sorted(str(t) for t in s.topology))
print("Etiketler:", sorted(s.tags))
\end{lstlisting}

\begin{lstlisting}[caption={Çıktı}]
Tasiyici: frozenset({0, 1})
Topoloji: ['frozenset()', 'frozenset({0, 1})', 'frozenset({1})']
Etiketler: ['compact', 'connected', 'finite', 't0']
\end{lstlisting}

\paragraph{Ne oldu?}
Çıktıdaki üç satırı tek tek okuyalım. \texttt{Tasiyici} satırı
$X=\{0,1\}$'i verir. \texttt{Topoloji} satırındaki üç küme tam olarak
Tanım~\ref{def:topological_space}'in istediği yapıdır: $\emptyset$ ve
$X$ (T1), tek ek açık küme $\{1\}$; $\{1\}\cap X=\{1\}$ ve
$\{1\}\cup X=X$ zaten listede olduğundan (T2)--(T3) de sağlanır.
\texttt{Etiketler} satırında \texttt{t0} var ama \texttt{t1} yok:
$1$'i içerip $0$'ı dışlayan açık küme vardır ($\{1\}$), fakat $0$'ı
içerip $1$'i dışlayan açık küme yoktur --- T0 sağlanıp T1'in
sağlanmamasının kaynağı budur (Bölüm~\ref{ch:separation}, Örnek~1'de
yüklemlerle test edilir). Şekil~\ref{fig:ch04:sierpinski} açık küme
yapısını gösterir.

\begin{figure}[ht]
  \centering
  \begin{tikzpicture}[font=\small]
  \fill (0,0) circle (2pt) node[below=3pt] {$0$};
  \fill (2,0) circle (2pt) node[below=3pt] {$1$};
  \draw[blue!70!black, thick] (1,0.05) ellipse [x radius=2.1, y radius=1.0];
  \node[blue!70!black] at (1,1.3) {$X=\{0,1\}$};
  \draw[violet!80!black, thick] (2,0.05) ellipse [x radius=0.62, y radius=0.5];
  \node[violet!80!black] at (2.05,0.82) {$\{1\}$};
  \draw[red!70!black, dashed] (0,0.05) ellipse [x radius=0.62, y radius=0.5];
  \node[red!70!black] at (-0.15,-1.05) {$\{0\}\notin\tau$ (açık değil)};
\end{tikzpicture}

  \caption{Sierpiński uzayının açıkları: $\emptyset$, $\{1\}$ ve $X$.
    Kesikli bölge $\{0\}$'ın açık \emph{olmadığını} vurgular; bu
    asimetri uzayı T0 yapar ama T1 yapmaz.}
  \label{fig:ch04:sierpinski}
\end{figure}

\subsection{Örnek 2: Ayrık Topoloji}

\begin{lstlisting}[language=Python, caption={Ayrık topoloji}]
from pytop import discrete_topology

d = discrete_topology(1, 2, 3)
print("|tau| =", len(d.topology))
print("Etiketler:", sorted(d.tags))
\end{lstlisting}

\begin{lstlisting}[caption={Çıktı}]
|tau| = 8
Etiketler: ['discrete', 'finite', 'hausdorff', 'metrizable', 'normal', 'regular']
\end{lstlisting}

\paragraph{Ne oldu?}
$n=3$ eleman için $|\tau|=2^3=8$: her alt küme açıktır. Her tekil küme
açık olduğundan herhangi iki nokta kendi tekil komşuluklarıyla ayrılır
--- \texttt{hausdorff}, \texttt{normal}, \texttt{regular} etiketlerinin
kaynağı budur. \texttt{metrizable}, 0--1 metriğinden gelir:
$d(x,y)=1$ $(x\neq y)$ metriği tam olarak ayrık topolojiyi üretir.
Listede \texttt{compact} yok; oysa uzay sonlu olduğu için kompakttır ve
\texttt{is\_compact} yüklemi \texttt{true} döner --- fark için Tag
Gerekçeleri tablosuna bakın.

\subsection{Örnek 3: \texttt{make\_topology} ile Manuel İnşa}

\begin{lstlisting}[language=Python, caption={Manuel topoloji inşası}]
from pytop import make_topology

sp = make_topology({1, 2, 3}, {1}, {2, 3})
print("|tau| =", len(sp.topology))
print("Acik kumeler:", sorted(str(t) for t in sp.topology))
\end{lstlisting}

\begin{lstlisting}[caption={Çıktı}]
|tau| = 4
Acik kumeler: ['frozenset()', 'frozenset({1, 2, 3})',
               'frozenset({1})', 'frozenset({2, 3})']
\end{lstlisting}

\paragraph{Ne oldu?}
\texttt{make\_topology} verilen $\{1\}$ ve $\{2,3\}$ açıklarına yalnızca
$\emptyset$ ve $X$'i ekledi; $|\tau|=4$. Bu örnekte sonuç geçerli bir
topolojidir, çünkü verilen iki küme ayrıktır: birleşimleri $X$,
kesişimleri $\emptyset$ zaten ailededir. \texttt{make\_topology} bu
kapanışları \emph{hesaplamaz} --- aşağıdaki kutu, kapalı olmayan bir
aile verildiğinde ne olduğunu gösterir.

\begin{dikkat}
\texttt{make\_topology} verdiğiniz aileye yalnızca $\emptyset$ ve $X$'i
ekler; birleşim/kesişim kapanışını \textbf{hesaplamaz} ve aksiyomları
\textbf{denetlemez}. Aşağıdaki çağrıda $\{1\}\cup\{2\}=\{1,2\}$ ailede
yoktur; dönen nesne (T3)'ü ihlal eden, geçersiz bir ``topoloji''dir.
Kapanışın hesaplanmasını ve ailenin denetlenmesini istiyorsanız
\texttt{topology\_from\_basis} kullanın --- o, baz koşullarını
sağlamayan aileyi \texttt{BasisConstructionError} ile reddeder.
\end{dikkat}

\begin{lstlisting}[language=Python, caption={Kapanış hesaplanmaz --- sorumluluk kullanıcıda}]
from pytop import make_topology, topology_from_basis

sessiz = make_topology({1, 2, 3}, {1}, {2})   # denetlemez!
print("make_topology:", sorted(sorted(t) for t in sessiz.topology))

try:
    topology_from_basis({1, 2, 3}, [{1}, {2}])  # B1 ihlali
except Exception as e:
    print("topology_from_basis:", type(e).__name__)
\end{lstlisting}

\begin{lstlisting}[style=output, caption={Çıktı}]
make_topology: [[], [1], [1, 2, 3], [2]]
topology_from_topology: BasisConstructionError
\end{lstlisting}

\subsection{Örnek 4: Alexandrov Zincir Uzayı}

\begin{lstlisting}[language=Python, caption={Zincir uzayı}]
from pytop import finite_chain_space

c = finite_chain_space(3)
print("Tasiyici:", c.carrier)
print("Topoloji:", sorted(str(t) for t in c.topology))
\end{lstlisting}

\begin{lstlisting}[caption={Çıktı}]
Tasiyici: (1, 2, 3)
Topoloji: ['set()', '{1, 2, 3}', '{1, 2}', '{1}']
\end{lstlisting}

\paragraph{Ne oldu?}
Çıktıdaki açıklar tam bir ``önek merdiveni''dir:
$\emptyset\subset\{1\}\subset\{1,2\}\subset\{1,2,3\}$. Alexandrov zincir
uzayında $1$ her boş olmayan açıkta yer alır (``en açık'' nokta), $3$
yalnız $X$'te görünür. Öneklerin kesişimi ve birleşimi yine önek
olduğundan (T2)--(T3) otomatik sağlanır.

\subsection{Örnek 5: Bazdan Topoloji Üretimi}

\begin{lstlisting}[language=Python, caption={Bazdan topoloji}]
from pytop import topology_from_basis

b = [{1}, {2, 3}, {4}]
ts = topology_from_basis({1, 2, 3, 4}, b)
print("Baz:", [set(x) for x in b])
print("|tau| =", len(ts.topology))
print("Topoloji:", sorted(str(t) for t in ts.topology))
\end{lstlisting}

\begin{lstlisting}[caption={Çıktı}]
Baz: [{1}, {2, 3}, {4}]
|tau| = 8
Topoloji: ['set()', '{1, 2, 3, 4}', '{1, 2, 3}', '{1, 4}',
           '{1}', '{2, 3, 4}', '{2, 3}', '{4}']
\end{lstlisting}

\paragraph{Ne oldu?}
Baz $\{\{1\},\{2,3\},\{4\}\}$ bir bölüntüdür: elemanları ikişer ikişer
ayrıktır. (B2) koşulu boş yere sağlanır (kesişimler boş olduğundan
kontrol edilecek nokta yoktur) ve topoloji tüm alt-aile
birleşimlerinden oluşur: $2^3=8$ açık küme. Çıktıdaki 8 küme, üç baz
elemanının $2^3$ alt kümesinin birleşimleriyle birebir eşleşir.

\subsection{Örnek 6: Gerçek Doğru}

\begin{lstlisting}[language=Python, caption={Gerçek doğru}]
from pytop import real_line_metric

rl = real_line_metric()
print("Tasiyici:", rl.carrier)
print("Etiketler:", sorted(rl.tags))
\end{lstlisting}

\begin{lstlisting}[caption={Çıktı}]
Tasiyici: R
Etiketler: ['complete', 'connected', 'first_countable', 'hausdorff',
            'infinite', 'lindelof', 'metric', 'not_compact',
            'path_connected', 'second_countable', 'separable', 't0', 't1', 'uncountable']
\end{lstlisting}

\paragraph{Ne oldu?}
\texttt{Tasiyici: R} satırı taşıyıcının \emph{simgesel} olduğunu söyler:
gerçek doğrunun noktaları bellekte tutulmaz, \texttt{topology=None}'dır.
Tüm topolojik bilgi etiketlerde kodlanmıştır: \texttt{connected},
\texttt{hausdorff}, \texttt{second\_countable} gibi olumlu özellikler
yanında \texttt{not\_compact} gibi olumsuz bilgi de taşınır
(Heine--Borel: $\mathbb{R}$ sınırlı değildir). Sonlu uzaylardaki
``hesapla ve doğrula'' yaklaşımının yerini burada ``bilinen teoremleri
etiketle'' yaklaşımı alır; metrik tarafı Bölüm~14'te derinleşir.

%% ============================================================
\section{Alıştırmalar}
%% ============================================================

\subsection*{Kodlama Alıştırmaları}

\begin{enumerate}[label=\textbf{K\arabic*.}]
  \item \label{ex:ch04:K1} \texttt{cofinite\_topology('a','b','c')} ile üç
        noktalı kosonlu uzayı kurun; topolojisini ve etiketlerini
        yazdırın, \texttt{is\_t1} ve \texttt{is\_t2} ile test edin.
        Gözlem: sonlu bir kümede kosonlu topoloji ayrık topolojiyle
        çakışır. T1 olup Hausdorff olmayan bir örnek için taşıyıcının
        neden sonsuz olması gerektiğini bir cümleyle açıklayın
        (krş.\ Bölüm~\ref{ch:separation}, Örnek~2).
        \ipucu{$|\tau|=2^3$ çıkacak; anahtar, Bölüm~\ref{ch:separation}'daki
        ``Sonlu T1 $\iff$ Ayrık'' teoremidir.}
        (\hyperref[sol:ch04:K1]{çözüm})

  \item \label{ex:ch04:K2} \texttt{topology\_from\_subbasis(\{1,2,3,4\},
        [\{1,2\},\{3,4\},\{2,3\}])} çağrısının ürettiği topolojinin kaç
        açık küme içerdiğini bulun ve açık kümeleri listeleyin.
        \ipucu{Önce alt-baz çiftlerinin kesişimlerini elle listeleyin
        ($\{2\}$, $\{3\}$, $\emptyset$); sonra birleşimleri sayın.}
        (\hyperref[sol:ch04:K2]{çözüm})

  \item \label{ex:ch04:K3} \texttt{finite\_chain\_space(5)} ile bir
        Alexandrov-5 zinciri oluşturun; topolojinin kaç açık küme
        içerdiğini ve hangi elemanın ``en açık'' nokta olduğunu bulun.
        \ipucu{Açıklar önek yapısındadır; ``en açık'' nokta her boş
        olmayan açıkta bulunandır.}
        (\hyperref[sol:ch04:K3]{çözüm})
\end{enumerate}

\subsection*{Teori Alıştırmaları}

\begin{enumerate}[label=\textbf{T\arabic*.}]
  \item \label{ex:ch04:T1} $X=\{1,2,3\}$ üzerinde kaç farklı topoloji
        tanımlanabilir? Bu sayıyı T1--T3 aksiyomlarını doğrudan kontrol
        ederek hesaplamak yerine neden baz teoremini kullanmak daha
        pratiktir?
        \ipucu{Aday aile sayısı $2^{2^3}=256$'dır; cevap 29. Baz teoremi
        adayları üretken küçük ailelere indirger.}
        (\hyperref[sol:ch04:T1]{çözüm})

  \item \label{ex:ch04:T2} Ayrık topolojinin her zaman diğer
        topolojilerden daha ince, indirgenmiş topolojinin ise daha kaba
        olduğunu kanıtlayın. Kanıt için T1--T3 aksiyomlarından hangisi
        gereklidir?
        \ipucu{Ayrıklık için aksiyom gerekmez
        ($\tau\subseteq\mathcal{P}(X)$ tanım gereğidir); indirgenmişlik
        için yalnız (T1) yeter.}
        (\hyperref[sol:ch04:T2]{çözüm})
\end{enumerate}

\chapter{Yüklemler ve Altküme Operatörleri}
\label{ch:predicates_operators}

%% ============================================================
\section{Konu}
%% ============================================================

\subsection{Altküme Yüklemleri}

$(X, \tau)$ bir topolojik uzay ve $A \subseteq X$ olsun.

\begin{table}[h]
\centering
\begin{tabular}{ll}
\toprule
\textbf{Yüklem} & \textbf{Tanım} \\
\midrule
Açık (open) & $A \in \tau$ \\
Kapalı (closed) & $X \setminus A \in \tau$ \\
Clopen & $A \in \tau$ ve $X \setminus A \in \tau$ \\
Yoğun (dense) & Her boş olmayan $U \in \tau$ için $U \cap A \neq \emptyset$ \\
Hiçbiryerde-yoğun & $\mathrm{int}(\overline{A}) = \emptyset$ \\
\bottomrule
\end{tabular}
\caption{Altküme yüklemleri}
\end{table}

\subsection{Altküme Operatörleri}

\begin{table}[h]
\centering
\begin{tabular}{lll}
\toprule
\textbf{Operatör} & \textbf{Gösterim} & \textbf{Tanım} \\
\midrule
Kapanış & $\overline{A} = \mathrm{cl}(A)$ & $A$'yı içeren en küçük kapalı küme \\
İç & $A^\circ = \mathrm{int}(A)$ & $A$'ya dahil en büyük açık küme \\
Sınır & $\partial A = \mathrm{bd}(A)$ & $\overline{A} \setminus A^\circ$ \\
Türev kümesi & $A' = d(A)$ & $A$'nın birikme noktaları kümesi \\
\bottomrule
\end{tabular}
\caption{Altküme operatörleri}
\end{table}

\noindent\textbf{Birikme noktası:} $x \in A'$ $\iff$ her $U \ni x$ açık için
$U \cap (A \setminus \{x\}) \neq \emptyset$.

\begin{sezgi}
Bir $A$ kümesini, $X$ üzerinde duran bir ``boya lekesi'' gibi düşünün.
\textbf{İç} $A^\circ$, lekenin tamamen içine sığabildiğiniz açık bir
komşuluğunuzun olduğu noktalardır --- ``rahatça içerideyim'' bölgesi.
\textbf{Kapanış} $\mathrm{cl}(A)$, lekeye dokunan her noktayı ekler:
leke artı onun ``gölgesi''. \textbf{Sınır} $\partial A$ ise tam kenardır:
her açık komşuluğu hem $A$'dan hem $X\setminus A$'dan kesen noktalar.
Bu üçlü, $X$'i ayrık üç parçaya böler: iç, sınır ve \textbf{dış}
$\mathrm{ext}(A) = (X\setminus A)^\circ$.
\end{sezgi}

\begin{figure}[ht]
  \centering
  \begin{tikzpicture}[font=\small,
    nokta/.style={circle, draw, minimum size=7mm, inner sep=1pt},
    ic/.style={nokta, fill=green!18},
    sin/.style={nokta, fill=orange!22},
    dis/.style={nokta, fill=blue!10}]
  % X = {1,2,3,4,5}; A = {1,2,3}
  % int(A)={1,2} (yesil), bd(A)={3} (turuncu), ext(A)={4,5} (mavi)
  \node[ic]  (n1) at (0,0)   {$1$};
  \node[ic]  (n2) at (1.5,0) {$2$};
  \node[sin] (n3) at (3,0)   {$3$};
  \node[dis] (n4) at (4.5,0) {$4$};
  \node[dis] (n5) at (6,0)   {$5$};

  % A kapanis cercevesi cl(A) = {1,2,3}
  \draw[red!70!black, thick, rounded corners]
    ($(n1.north west)+(-0.25,0.55)$) rectangle ($(n3.south east)+(0.25,-0.55)$);
  \node[red!70!black] at (1.5,1.05) {$\mathrm{cl}(A)=\{1,2,3\}$};

  % int(A) cercevesi
  \draw[green!55!black, thick, dashed, rounded corners]
    ($(n1.north west)+(-0.15,0.3)$) rectangle ($(n2.south east)+(0.15,-0.3)$);

  % Etiketler
  \node[green!45!black, font=\scriptsize] at (0.75,-1.0) {iç: $A^\circ=\{1,2\}$};
  \node[orange!75!black, font=\scriptsize] at (3,-1.0) {sınır: $\partial A=\{3\}$};
  \node[blue!65!black, font=\scriptsize] at (5.25,-1.0) {dış: $\mathrm{ext}(A)=\{4,5\}$};

  \node[font=\scriptsize, align=center] at (3,-1.7)
    {$X = A^\circ \;\cup\; \partial A \;\cup\; \mathrm{ext}(A)$ (ayrık parçalanış)};
\end{tikzpicture}

  \caption{$A=\{1,2,3\}$ için iç $A^\circ=\{1,2\}$ (yeşil), sınır
    $\partial A=\{3\}$ (turuncu) ve dış $\mathrm{ext}(A)=\{4,5\}$ (mavi);
    kapanış $\mathrm{cl}(A)$ kırmızı çerçevedir. $X$ bu üç parçaya ayrık olarak bölünür.}
  \label{fig:ch05:ic_kapanis_sinir}
\end{figure}

\subsection{Komşuluk Sistemi}

\begin{definition}[Komşuluk Sistemi]
$x \in X$ için:
\[
  N(x) = \{N \subseteq X : \exists U \in \tau,\, x \in U \subseteq N\}
\]
\end{definition}

\noindent\textbf{Komşuluk Aksiyomları (N1--N4):}
\begin{enumerate}
  \item[(N1)] $x \in N$, her $N \in N(x)$ için
  \item[(N2)] $X \in N(x)$
  \item[(N3)] $N_1, N_2 \in N(x) \Rightarrow N_1 \cap N_2 \in N(x)$
  \item[(N4)] $N \in N(x)$, $N \subseteq M \Rightarrow M \in N(x)$
\end{enumerate}

%% ============================================================
\section{Teoremler}
%% ============================================================

\begin{theorem}[Kuratowski Kapanış Aksiyomları]
\label{thm:kuratowski}
$\mathrm{cl}: \mathcal{P}(X) \to \mathcal{P}(X)$ operatörü şu dört özelliği sağlar:
\begin{enumerate}
  \item[(K1)] $\mathrm{cl}(\emptyset) = \emptyset$
  \item[(K2)] $A \subseteq \mathrm{cl}(A)$
  \item[(K3)] $\mathrm{cl}(\mathrm{cl}(A)) = \mathrm{cl}(A)$
  \item[(K4)] $\mathrm{cl}(A \cup B) = \mathrm{cl}(A) \cup \mathrm{cl}(B)$
\end{enumerate}
Bu dört özelliği sağlayan her operatör bir topoloji tanımlar.
\end{theorem}

\begin{figure}[ht]
  \centering
  \begin{tikzpicture}[font=\small, node distance=7mm,
    aks/.style={draw, rounded corners, fill=blue!6, inner sep=6pt,
                text width=4.6cm, align=center}]
  \node[aks] (k1) {\textbf{(K1)}\\ $\mathrm{cl}(\emptyset)=\emptyset$};
  \node[aks, right=of k1] (k2) {\textbf{(K2)}\\ $A \subseteq \mathrm{cl}(A)$};
  \node[aks, below=of k1] (k3) {\textbf{(K3)}\\ $\mathrm{cl}(\mathrm{cl}(A))=\mathrm{cl}(A)$};
  \node[aks, below=of k2] (k4) {\textbf{(K4)}\\ $\mathrm{cl}(A\cup B)=\mathrm{cl}(A)\cup \mathrm{cl}(B)$};

  \node[draw, rounded corners, fill=green!10, inner sep=6pt,
        below=10mm of $(k3)!0.5!(k4)$, text width=7.2cm, align=center] (sonuc)
    {Bu dört aksiyomu sağlayan her $\mathrm{cl}$ operatörü\\
     \emph{tek bir} topoloji belirler: $\tau=\{X\setminus \mathrm{cl}(A)\}$.};

  \draw[-{Stealth[length=2.4mm]}, thick] (k3.south) -- (k3.south |- sonuc.north);
  \draw[-{Stealth[length=2.4mm]}, thick] (k4.south) -- (k4.south |- sonuc.north);
\end{tikzpicture}

  \caption{Kuratowski aksiyomları K1--K4 ve sonucu: bu dördünü sağlayan her
    $\mathrm{cl}$ operatörü tek bir topoloji belirler.}
  \label{fig:ch05:kuratowski}
\end{figure}

\begin{proof}[İspat eskizi]
Kapanış, $A$'yı içeren tüm kapalı kümelerin kesişimidir:
$\mathrm{cl}(A) = \bigcap\{C : C \text{ kapalı},\, A\subseteq C\}$.
\textbf{(K1)} $\emptyset$ zaten kapalı, en küçük kapalı üst-küme kendisidir.
\textbf{(K2)} her $C\supseteq A$ içerdiğinden $A\subseteq \mathrm{cl}(A)$.
\textbf{(K3)} $\mathrm{cl}(A)$ kapalıdır; kapalı bir kümenin kapanışı kendisidir
(idempotentlik). \textbf{(K4)} $\subseteq$ yönü monotonluktan, $\supseteq$ yönü
$\mathrm{cl}(A)\cup\mathrm{cl}(B)$'nin $A\cup B$'yi içeren kapalı bir küme olmasından gelir.
Sonlu uzaylarda dördü de \texttt{kuratowski\_closure\_check} ile doğrulanır (Örnek~2).
\end{proof}

\begin{theorem}[İç-Kapanış Dualitesi]
\label{thm:interior_closure_duality}
Her $A \subseteq X$ için:
\[
  A^\circ = X \setminus \mathrm{cl}(X \setminus A), \qquad
  \mathrm{cl}(A) = X \setminus \mathrm{int}(X \setminus A).
\]
\end{theorem}

\begin{proof}[İspat eskizi]
$A^\circ$ tanım gereği $A$'ya dahil \emph{en büyük açık} kümedir; $\mathrm{cl}(B)$ ise
$B$'yi içeren \emph{en küçük kapalı} kümedir. Tümleme alma açık~$\leftrightarrow$~kapalı
ve ``en büyük''~$\leftrightarrow$~``en küçük'' ile $\subseteq$ ilişkisini ters çevirir.
$B=X\setminus A$ alın: $X\setminus A$'yı içeren en küçük kapalı kümenin tümleyeni,
$A$'ya dahil en büyük açık kümedir; yani $A^\circ = X\setminus\mathrm{cl}(X\setminus A)$.
İkinci eşitlik $A\mapsto X\setminus A$ ikamesiyle simetrik olarak çıkar.
\end{proof}

\begin{theorem}[Komşuluk Sistemi Round-Trip]
\label{thm:neighborhood_round_trip}
N1--N4'ü sağlayan her $\{N(x)\}_{x \in X}$ ailesi,
$\tau = \{U : \forall x \in U, U \in N(x)\}$ şeklinde tek bir topolojiyi geri üretir.
\end{theorem}

\begin{proof}[İspat eskizi]
İleri yön: $\tau$'dan $N(x)=\{N : \exists U\in\tau,\, x\in U\subseteq N\}$ tanımlayın.
(N1) $x\in U\subseteq N$ olduğundan $x\in N$. (N2) $X\in\tau$ her noktayı içerir.
(N3) $U_1\cap U_2\in\tau$ açık olduğundan kesişim de komşuluktur. (N4) üst-küme almak
tanımı bozmaz. Geri yön: $\tau=\{U : \forall x\in U,\, U\in N(x)\}$ ailesinin topoloji
aksiyomlarını sağladığı ve başlangıç sistemini geri verdiği doğrulanır; teklik bu
eşleşmenin tersinir olmasından gelir.
\end{proof}

%% ============================================================
\section{Algoritmalar}
%% ============================================================

\begin{algorithm}
\caption{Sonlu Kapanış Hesabı}
\begin{algorithmic}[1]
\Require $A \subseteq X$, $\tau$
\Ensure $\mathrm{cl}(A)$
\State $\textit{closed} \leftarrow \{X \setminus U : U \in \tau\}$
\State $\textit{result} \leftarrow X$
\For{each $C \in \textit{closed}$}
    \If{$A \subseteq C$ \textbf{and} $|C| < |\textit{result}|$}
        \State $\textit{result} \leftarrow C$
    \EndIf
\EndFor
\Return $\textit{result}$
\end{algorithmic}
\end{algorithm}

\noindent\textbf{Karmaşıklık:} $O(|\tau| \cdot |X|)$.

\begin{algorithm}
\caption{Türev Kümesi Hesabı}
\begin{algorithmic}[1]
\Require $A \subseteq X$, $\tau$
\Ensure $d(A)$
\State $\textit{acc} \leftarrow \emptyset$
\For{each $x \in X$}
    \If{her açık $U \ni x$ için $U \cap (A \setminus \{x\}) \neq \emptyset$}
        \State $\textit{acc}.\text{add}(x)$
    \EndIf
\EndFor
\Return $\textit{acc}$
\end{algorithmic}
\end{algorithm}

\noindent\textbf{Karmaşıklık:} $O(|X| \cdot |\tau|)$.

%% ============================================================
\section{pytop API}
%% ============================================================

\begin{lstlisting}[language=Python, caption={Subset operatörleri importları}]
from pytop import (
    analyze_predicate,
    closure_of_subset, interior_of_subset, boundary_of_subset,
    derived_set_of_subset,
    is_open_subset, is_closed_subset, is_dense_subset, is_nowhere_dense_subset,
    neighborhood_system, character_at_point, analyze_neighborhood_system,
)
\end{lstlisting}

\begin{table}[h]
\centering
\begin{tabular}{lll}
\toprule
\textbf{Fonksiyon} & \textbf{İmza} & \textbf{Açıklama} \\
\midrule
\texttt{analyze\_predicate} & \texttt{(space, predicate, subset)} & Yüklem sorgula \\
\texttt{closure\_of\_subset} & \texttt{(space, subset)} & Kapanış \\
\texttt{interior\_of\_subset} & \texttt{(space, subset)} & İç \\
\texttt{boundary\_of\_subset} & \texttt{(space, subset)} & Sınır \\
\texttt{derived\_set\_of\_subset} & \texttt{(space, subset)} & Türev kümesi \\
\texttt{neighborhood\_system} & \texttt{(carrier, topology, point)} & $N(x)$ \\
\texttt{character\_at\_point} & \texttt{(carrier, topology, point)} & $\chi(x)$ \\
\texttt{analyze\_neighborhood\_system} & \texttt{(carrier, topology, point=None)} & Tam analiz \\
\bottomrule
\end{tabular}
\caption{Altküme operatör fonksiyonları}
\end{table}

\noindent\textbf{Not:} \texttt{neighborhood\_system} ve ilgili fonksiyonlar \texttt{carrier} (liste)
ve \texttt{topology} (liste) alır. Space nesnesi için: \texttt{list(space.carrier)},
\texttt{list(space.topology)}.

%% ============================================================
\section{Örnekler}
%% ============================================================

\subsection{Örnek 1: Kapanış ve İç — Sierpiński Uzayı}

\begin{lstlisting}[language=Python]
from pytop import sierpinski_space, closure_of_subset, interior_of_subset, boundary_of_subset

s = sierpinski_space()
print("cl({0})  =", closure_of_subset(s, {0}).value)
print("cl({1})  =", closure_of_subset(s, {1}).value)
print("int({0}) =", interior_of_subset(s, {0}).value)
print("int({1}) =", interior_of_subset(s, {1}).value)
print("bd({1})  =", boundary_of_subset(s, {1}).value)
\end{lstlisting}

\begin{lstlisting}[caption={Çıktı}]
cl({0})  = frozenset({0})
cl({1})  = frozenset({0, 1})
int({0}) = frozenset()
int({1}) = frozenset({1})
bd({1})  = frozenset({0})
\end{lstlisting}

$\{0\}$'ın kapanışı kendisidir: $\{0\} = X \setminus \{1\}$ zaten kapalıdır.
$\{1\}$'in kapanışı $X$'tir: $\{1\}$ kapalı değildir.

\subsection{Örnek 2: Türev Kümesi}

\begin{lstlisting}[language=Python]
from pytop import sierpinski_space, derived_set_of_subset, is_nowhere_dense_subset

s = sierpinski_space()
print("d({0}) =", derived_set_of_subset(s, {0}).value)
print("d({1}) =", derived_set_of_subset(s, {1}).value)
print("{0} nowhere dense? =", is_nowhere_dense_subset(s, {0}).status)
\end{lstlisting}

\begin{lstlisting}[caption={Çıktı}]
d({0}) = frozenset()
d({1}) = frozenset({0})
{0} nowhere dense? = true
\end{lstlisting}

$d(\{1\}) = \{0\}$: $0$'ı içeren her $\{0,1\}$, $\{1\} \setminus \{0\} = \{1\}$ ile kesişir.

\subsection{Örnek 3: Komşuluk Sistemi}

\begin{lstlisting}[language=Python]
from pytop import sierpinski_space, neighborhood_system, character_at_point

s = sierpinski_space()
carrier = list(s.carrier)
topology = list(s.topology)
print("N(0) =", neighborhood_system(carrier, topology, 0).value)
print("N(1) =", neighborhood_system(carrier, topology, 1).value)
print("chi(0) =", character_at_point(carrier, topology, 0).value)
print("chi(1) =", character_at_point(carrier, topology, 1).value)
\end{lstlisting}

\begin{lstlisting}[caption={Çıktı}]
N(0) = [['0', '1']]
N(1) = [['1'], ['0', '1']]
chi(0) = 1
chi(1) = 2
\end{lstlisting}

$N(0) = \{X\}$: $0$ yalnızca $X$'te açık. $N(1) = \{\{1\}, X\}$: yerel baz büyüklükleri
$\chi(0)=1$, $\chi(1)=2$.

\subsection{Örnek 4: İç / Kapanış / Sınır / Dış Parçalanışı}

\begin{lstlisting}[language=Python]
from pytop import (make_topology, interior_of_subset, closure_of_subset,
                   boundary_of_subset, exterior_of_subset)

sp = make_topology({1, 2, 3, 4, 5}, {1, 2}, {4, 5}, {1, 2, 4, 5})
A = {1, 2, 3}
print("int(A) =", interior_of_subset(sp, A).value)
print("cl(A)  =", closure_of_subset(sp, A).value)
print("bd(A)  =", boundary_of_subset(sp, A).value)
print("ext(A) =", exterior_of_subset(sp, A).value)
\end{lstlisting}

\begin{lstlisting}[caption={Çıktı}]
int(A) = frozenset({1, 2})
cl(A)  = frozenset({1, 2, 3})
bd(A)  = frozenset({3})
ext(A) = frozenset({4, 5})
\end{lstlisting}

İç $\{1,2\}$, sınır $\{3\}$ ve dış $\{4,5\}$ birlikte $X$'i ayrık olarak kaplar:
$X = A^\circ \cup \partial A \cup \mathrm{ext}(A)$ (Şekil~\ref{fig:ch05:ic_kapanis_sinir}).

\subsection{Örnek 5: Kuratowski Aksiyomlarının Doğrulanması}

\begin{lstlisting}[language=Python]
from pytop import finite_chain_space, kuratowski_closure_check

s = finite_chain_space(3)
report = kuratowski_closure_check(list(s.carrier), list(s.topology))
for k in sorted(report):
    print(f"{k:<13}= {report[k]}")
\end{lstlisting}

\begin{lstlisting}[caption={Çıktı}]
all          = True
empty        = True
extensive    = True
finite_union = True
idempotent   = True
\end{lstlisting}

\texttt{empty}~$\to$~(K1), \texttt{extensive}~$\to$~(K2), \texttt{idempotent}~$\to$~(K3),
\texttt{finite\_union}~$\to$~(K4); \texttt{all} hepsinin sağlandığını özetler
(Şekil~\ref{fig:ch05:kuratowski}).

\subsection{Örnek 6: Yoğun Nokta ve Birikme Kümesi}

\begin{lstlisting}[language=Python]
from pytop import finite_chain_space, closure_of_subset, is_dense_subset, derived_set_of_subset

c = finite_chain_space(4)
print("cl({1})    =", closure_of_subset(c, {1}).value)
print("{1} dense? =", is_dense_subset(c, {1}).status)
print("d({3})     =", derived_set_of_subset(c, {3}).value)
print("cl({4})    =", closure_of_subset(c, {4}).value)
\end{lstlisting}

\begin{lstlisting}[caption={Çıktı}]
cl({1})    = frozenset({1, 2, 3, 4})
{1} dense? = true
d({3})     = frozenset({4})
cl({4})    = frozenset({4})
\end{lstlisting}

$1$ jeneriktir ($\mathrm{cl}(\{1\})=X$, yoğun), $4$ kapalı bir noktadır
($\mathrm{cl}(\{4\})=\{4\}$); Şekil~\ref{fig:ch05:yogunluk}.

\begin{figure}[ht]
  \centering
  \begin{tikzpicture}[font=\small,
    nokta/.style={circle, draw, minimum size=7mm, inner sep=1pt, fill=gray!8},
    yogun/.style={circle, draw, minimum size=7mm, inner sep=1pt, fill=green!22}]
  % finite_chain_space(4): acik kumeler {1} subset {1,2} subset {1,2,3} subset X
  % cl({1}) = X  -> 1 yogun nokta;  cl({4}) = {4} -> 4 kapali nokta
  \node[yogun] (n1) at (0,0)   {$1$};
  \node[nokta] (n2) at (2,0)   {$2$};
  \node[nokta] (n3) at (4,0)   {$3$};
  \node[nokta] (n4) at (6,0)   {$4$};

  \foreach \a/\b in {n1/n2, n2/n3, n3/n4}
    \draw[-{Stealth[length=2.4mm]}, thick] (\a) -- (\b);

  \node[font=\scriptsize] at (1,0.45) {özelleşme};

  % cl({1}) = X kapanis cercevesi
  \draw[green!55!black, thick, rounded corners]
    ($(n1.north west)+(-0.25,0.95)$) rectangle ($(n4.south east)+(0.25,-0.3)$);
  \node[green!45!black, font=\scriptsize] at (3,1.25) {$\mathrm{cl}(\{1\})=X$ — $\{1\}$ yoğun};

  \node[font=\scriptsize, align=center] at (6,-0.95)
    {$\mathrm{cl}(\{4\})=\{4\}$\\ kapalı nokta};
  \node[font=\scriptsize, align=center] at (0,-0.95)
    {jenerik\\ (yoğun)};
\end{tikzpicture}

  \caption{\texttt{finite\_chain\_space(4)} zincirinde özelleşme oku: $1$ jenerik
    (yoğun) noktadır, $\mathrm{cl}(\{1\})=X$; $4$ ise kapalı noktadır, $\mathrm{cl}(\{4\})=\{4\}$.}
  \label{fig:ch05:yogunluk}
\end{figure}

\begin{karsiornek}
``İç küme her zaman kapalıdır'' yanılgısı. Örnek~4'te $A^\circ=\{1,2\}$ açıktır ama
\textbf{kapalı değildir} ($X\setminus\{1,2\}=\{3,4,5\}$ açık değil). İç ve dış daima
açıktır; kapanış ve sınır daima kapalıdır --- bu iki aileyi karıştırmak en sık yapılan
hatadır. Ayrıca $\partial A=\emptyset$ ancak ve ancak $A$ clopen ise olur: Örnek~``Sınır
Hesabı''nda açık $\{1\}$ için $\partial\{1\}=\emptyset$, fakat clopen olmayan $\{1,2\}$
için sınır boş değildir.
\end{karsiornek}

%% ============================================================
\section{Alıştırmalar}
%% ============================================================

\subsection*{Kodlama Alıştırmaları}
\begin{enumerate}[label=\textbf{K\arabic*.}]
  \item İndirgenmiş iki-noktalı topolojide her noktanın komşuluk sistemini hesaplayın.
  \item $X = \{1,2,3,4\}$, $\tau = \{\emptyset, \{1,2\}, \{3,4\}, X\}$ topolojisinde
        $\mathrm{bd}(\{1,2,3\})$ ve $\mathrm{cl}(\{1,2,3\})$ hesaplayın.
  \item \texttt{finite\_chain\_space(4)} zincirinde $\{1,2\}$'nin iç, kapanış ve sınır
        kümelerini bulun.
  \item \texttt{make\_topology(\{1,2,3,4,5\}, \{1,2\}, \{4,5\}, \{1,2,4,5\})} uzayında
        $A=\{1,2,3\}$ için \texttt{interior\_of\_subset}, \texttt{boundary\_of\_subset} ve
        \texttt{exterior\_of\_subset} çalıştırıp $X = A^\circ \cup \partial A \cup
        \mathrm{ext}(A)$ ayrık parçalanışını doğrulayın.
  \item \texttt{finite\_chain\_space(4)} için \texttt{kuratowski\_closure\_check} çağırıp
        \texttt{all} alanının \texttt{True} döndüğünü gösterin; ardından $\{1\}$'in yoğun
        olduğunu \texttt{is\_dense\_subset} ile teyit edin.
\end{enumerate}

\subsection*{Teori Alıştırmaları}
\begin{enumerate}[label=\textbf{T\arabic*.}]
  \item $A^\circ = X \setminus \mathrm{cl}(X \setminus A)$ eşitliğini Kuratowski
        aksiyomlarını kullanarak kanıtlayın.
  \item $A$'nın yoğun olması için gerek ve yeter koşulun $\mathrm{cl}(A) = X$ olduğunu
        gösterin.
  \item Her $A \subseteq X$ için $X = A^\circ \cup \partial A \cup \mathrm{ext}(A)$
        olduğunu ve üç kümenin ikişer ikişer ayrık olduğunu ispatlayın.
        \ipucu{$\mathrm{ext}(A)=(X\setminus A)^\circ$ ve $\partial A=\mathrm{cl}(A)\setminus A^\circ$.}
  \item $A^\circ$ kümesinin daima açık olduğunu, fakat genel olarak kapalı
        \textbf{olmadığını} gösteren bir karşı-örnek verin. \ipucu{Örnek~4'teki $\{1,2\}$.}
\end{enumerate}

\chapter{Ayrılma Aksiyomları}
\label{ch:separation}

%% ============================================================
\section{Konu}
%% ============================================================

Ayrılma aksiyomları, bir topolojik uzaydaki noktaların ve kapalı kümelerin birbirinden
açık kümeler aracılığıyla ne ölçüde ``ayrılabildiğini'' ölçer.

\begin{sezgi}
Ayrılma aksiyomlarını bir mikroskobun çözünürlük kademeleri gibi
düşünün: T0'da iki noktayı \emph{en az bir yönden} ayırt edebilirsiniz;
T1'de her iki yönden; T2'de noktaları çakışmayan iki ayrı ``görüş
alanına'' koyabilirsiniz; T3 ve T4'te artık nokta--kapalı küme ve
kapalı--kapalı çiftleri bile ayrışır. Matematiksel karşılığı: her
aksiyom, $\tau$'nun ``ayırma gücü'' üzerine gittikçe güçlenen bir
varsayımdır ve zincir boyunca her adım bir öncekini gerektirir.
\end{sezgi}

\begin{table}[h]
\centering
\begin{tabular}{lll}
\toprule
\textbf{Aksiyom} & \textbf{İsim} & \textbf{Koşul} \\
\midrule
T0 & Kolmogorov & $\forall x \neq y$: $\exists U \in \tau$ ayıran \\
T1 & Fréchet & $\forall x \neq y$: iki yönlü ayrım \\
T2 & Hausdorff & $\forall x \neq y$: disjoint $U \ni x$, $V \ni y$ \\
T2.5 & Urysohn & disjoint kapanışlı komşuluklar \\
T3 & Regüler & T1 + nokta/kapalı ayırma \\
T3.5 & Tychonoff & T1 + sürekli fonksiyon ayırma \\
T4 & Normal & T1 + disjoint kapalılar ayırma \\
\bottomrule
\end{tabular}
\caption{Ayrılma aksiyomları}
\end{table}

\noindent\textbf{Sıralama:} T4 $\Rightarrow$ T3.5 $\Rightarrow$ T3 $\Rightarrow$ T2.5
$\Rightarrow$ T2 $\Rightarrow$ T1 $\Rightarrow$ T0.

\begin{dikkat}
T3 ve T4'ün tanımı kaynaktan kaynağa değişir: bazı kitaplar
``regüler''i T1 olmadan tanımlar. Bu kılavuzda ve \texttt{pytop}'ta
tablodaki konvansiyon geçerlidir: \textbf{T3 = T1 + regüler,
T4 = T1 + normal}. Fark gerçektir: iki noktalı indirgenmiş uzay
regüler \emph{ve} normaldir (ayrılacak uygun çift yoktur), ama T1
olmadığından T3 de T4 de değildir; Örnekler bölümünde
\texttt{is\_regular}/\texttt{is\_t3} ile doğrulanır.
\end{dikkat}

Şekil~\ref{fig:ch06:t2} Hausdorff koşulunu, Şekil~\ref{fig:ch06:t3}
regülerlik koşulunu resmeder.

\begin{figure}[ht]
  \centering
  \begin{tikzpicture}[font=\small]
  \draw[blue!70!black, thick, fill=blue!8] (0,0) ellipse [x radius=1.25, y radius=0.85];
  \draw[violet!80!black, thick, fill=violet!10] (3.4,0) ellipse [x radius=1.25, y radius=0.85];
  \fill (0,0) circle (1.8pt) node[below=2pt] {$x$};
  \fill (3.4,0) circle (1.8pt) node[below=2pt] {$y$};
  \node[blue!70!black] at (0,1.15) {$U$};
  \node[violet!80!black] at (3.4,1.15) {$V$};
  \node at (1.7,-1.3) {$U \cap V = \emptyset$};
\end{tikzpicture}

  \caption{T2 (Hausdorff): farklı $x,y$ noktaları, kesişmeyen $U\ni x$
    ve $V\ni y$ açıklarıyla ayrılır.}
  \label{fig:ch06:t2}
\end{figure}

\begin{figure}[ht]
  \centering
  \begin{tikzpicture}[font=\small]
  \fill (0,0) circle (1.8pt) node[below=3pt] {$x$};
  \draw[blue!70!black, thick] (0,0.05) ellipse [x radius=1.0, y radius=0.72];
  \node[blue!70!black] at (0,1.05) {$U$};
  \draw[red!70!black, thick, fill=red!6, pattern=north east lines, pattern color=red!40]
    plot[smooth cycle, tension=0.9] coordinates
    {(3.3,-0.3) (4.4,-0.5) (5.0,0.4) (4.0,0.9) (3.2,0.5)};
  \node[red!70!black] at (4.15,-1.0) {$F$ (kapalı)};
  \draw[violet!80!black, thick] (4.1,0.15) ellipse [x radius=1.55, y radius=1.2];
  \node[violet!80!black] at (4.1,1.6) {$V$};
  \node at (2.05,-1.75) {$x\notin F,\qquad x\in U,\quad F\subseteq V,\quad U\cap V=\emptyset$};
\end{tikzpicture}

  \caption{Regülerlik: $x$ noktası ile onu içermeyen kapalı $F$ kümesi,
    ayrık $U$ ve $V$ açıklarına konabilir. T3 bunun üstüne T1 ister.}
  \label{fig:ch06:t3}
\end{figure}

\begin{karsiornek}
Hiçbir ayrılma aksiyomunu sağlamayan uzay: iki noktalı indirgenmiş
uzay. Açıklar yalnız $\emptyset$ ve $X$ olduğundan iki noktayı ayıran
\emph{hiçbir} açık yoktur --- uzay T0 bile değildir. Örnekler
bölümünde \texttt{two\_point\_indiscrete\_space()} ile test edilir;
zincirin en altının da boş kalabileceğini gösterir.
\end{karsiornek}

%% ============================================================
\section{Teoremler}
%% ============================================================

\begin{theorem}[Ayrılma Zinciri]
T4 $\Rightarrow$ T3.5 $\Rightarrow$ T3 $\Rightarrow$ T2.5 $\Rightarrow$ T2 $\Rightarrow$ T1
$\Rightarrow$ T0. Tersi genel olarak doğru değildir.
\end{theorem}

\begin{proof}[Rehberli Kanıt --- model halka: T2 $\Rightarrow$ T1]
$x\neq y$ verilsin. T2 ile ayrık $U\ni x$, $V\ni y$ açıkları vardır.
$U\cap V=\emptyset$ olduğundan $y\notin U$ ve $x\notin V$: hem
``$x$'i içerip $y$'yi dışlayan'' hem ``$y$'yi içerip $x$'i dışlayan''
açık bulundu --- T1'in iki yönlü koşulu sağlandı. Kalan halkalar
(T1 $\Rightarrow$ T0 dahil) aynı şablonla Alıştırma~T1'de
kanıtlanır.
\end{proof}

\begin{figure}[ht]
  \centering
  \begin{tikzpicture}[font=\small, node distance=9mm,
    aks/.style={draw, rounded corners, fill=blue!6, inner sep=5pt}]
  \node[aks] (t4) {T4};
  \node[aks, right=of t4] (t35) {T3.5};
  \node[aks, right=of t35] (t3) {T3};
  \node[aks, right=of t3] (t25) {T2.5};
  \node[aks, right=of t25] (t2) {T2};
  \node[aks, right=of t2] (t1) {T1};
  \node[aks, right=of t1] (t0) {T0};
  \foreach \a/\b in {t4/t35, t35/t3, t3/t25, t25/t2, t2/t1, t1/t0}
    \draw[-{Stealth[length=2.6mm]}, thick] (\a) -- (\b);
  \draw[red!70!black, -{Stealth[length=2.2mm]}, dashed]
    (t1.south) to[bend right=40] node[below=14pt, font=\scriptsize, align=center]
    {geçmez: kosonlu $\mathbb{N}$\\ (T1, T2 değil)} (t2.south);
  \draw[red!70!black, -{Stealth[length=2.2mm]}, dashed]
    (t0.south) to[bend right=40] node[below=35pt, font=\scriptsize, align=center]
    {geçmez: Sierpiński\\ (T0, T1 değil)} (t1.south);
\end{tikzpicture}

  \caption{Ayrılma zinciri: düz oklar implikasyonları gösterir.
    Kesikli kırmızı oklar terslerin \emph{geçmediği} iki kanonik
    karşı-örneği işaretler; her ikisi de bu kılavuzun örnekleridir
    (Örnek~1 ve Örnek~2).}
  \label{fig:ch06:zincir}
\end{figure}

\begin{theorem}[Urysohn Fonksiyon Teoremi]
$X$ normal ise ve $C, D$ disjoint kapalı kümeler ise, $f: X \to [0,1]$ sürekli vardır:
$f|_C \equiv 0$, $f|_D \equiv 1$.
\end{theorem}

\begin{proof}[İspat eskizi]
Normallikle $C \subseteq U_{1/2}$, $\overline{U_{1/2}} \cap D = \emptyset$
olan açık $U_{1/2}$ seç. Aynı adımı tekrarlayıp her ikili kesir
$q = k/2^n \in [0,1]$ için, $p < q \Rightarrow \overline{U_p} \subseteq U_q$
koşulunu sağlayan iç-içe açıklar ailesi $\{U_q\}$ kur (normallik her ekleme
adımını mümkün kılar). Sonra $f(x) = \inf\{q : x \in U_q\}$ (hiç içermiyorsa
$1$) tanımla; iç-içe geçme süreklilik verir, $C$ üzerinde $0$, $D$ üzerinde
$1$ alınır.
\end{proof}

\begin{figure}[ht]
  \centering
  \begin{tikzpicture}[font=\small]
  \draw[rounded corners=8pt, thick] (-0.4,-1.15) rectangle (6.4,1.45);
  \node at (5.95,1.2) {$X$};
  \draw[blue!70!black, thick, fill=blue!10] (0.8,0.15) ellipse [x radius=0.85, y radius=0.6];
  \node[blue!70!black] at (0.8,0.15) {$C$};
  \draw[red!70!black, thick, fill=red!10] (5.2,0.15) ellipse [x radius=0.85, y radius=0.6];
  \node[red!70!black] at (5.2,0.15) {$D$};
  \draw[-{Stealth}] (7.7,-0.9) -- (7.7,1.3);
  \draw (7.6,-0.6) -- (7.8,-0.6) node[right] {$0$};
  \draw (7.6,1.0) -- (7.8,1.0) node[right] {$1$};
  \draw[blue!70!black, -{Stealth}, dashed]
    (1.7,0.0) to[bend right=14] node[below, font=\scriptsize, pos=0.45] {$f$} (7.58,-0.6);
  \draw[red!70!black, -{Stealth}, dashed] (6.1,0.3) to[bend left=10] (7.58,1.0);
  \node[font=\scriptsize, align=center] at (3.0,-1.65)
    {$f$ sürekli; $f|_C\equiv 0$, $f|_D\equiv 1$; ara bölgede $0$ ile $1$ arası değerler};
\end{tikzpicture}

  \caption{Urysohn fonksiyonu: normal uzayda ayrık kapalı $C$ ve $D$
    için $f|_C\equiv 0$, $f|_D\equiv 1$ olan sürekli $f:X\to[0,1]$
    vardır; uzay, iki kapalı küme arasında ``sürekli geçiş'' yapacak
    kadar zengindir.}
  \label{fig:ch06:urysohn}
\end{figure}

\begin{nedenonemli}
Urysohn teoremi, küme dilindeki bir ayırma özelliğini (normallik)
\emph{sürekli fonksiyon üretimine} çevirir. Tietze Genişleme bunun
doğrudan sonucudur ve metrikleştirme teoremlerinin (Bölüm~9 sonrası
ufuk) temel aracıdır: fonksiyon üretebilen uzay, metrik benzeri
yapılar kurabilen uzaydır.
\end{nedenonemli}

\begin{theorem}[Tietze Genişleme]
$X$ normal ise her kapalı $A \subseteq X$ üzerinde $f: A \to [a,b]$ süreklisi
tüm $X$'e sürekli genişletilebilir.
\end{theorem}

\begin{theorem}[Tychonoff Karakterizasyonu]
$X$, T3.5'tir $\iff$ $X$ bir küp $[0,1]^I$'ya homeomorf gömülebilir.
\end{theorem}

\begin{proof}[İspat eskizi]
($\Leftarrow$) $[0,1]^I$ bir kompakt Hausdorff çarpımdır, dolayısıyla
T3.5'tir; T3.5 kalıtsaldır, alt-uzaya geçer. ($\Rightarrow$) $X$ tam regüler
ise $I = C(X,[0,1])$ (tüm sürekli $[0,1]$-değerli fonksiyonlar) indeks kümesini
al; değerlendirme gömmesi $e(x) = (f(x))_{f\in I}$ tanımla. Tam regülerlik,
$e$'nin birebir ve homeomorfik bir gömme olmasını sağlar --- her nokta--kapalı
çift bir $f$ ile ayrıldığından.
\end{proof}

\begin{theorem}[Sonlu T1 $\iff$ Ayrık]
Sonlu bir uzayda T1 $\iff$ ayrık topoloji.
\end{theorem}
\begin{proof}[Rehberli Kanıt]
\textbf{Strateji:} T1'den tekillerin kapalılığına, oradan sonlu
birleşimle her alt kümenin kapalılığına yürünür.
\begin{enumerate}
  \item ($\Rightarrow$) T1 gereği her $y\neq x$ için $y\in U_y$,
        $x\notin U_y$ olan açık $U_y$ vardır;
        $X\setminus\{x\}=\bigcup_{y\neq x}U_y$ açıktır, yani $\{x\}$
        kapalıdır.
  \item Herhangi $A\subseteq X$, \emph{sonlu} sayıda tekilin birleşimi
        olarak kapalıdır (kapalıların sonlu birleşimi kapalıdır).
  \item Her $A$ kapalı ise her $X\setminus A$ açıktır; $A$ keyfî
        olduğundan her alt küme açıktır: topoloji ayrıktır.
  \item ($\Leftarrow$) Ayrık topolojide her $\{x\}$ açıktır; $x\neq y$
        çifti $\{x\}$ ve $\{y\}$ ile iki yönlü ayrılır: T1 (hatta T2).
\end{enumerate}
Sonsuzlukta 2.~adım çöker: sonsuz birleşim kapalılığı korumaz ---
kosonlu $\mathbb{N}$ tam bu nedenle T1 olup ayrık değildir.
\end{proof}

\noindent Bu teorem, Bölüm~\ref{ch:topological_spaces} K1
alıştırmasındaki gözlemin --- üç noktalı kosonlu uzayın ayrık çıkması
--- genel açıklamasıdır.

\begin{theorem}[Kompakt $+$ Hausdorff $\Rightarrow$ T4]
Her kompakt Hausdorff uzay normaldir (T4).
\end{theorem}
\begin{proof}[İspat eskizi]
Önce ``kompakt Hausdorff $\Rightarrow$ regüler''i kur: $x \notin C$ (kapalı,
dolayısıyla kompakt) ise, her $c \in C$ için Hausdorff ayrık $U_c \ni x$,
$V_c \ni c$ verir; $\{V_c\}$ ailesi $C$'yi örter, kompaktlıkla sonlu alt-örtü
$V_{c_1},\dots,V_{c_n}$ al. $U = \bigcap_i U_{c_i}$ ve $V = \bigcup_i V_{c_i}$
aradığın ayrık açıkları verir. Sonra aynı argümanı bir nokta yerine ikinci bir
kapalı (kompakt) $D$ kümesine uygula: her $d \in D$ için regülerlik ayrık
$U_d \supseteq C$, $W_d \ni d$ verir; $D$ kompakt olduğundan sonlu alt-örtüyle
$C$ ve $D$ ayrık açıklara konur --- normallik tam budur.
\end{proof}

\noindent Bu, Kompaktlık bölümündeki ``kompakt $+$ Hausdorff
$\Rightarrow$ T4'' teoreminin ayrılma-aksiyomu tarafıdır; kompaktlık, sonsuz
Hausdorff ayrımlarını \emph{sonlu} sayıya indirgeyerek normalliği mümkün kılar.

%% ============================================================
\section{Algoritmalar}
%% ============================================================

\begin{algorithm}
\caption{Sonlu T0 Karar Prosedürü}
\begin{algorithmic}[1]
\For{each pair $(x, y)$ with $x \neq y$}
    \If{$\nexists U \in \tau$: $(x \in U \wedge y \notin U)$ or $(y \in U \wedge x \notin U)$}
        \Return \textbf{false}
    \EndIf
\EndFor
\Return \textbf{true}
\end{algorithmic}
\end{algorithm}

\noindent\textbf{Karmaşıklık T0:} $O(|X|^2 \cdot |\tau|)$. \quad
\textbf{Karmaşıklık T2:} $O(|X|^2 \cdot |\tau|^2)$.

\subsection{İz Sürme: T0 Prosedürü Sierpiński Üzerinde}

$X=\{0,1\}$, $\tau=\{\emptyset,\{1\},X\}$:

\begin{table}[ht]
\centering
\begin{tabular}{llccl}
\toprule
\textbf{Çift} $(x,y)$ & \textbf{Denenen} $U$ &
$x\in U \wedge y\notin U$ & $y\in U \wedge x\notin U$ & \textbf{Karar} \\
\midrule
$(0,1)$ & $\emptyset$ & hayır & hayır & devam \\
$(0,1)$ & $\{1\}$ & hayır & \textbf{evet} & çift ayrıldı \\
--- & --- & --- & --- & tüm çiftler bitti $\to$ \textbf{true} \\
\bottomrule
\end{tabular}
\caption{Tek çift tek açıkla ayrılır; genel sınır
  $O(|X|^2\cdot|\tau|)$'dur.}
\end{table}

%% ============================================================
\section{pytop API}
%% ============================================================

\begin{lstlisting}[language=Python, caption={Ayrılma aksiyom importları}]
from pytop import (
    is_t0, is_t1, is_t2, is_t2_5, is_t3, is_t4,
    is_hausdorff, is_regular, is_normal, is_perfectly_normal,
    separation_chain, analyze_separation,
)
\end{lstlisting}

\begin{table}[h]
\centering
\begin{tabular}{lll}
\toprule
\textbf{Fonksiyon} & \textbf{İmza} & \textbf{Açıklama} \\
\midrule
\texttt{is\_t0} \ldots \texttt{is\_t4} & \texttt{(space)} & Aksiyom kontrolü \\
\texttt{separation\_chain} & \texttt{(space)} & Tüm zincir: \texttt{dict[str, Result]} \\
\texttt{analyze\_separation} & \texttt{(space, property='hausdorff')} & Tek aksiyom sorgu \\
\bottomrule
\end{tabular}
\caption{Ayrılma fonksiyonları}
\end{table}

\begin{nedenonemli}
\texttt{is\_*} yüklemleri ham \texttt{bool} değil, \texttt{.status}
alanı \texttt{true} / \texttt{false} / \texttt{unknown} olabilen bir
\texttt{Result} döndürür. Üçüncü değer dürüstlüktür: örneğin T3.5
(Tychonoff) sürekli fonksiyonlarla tanımlanır ve sonlu açık-küme
taramasıyla karar verilemez; \texttt{separation\_chain} bu durumda
\texttt{tychonoff: unknown} raporlar. Sembolik uzaylarda da
bilinmeyen özellikler \texttt{unknown} kalır --- kütüphane bilmediğini
\emph{bilmediğini söyleyerek} belirtir.
\end{nedenonemli}

%% ============================================================
\section{Örnekler}
%% ============================================================

\subsection{Örnek 1: Sierpiński — T0 $\checkmark$, T1 $\times$}

\begin{lstlisting}[language=Python]
from pytop import sierpinski_space, is_t0, is_t1, is_t2
s = sierpinski_space()
print("T0:", is_t0(s).status)
print("T1:", is_t1(s).status)
print("T2:", is_t2(s).status)
\end{lstlisting}
\begin{lstlisting}[caption={Çıktı}]
T0: true
T1: false
T2: false
\end{lstlisting}

\paragraph{Ne oldu?}
\texttt{T0: true} --- $(0,1)$ çifti için $\{1\}$ açığı $1$'i içerir,
$0$'ı dışlar; tek çift bu olduğundan T0 tamam. \texttt{T1: false} ---
ters yön yok: $0$'ı içerip $1$'i dışlayan açık küme yoktur
($\{0\}\notin\tau$; bkz.\ Bölüm~\ref{ch:topological_spaces},
Örnek~1). T1 düşünce T2 de düşer. Sierpiński, zincirin ``T0'da
takılan'' kanonik örneğidir (Şekil~\ref{fig:ch06:zincir}).

\subsection{Örnek 2: Kosonlu $\mathbb{N}$ — T1 $\checkmark$, T2 $\times$}

\begin{lstlisting}[language=Python]
from pytop import naturals_cofinite, is_t1, is_t2
nc = naturals_cofinite()
print("T1:", is_t1(nc).status)
print("T2:", is_t2(nc).status)
\end{lstlisting}
\begin{lstlisting}[caption={Çıktı}]
T1: true
T2: false
\end{lstlisting}

\paragraph{Ne oldu?}
\texttt{T1: true} --- her $n$ için $\mathbb{N}\setminus\{n\}$
kosonludur, dolayısıyla açıktır; bu, her noktayı diğerinden iki yönlü
ayırır. \texttt{T2: false} --- boş olmayan iki kosonlu açık $U,V$
daima kesişir: $\mathbb{N}\setminus(U\cap V)=
(\mathbb{N}\setminus U)\cup(\mathbb{N}\setminus V)$ iki sonlu kümenin
birleşimi olarak sonludur; $\mathbb{N}$ sonsuz olduğundan
$U\cap V\neq\emptyset$. Bu örnek,
Bölüm~\ref{ch:topological_spaces}~K1'in ``neden sonsuz taşıyıcı
gerekir'' sorusunun cevabıdır.

\subsection{Örnek 3: Ayrılma Zinciri — Ayrık}

\begin{lstlisting}[language=Python]
from pytop import discrete_topology, separation_chain
d = discrete_topology(1, 2, 3)
for prop, result in separation_chain(d).items():
    print(f"  {prop:20s}: {result.status}")
\end{lstlisting}
\begin{lstlisting}[style=output, caption={Çıktı}]
  t0                  : true
  t1                  : true
  hausdorff           : true
  urysohn             : true
  t3                  : true
  tychonoff           : true
  t4                  : true
  completely_normal   : true
  perfectly_normal    : true
\end{lstlisting}

\paragraph{Ne oldu?}
Ayrık uzayda her tekil açık olduğundan her ayırma görevi tekil
komşuluklarla çözülür; dokuz yüklemin tümü \texttt{true} döner.
\texttt{tychonoff} bile \texttt{true}'dur: ayrık uzayda her fonksiyon
süreklidir, ayırıcı fonksiyon doğrudan yazılır.
\texttt{separation\_chain}'in anahtar sırası zincirin mantıksal
sırasıdır --- bir uzayın ``nerede takıldığını'' yukarıdan aşağı
okuyabilirsiniz (krş.\ Alıştırma~K2).

\subsection{Örnek 4: Regüler Ama T3 Değil --- Konvansiyon Testi}

\begin{lstlisting}[language=Python]
from pytop import two_point_indiscrete_space, is_regular, is_normal, is_t3, is_t4

tp = two_point_indiscrete_space()
print("regular:", is_regular(tp).status, "| normal:", is_normal(tp).status)
print("t3     :", is_t3(tp).status, "| t4    :", is_t4(tp).status)
\end{lstlisting}

\begin{lstlisting}[style=output, caption={Çıktı}]
regular: true | normal: true
t3     : false | t4    : false
\end{lstlisting}

\paragraph{Ne oldu?}
İndirgenmiş iki noktalı uzayda kapalı kümeler yalnız $\emptyset$ ve
$X$'tir; ``nokta--kapalı'' ve ``kapalı--kapalı'' ayırma koşullarının
öncülü hiç gerçekleşmez, koşullar boş yere sağlanır:
\texttt{regular} ve \texttt{normal} \texttt{true}. Ama uzay T1
olmadığından (tekiller kapalı değil) \texttt{t3} ve \texttt{t4}
\texttt{false} kalır --- ``Dikkat'' kutusundaki konvansiyonun
davranışsal kanıtı.

\subsection{Örnek 5: Dışlanan-Nokta Topolojisi --- İkinci Bir ``T0'da Takılan''}

Sierpiński, T0-olup-T1-olmayan tek örnek değildir. $X=\{0,1,2\}$ üzerinde
\textbf{dışlanan-nokta topolojisi} ($0$ dışlanır): bir küme ya $0$'ı içermez ya
da $X$'in kendisidir.

\begin{lstlisting}[language=Python]
from pytop import excluded_point_topology, is_t0, is_t1, separation_chain

ep = excluded_point_topology(3, 0)   # X = {0,1,2}, dislanan nokta 0
print("T0:", is_t0(ep).status, "| T1:", is_t1(ep).status)
for prop, result in separation_chain(ep).items():
    print(f"  {prop:20s}: {result.status}")
\end{lstlisting}
\begin{lstlisting}[style=output, caption={Çıktı}]
T0: true | T1: false
  t0                  : true
  t1                  : false
  hausdorff           : false
  urysohn             : false
  t3                  : false
  tychonoff           : false
  t4                  : false
  completely_normal   : false
  perfectly_normal    : false
\end{lstlisting}

\paragraph{Ne oldu?}
\texttt{T0: true} --- $1$ ve $2$, $0$-içermeyen $\{1\}$/$\{2\}$ açıklarıyla iki
yönlü ayrılır; $0$ ile herhangi bir nokta arasında ise $0$'ı dışlayan açık tek
yönlü ayrım verir. \texttt{T1: false} --- $0$'ı içeren tek açık $X$'tir,
dolayısıyla $\{0\}$ kapalı değildir. Sierpiński'den farklı bir mekanizmayla
zincir yine T0'da takılır.

\subsection{Örnek 6: Metrik Uzaylar --- Zincirin Tepesine Kadar}

Her metrik uzay normaldir (T4) ve zincirin tüm üst basamaklarını sağlar; bunu
$\mathbb{R}$ ve $[0,1]$ üzerinde gözlemleyelim.

\begin{lstlisting}[language=Python]
from pytop import real_line_metric, closed_unit_interval_metric
from pytop import is_t2, is_urysohn, is_t3, is_t4

rl = real_line_metric()
ui = closed_unit_interval_metric()
print("R         t2/urysohn/t3/t4:", is_t2(rl).status, is_urysohn(rl).status, is_t3(rl).status, is_t4(rl).status)
print("[0,1]     t2/urysohn/t3/t4:", is_t2(ui).status, is_urysohn(ui).status, is_t3(ui).status, is_t4(ui).status)
\end{lstlisting}
\begin{lstlisting}[style=output, caption={Çıktı}]
R         t2/urysohn/t3/t4: true true true true
[0,1]     t2/urysohn/t3/t4: true true true true
\end{lstlisting}

\paragraph{Ne oldu?}
Metrik mesafe kendisi bir ayırıcıdır: ayrık kapalı $C, D$ için
$f(x) = d(x,C)/(d(x,C)+d(x,D))$ sürekli Urysohn fonksiyonunu doğrudan verir ---
yani metrik uzaylar yalnız T2 değil, T4'e kadar her aksiyomu sağlar
(Teorem~2.2'nin metrik durumda elle yazılabilen tanığı). Ayrılma zincirinin
sonlu/sembolik örneklerle metrik örnekler arasındaki keskin farkı budur.

%% ============================================================
\section{Alıştırmalar}
%% ============================================================

\subsection*{Kodlama}
\begin{enumerate}[label=\textbf{K\arabic*.}]
  \item \label{ex:ch06:K1} \texttt{make\_topology(\{1,2,3\},\{1\},\{2\},\{1,2\})}
        için \texttt{separation\_chain} çalıştırın; her sonucu yorumlayın.
        \ipucu{Önce açıkları listeleyin:
        $\emptyset,\{1\},\{2\},\{1,2\},X$. Her çifti ayıran açık
        arayın; $3$'ü içerip $1$'i dışlayan açık var mı?}
        (\hyperref[sol:ch06:K1]{çözüm})

  \item \label{ex:ch06:K2} \texttt{finite\_chain\_space(3)} üzerinde en
        yüksek sağlanan aksiyomu bulun.
        \ipucu{Çıktıda \texttt{true} olan en güçlü anahtarı arayın;
        \texttt{unknown} değerlerini ``Neden önemli?'' kutusuna göre
        yorumlayın.}
        (\hyperref[sol:ch06:K2]{çözüm})

  \item \label{ex:ch06:K3} \texttt{two\_point\_indiscrete\_space()}
        üzerinde T0, T1, T2 test edin.
        \ipucu{Tek boş olmayan açık $X$ iken herhangi bir çift nasıl
        ayrılabilir?}
        (\hyperref[sol:ch06:K3]{çözüm})

  \item \label{ex:ch06:K4} \texttt{excluded\_point\_topology(4, 0)} (yani
        $X=\{0,1,2,3\}$, dışlanan nokta $0$) üzerinde
        \texttt{separation\_chain} çalıştırın; sağlanan en yüksek aksiyomu
        belirleyin ve neden T1'in düştüğünü açıklayın.
        \ipucu{$0$'ı içeren tek açık $X$'tir; $\{0\}$ kapalı mı?}
        (\hyperref[sol:ch06:K4]{çözüm})
\end{enumerate}

\subsection*{Teori}
\begin{enumerate}[label=\textbf{T\arabic*.}]
  \item \label{ex:ch06:T1} T2 $\Rightarrow$ T1 $\Rightarrow$ T0
        implikasyonlarını kanıtlayın.
        \ipucu{T2 $\Rightarrow$ T1: ayrık $U\ni x$, $V\ni y$
        verildiğinde $U$, $y$'yi; $V$, $x$'i dışlar --- iki yönlü ayrım
        hazır. T1 $\Rightarrow$ T0: iki yönlü ayrım tek yönlüyü içerir.}
        (\hyperref[sol:ch06:T1]{çözüm})

  \item \label{ex:ch06:T2} Sonlu uzayda T1 $\iff$ ayrık topoloji
        olduğunu gösterin.
        \ipucu{T1 $\Rightarrow$ tekiller kapalı $\Rightarrow$ (sonlu
        birleşim) her alt küme kapalı $\Rightarrow$ her alt küme açık.}
        (\hyperref[sol:ch06:T2]{çözüm})

  \item \label{ex:ch06:T3} Her metrik uzayın normal (T4) olduğunu
        kanıtlayın.
        \ipucu{Ayrık kapalı $C, D$ için
        $f(x) = d(x,C)/(d(x,C)+d(x,D))$ sürekli Urysohn fonksiyonunu yazın;
        $f^{-1}([0,\tfrac12))$ ve $f^{-1}((\tfrac12,1])$ aradığınız ayrık
        açıklardır.}
        (\hyperref[sol:ch06:T3]{çözüm})
\end{enumerate}

\chapter{Kompaktlık}

Kompaktlık, sonsuz yapılarda ``sınırlılık'' fikrinin topolojik soyutlamasıdır.
Açık örtü tanımı temel alınarak incelenmekte; varyantlar ve lokal kompaktlık da ele alınmaktadır.

%-------------------------------------------------------------------
\section{Konu}

\subsection{Sezgisel Giriş — ``Gizli Sonluluk''}

Kompaktlığı tek cümlede özetlemek gerekirse: \emph{sonsuz bir uzayın, sonlu bir
uzay gibi davranmasıdır.} Sonlu bir kümede her şey kolaydır; kompaktlık bu
``sonluluk konforunu'' sonsuz uzaylara taşıyan köprüdür: uzayı ne kadar çok açık
kümeyle örtersek örtelim, \textbf{her zaman sonlu sayıda tanesi yeter.}

\begin{figure}[ht]
  \centering
  \begin{tikzpicture}[font=\small,
    kapsul/.style={rounded corners=6pt, thick, fill opacity=0.18, draw opacity=0.9},
    etiket/.style={font=\scriptsize}]
  % --- Ust: sonsuz acik ortu ---
  \node[etiket, align=left] at (-1.7,1.05) {Açık örtü\\(sonsuz olabilir)};
  \draw[very thick] (0,0.6) -- (8,0.6);
  \node[left] at (0,0.6) {$X$};
  \foreach \x/\c in {0.2/blue, 1.4/red, 2.6/teal, 3.8/orange, 5.0/violet, 6.2/blue, 7.0/red} {
    \fill[\c!60, fill opacity=0.18, rounded corners=5pt] (\x,0.35) rectangle ++(1.6,0.5);
    \draw[\c!70!black, thick, rounded corners=5pt] (\x,0.35) rectangle ++(1.6,0.5);
  }

  % --- ok ---
  \draw[-{Stealth[length=3mm]}, very thick] (4,-0.05) -- (4,-0.7)
    node[midway, right=2pt, etiket] {kompakt $\Rightarrow$ sonlu alt-örtü};

  % --- Alt: sonlu alt-ortu ---
  \node[etiket, align=left] at (-1.7,-1.45) {Sonlu alt-örtü\\(3 küme yeter)};
  \draw[very thick] (0,-1.4) -- (8,-1.4);
  \node[left] at (0,-1.4) {$X$};
  \foreach \x/\c in {0.2/blue, 2.8/teal, 5.6/violet} {
    \fill[\c!60, fill opacity=0.22, rounded corners=5pt] (\x,-1.65) rectangle ++(2.6,0.5);
    \draw[\c!70!black, very thick, rounded corners=5pt] (\x,-1.65) rectangle ++(2.6,0.5);
  }
\end{tikzpicture}

  \caption{Açık örtü $\to$ sonlu alt-örtü: uzayı sonsuz çoklukta açık kümeyle
    örtsek bile, kompaktsa sonlu sayıda tanesi onu kapatmaya yeter.}
\end{figure}

\begin{sezgi}
``Açık örtü'' uzayı kaplayan açık küme ailesidir; ``sonlu alt-örtü'' bu aileden
seçilmiş, hâlâ kaplamayı başaran sonlu bir parçadır. Kompakt olmak, bunu
\emph{her} örtü için yapabilmektir. Bu özellik, analizdeki maksimum değer
teoremi, düzgün süreklilik ve Heine-Borel'in arkasındaki sessiz kahramandır.
\end{sezgi}

\subsection{Kompaktlık Tanımı}

\begin{definition}[Kompakt Uzay]
$(X, \tau)$ topolojik uzayı \textbf{kompakt} ise:
Her açık örtü $\{U_\alpha\}$ için sonlu alt-örtü vardır:
\[
X \subseteq \bigcup_{\alpha} U_\alpha
\implies
\exists \text{ sonlu } I:\; X \subseteq \bigcup_{\alpha \in I} U_\alpha.
\]
\end{definition}

\subsection{Neden ``kapalı ve sınırlı'' yetmez? — Karşı-örnekler}

Reel doğruda kompaktlık tam olarak ``kapalı + sınırlı''ya denktir (Heine-Borel).
Bu denkliğin neden \emph{hem} kapalılığa \emph{hem} sınırlılığa ihtiyaç
duyduğunu en iyi açık aralık gösterir.

\begin{figure}[ht]
  \centering
  \begin{tikzpicture}[font=\small, etiket/.style={font=\scriptsize}]
  % open interval (0,1) on a number line scaled to width 9
  \def\W{9}
  \draw[very thick] (0,0) -- (\W,0);
  \draw[fill=white, thick] (0,0) circle (2.2pt);   % open endpoint 0
  \draw[fill=white, thick] (\W,0) circle (2.2pt);  % open endpoint 1
  \node[below=3pt] at (0,0) {$0$};
  \node[below=3pt] at (\W,0) {$1$};
  \node[etiket, align=center] at (\W/2,3.15)
    {Açık örtü $U_n=\left(\tfrac{1}{n},\,1\right)$, \; $n=2,3,4,\dots$};

  % nested intervals U_n = (1/n, 1): left endpoint at W/n
  \foreach \n/\y/\c in {2/0.55/blue, 3/0.95/teal, 4/1.35/orange, 5/1.75/violet, 6/2.15/red} {
    \pgfmathsetmacro\lx{\W/\n}
    \draw[\c!75!black, very thick] (\lx,\y) -- (\W,\y);
    \draw[\c!75!black, fill=white, thick] (\lx,\y) circle (1.8pt);
    \node[\c!60!black, etiket, left=1pt] at (\lx,\y) {$U_{\n}$};
  }
  % arrow showing escape toward 0
  \draw[-{Stealth[length=2.6mm]}, red!70!black, thick]
    (2.0,2.45) -- (0.35,2.45);
  \node[red!70!black, etiket, align=center] at (4.6,2.5)
    {sol uç $\tfrac{1}{n}\to 0$ ama $0$'a hiç ulaşmaz};

  \node[etiket, align=center, fill=red!6, rounded corners, inner sep=4pt]
    at (\W/2,-1.0)
    {Hiçbir \emph{sonlu} alt-örtü $0$'a yaklaşan kısmı kapatamaz
     $\;\Rightarrow\;(0,1)$ kompakt \textbf{değil}.\\[2pt]
     Kapalı $[0,1]$'de $0$ dahil olduğundan sonlu alt-örtü bulunur $\Rightarrow$ kompakt.};
\end{tikzpicture}

  \caption{Kaçan örtü: $(0,1)$ açık aralığını $U_n=(1/n,1)$ ailesi örter; sol
    uçlar $0$'a yaklaşır ama $0$'a hiç ulaşmaz, bu yüzden hiçbir sonlu alt-örtü
    işe yaramaz.}
\end{figure}

\begin{karsiornek}
\textbf{$(0,1)$ kompakt değil.} $U_n = \left(\tfrac{1}{n}, 1\right)$ ailesi
$(0,1)$'i örter. Sonlu bir alt-aile $\{U_{n_1}, \dots, U_{n_k}\}$ seçersek, en
büyük indis $N=\max n_i$ için birleşim yalnızca $\left(\tfrac{1}{N}, 1\right)$
olur; $0$'a yakın noktalar açıkta kalır. Sonlu alt-örtü yoktur. Kapalı $[0,1]$'de
$0$ dahil olduğundan sonlu alt-örtü her zaman bulunur.

\medskip
\textbf{$\mathbb{R}$ kompakt değil.} $V_n = (-n, n)$ ailesi $\mathbb{R}$'yi örter
ama hiçbir sonlu alt-aile sınırsız $\mathbb{R}$'yi kaplayamaz. Burada sorun
\emph{sınırsızlık}.
\end{karsiornek}

\subsection{Kompaktlık Varyantları}

\begin{table}[h]
\centering
\begin{tabular}{ll}
\toprule
\textbf{Varyant} & \textbf{Tanım} \\
\midrule
Kompakt & Her açık örtünün sonlu alt-örtüsü \\
Sayılabilir kompakt & Her sayılabilir örtünün sonlu alt-örtüsü \\
Ardışımlı kompakt & Her dizi yakınsak alt-dizi içerir \\
Sözde kompakt & Her sürekli $f: X \to \mathbb{R}$ sınırlıdır \\
Lindelöf & Her açık örtünün sayılabilir alt-örtüsü \\
$\sigma$-kompakt & Sayılabilir kompakt kümelerin birleşimi \\
Lokal kompakt & Her noktanın kompakt komşuluğu var \\
\bottomrule
\end{tabular}
\caption{Kompaktlık varyantları}
\end{table}

\begin{figure}[ht]
  \centering
  \begin{tikzpicture}[font=\small, node distance=12mm and 18mm,
    kut/.style={draw, rounded corners, fill=blue!6, inner sep=6pt, align=center},
    vurgu/.style={draw, rounded corners, fill=green!8, inner sep=6pt, align=center, very thick}]
  \node[vurgu] (k) {Kompakt};
  \node[kut, above right=of k] (lind) {Lindelöf};
  \node[kut, below right=of k] (say) {Sayılabilir\\kompakt};
  \node[kut, right=28mm of k] (ard) {Ardışımlı\\kompakt};

  \draw[-{Stealth[length=2.6mm]}, thick] (k) -- (lind);
  \draw[-{Stealth[length=2.6mm]}, thick] (k) -- (say);
  \draw[-{Stealth[length=2.6mm]}, thick] (ard) -- (say);

  % metric equivalence
  \draw[{Stealth[length=2.4mm]}-{Stealth[length=2.4mm]}, dashed, teal!70!black]
    (k.east) to[bend left=10] node[above, font=\scriptsize, align=center] {metrikte\\ $\Leftrightarrow$} (ard.west);

  \node[font=\scriptsize, align=left, anchor=north west] at (-2.2,-2.0)
    {Oklar: $\Rightarrow$ (her zaman geçer).\\
     Kesik çift ok: yalnız \emph{metrik} uzaylarda denklik.};
\end{tikzpicture}

  \caption{Kompaktlık varyantları arasındaki implikasyonlar: Kompakt
    $\Rightarrow$ Lindelöf ve Kompakt $\Rightarrow$ Sayılabilir kompakt her zaman
    geçer; Kompakt $\Leftrightarrow$ Ardışımlı kompakt yalnız metrik uzaylarda.}
\end{figure}

\textbf{Her zaman geçen implikasyonlar:} Kompakt $\Rightarrow$ Sayılabilir
kompakt, Kompakt $\Rightarrow$ Lindelöf, Ardışımlı kompakt $\Rightarrow$
Sayılabilir kompakt. \textbf{Metrik uzaylarda} üç kavram (kompakt, sayılabilir
kompakt, ardışımlı kompakt) çakışır.

%-------------------------------------------------------------------
\section{Teoremler}

\begin{theorem}[Heine-Borel]
$\mathbb{R}^n$'de kompakt $\Leftrightarrow$ kapalı ve sınırlı.
\end{theorem}

\begin{proof}[İspat eskizi]
($\Leftarrow$) $[a,b]$ aralığını ele alalım ve sonlu alt-örtüsü olmayan bir açık
örtü bulunduğunu varsayalım. Aralığı ikiye böleriz; en az bir yarısının da sonlu
alt-örtüsü yoktur. Bunu yineleyerek uzunluğu $0$'a giden iç içe kapalı aralıklar
dizisi elde ederiz; kesişimleri tek bir $p$ noktasıdır. $p$'yi içeren bir örtü
elemanı $U$, yeterince küçük aralıkları tümüyle içerir --- çelişki. $\mathbb{R}^n$
için Tychonoff ile $[a,b]^n$'e geçilir. ($\Rightarrow$) Kompakt küme Hausdorff
uzayda kapalıdır; sınırlılık, $\{(-n,n)^n\}$ örtüsünün sonlu alt-örtüsünden gelir.
\end{proof}

\begin{theorem}
Kompakt $+$ Hausdorff $\Rightarrow$ Normal (T4).
\end{theorem}

\begin{proof}[İspat eskizi]
$A, B$ ayrık kapalı kümeler olsun (kompakt uzayda kapalı $\Rightarrow$ kompakt).
Sabit $a\in A$ için Hausdorff'luk her $b\in B$'yi $a$'dan açık $U_b \ni a$,
$V_b \ni b$ ile ayırır. $\{V_b\}$, kompakt $B$'yi örter; sonlu alt-örtü
$V_{b_1},\dots,V_{b_k}$ alıp $U_a=\bigcap U_{b_i}$, $V_a=\bigcup V_{b_i}$ koyarız.
Burada \emph{sonlu} kesişimin açık kalması kritiktir --- bu yüzden $B$'nin kompakt
olması şarttır. Sonra $\{U_a\}_{a\in A}$ kompakt $A$'yı örter; yine sonlu
alt-örtüden $A$ ile $B$'yi ayıran açıklar kurulur.
\end{proof}

\begin{theorem}
Kompakt uzaydan Hausdorff uzaya sürekli bijeksiyon $\Rightarrow$ homeomorfizma.
\end{theorem}

\begin{theorem}[Tychonoff]
Keyfi kompakt uzayların kartezyen çarpımı kompakttır.
\end{theorem}

\begin{theorem}[Alexandroff Tek-Nokta Kompaktifikasyonu]
Her lokal kompakt Hausdorff uzay $X$ için $X^* = X \cup \{\infty\}$ kompakt
Hausdorff'tur. (Örn.\ $\mathbb{R}^* \cong S^1$, $(\mathbb{R}^2)^* \cong S^2$.)
\end{theorem}

\begin{theorem}[Kompakt görüntü]
$f: X \to Y$ sürekli ve $X$ kompakt ise $f(X)$ kompakttır.
\end{theorem}

\begin{proof}[İspat eskizi]
$\{V_\alpha\}$, $f(X)$'i örten açık aile olsun. Süreklilikten
$\{f^{-1}(V_\alpha)\}$, $X$'i örten açık ailedir; $X$ kompakt olduğundan sonlu
alt-örtü $f^{-1}(V_{\alpha_1}),\dots,f^{-1}(V_{\alpha_k})$ vardır. O hâlde
$V_{\alpha_1},\dots,V_{\alpha_k}$, $f(X)$'i örter.
\end{proof}

%-------------------------------------------------------------------
\section{Algoritmalar}

\subsection{Sonlu Uzayda Kompaktlık}

Sonlu uzay her zaman kompakttır: her örtü zaten sonlu.
Algoritma: $O(1)$.

\subsection{analyze\_compactness\_variants — Tag Tabanlı Çıkarım}

\begin{algorithm}
\caption{AnalyzeVariants$(X, \tau)$}
\begin{algorithmic}[1]
\If{$X$ sonlu}
    \State tüm varyantlar $\leftarrow$ True; \Return VariantProfile(...)
\EndIf
\If{\texttt{'compact'} $\in$ tags}
    \State countably\_compact, sequentially\_compact, lindelof $\leftarrow$ True
\EndIf
\If{\texttt{'metric'} $\wedge$ \texttt{'compact'} $\in$ tags}
    \State sequentially\_compact $\leftarrow$ True
\EndIf
\end{algorithmic}
\end{algorithm}

%-------------------------------------------------------------------
\section{pytop API}

\begin{lstlisting}[language=Python]
from pytop import (
    is_compact,
    is_lindelof,
    is_locally_compact,
    naturals_cofinite,
    one_point_compactification_of_reals,
    one_point_compactification_of_N,
)
from pytop.compactness_variants import (
    is_countably_compact,
    is_sequentially_compact,
    analyze_compactness_variants,
)
\end{lstlisting}

\texttt{analyze\_compactness\_variants(space)} $\to$ \texttt{Result}; \texttt{.value} bir \texttt{dict} içerir.

%-------------------------------------------------------------------
\section{Örnekler}

\subsection*{Örnek 5.1 — Sonlu Uzaylar: Otomatik Kompakt}

\begin{lstlisting}[language=Python]
from pytop import sierpinski_space, discrete_topology, finite_chain_space, is_compact

s = sierpinski_space()
print("Sierpinski compact?", is_compact(s).status)
d = discrete_topology(1, 2, 3)
print("Discrete(3) compact?", is_compact(d).status)
c = finite_chain_space(3)
print("Chain(3) compact?", is_compact(c).status)
\end{lstlisting}

\begin{verbatim}
Sierpinski compact?      true
Chain(3) compact?        true
Discrete(3) compact?     true
\end{verbatim}

\subsection*{Örnek 5.2 — Gerçek Doğru $\mathbb{R}$: Kompakt Değil}

\begin{lstlisting}[language=Python]
rl = real_line_metric()
print("compact?", is_compact(rl).status)
print("lindelof?", is_lindelof(rl).status)
print("locally_compact?", is_locally_compact(rl).status)
\end{lstlisting}

\begin{verbatim}
compact?          false
lindelof?         true
locally_compact?  conditional
\end{verbatim}

\subsection*{Örnek 5.3 — analyze\_compactness\_variants: Sierpiński}

\begin{lstlisting}[language=Python]
r = analyze_compactness_variants(s)
for key, val in r.value.items():
    if hasattr(val, 'status'):
        print(f"  {key:<25}: {val.status}")
\end{lstlisting}

\begin{verbatim}
  representation           : finite
  countably_compact        : true
  sequentially_compact     : true
  pseudocompact            : true
  lindelof                 : true
\end{verbatim}

\subsection*{Örnek 5.4 — Sonsuz Ama Kompakt: Kosonlu Topoloji}

Kompaktlık sınırlılığa değil, \emph{örtü davranışına} bağlıdır. Sonsuz bir küme
üzerindeki kosonlu topoloji buna örnektir: bir açık örtüde tek bir eleman
seçtiğimizde dışarıda yalnız sonlu çoklukta nokta kalır; onları da sonlu sayıda
elemanla kapatırız.

\begin{lstlisting}[language=Python]
from pytop import naturals_cofinite, is_compact, is_lindelof
from pytop.compactness_variants import is_countably_compact

nc = naturals_cofinite()
print("compact?", is_compact(nc).status)
print("lindelof?", is_lindelof(nc).status)
print("countably?", is_countably_compact(nc).status)
\end{lstlisting}

\begin{verbatim}
compact?   true
lindelof?  true
countably? true
\end{verbatim}

Sonsuz olmasına karşın kompakt --- ``kompakt $=$ sonlu'' sezgisinin neden yalnız
\emph{metrik} uzaylarda işlediğini gösterir.

\subsection*{Örnek 5.5 — Tek-Nokta Kompaktifikasyon (Alexandroff)}

Lokal kompakt ama kompakt olmayan bir uzaya tek bir $\infty$ noktası ekleyerek onu
kompakt hâle getirebiliriz. $\mathbb{R}$ için sonuç çembere ($S^1$) denktir.

\begin{lstlisting}[language=Python]
from pytop import (
    one_point_compactification_of_reals,
    one_point_compactification_of_N,
    is_compact,
)

r_star = one_point_compactification_of_reals()   # R u {inf} ~= S^1
n_star = one_point_compactification_of_N()
print("R* compact?", is_compact(r_star).status)
print("N* compact?", is_compact(n_star).status)
\end{lstlisting}

\begin{verbatim}
R* compact? true
N* compact? true
\end{verbatim}

$\mathbb{R}$ kompakt değildi (Örnek 5.2); tek bir nokta eklemek onu kompakt yaptı.

%-------------------------------------------------------------------
\section{Alıştırmalar}

\subsection*{Kodlama}

\begin{enumerate}[label=K\arabic*.]
\item \texttt{finite\_chain\_space(5)} ve \texttt{discrete\_topology(1,2,3,4,5)} için
      \texttt{is\_compact} karşılaştırın.
\item \texttt{naturals\_cofinite()} için \texttt{is\_compact}, \texttt{is\_lindelof},
      \texttt{is\_countably\_compact} hesaplayın.
\item \texttt{closed\_unit\_interval\_metric()} ile \texttt{real\_line\_metric()} için
      \texttt{analyze\_compactness\_variants} çalıştırıp karşılaştırın.
\item \texttt{real\_line\_metric()} kompakt değildi.
      \texttt{one\_point\_compactification\_of\_reals()} çıktısının
      \texttt{is\_compact} durumunu kontrol edip tek noktanın kompaktlığı nasıl
      kurtardığını açıklayın.
\end{enumerate}

\subsection*{Teori}

\begin{enumerate}[label=T\arabic*.]
\item Kompakt + sürekli $\Rightarrow$ görüntü kompakt olduğunu ispatlayın.
\item Heine-Borel teoreminin neden geçerli olduğunu: $[0,1]$ neden kompakt,
      $(0,1)$ neden kompakt değil? (İpucu: kaçan örtü $U_n=(1/n,1)$.)
\item Kompakt $+$ Hausdorff $\Rightarrow$ Normal ispatında, $B$'yi örten
      $\{V_b\}$ ailesinden sonlu alt-örtü çekme adımının neden kompaktlığa
      ihtiyaç duyduğunu açıklayın.
\end{enumerate}

\chapter{Bağlantılılık}

Bağlantılılık, bir topolojik uzayın iki ayrı açık parçaya bölünememesi özelliğidir.
Yol-bağlantılılık daha güçlü bir kavramdır ve sürekli yolların varlığına dayanır.

%-------------------------------------------------------------------
\section{Konu}

\subsection{Sezgisel Giriş — ``Tek Parça mı, Kopuk mu?''}

Bağlantılılığı tek cümlede özetlemek gerekirse: \emph{bir uzayın iki ayrı açık
parçaya kesilememesidir.} Bir kâğıt parçasını düşünün; onu hiçbir yeri yırtmadan
iki ayrı açık parçaya bölemiyorsanız o kâğıt bağlantılıdır. Baştan iki ayrı
kâğıdınız varsa uzay bağlantısızdır.

\begin{figure}[ht]
  \centering
  \begin{tikzpicture}[font=\small, etiket/.style={font=\scriptsize}]
  % --- Sol: baglantili uzay (tek parca) ---
  \node[etiket] at (1.6,2.55) {\textbf{Bağlantılı}};
  \draw[blue!70!black, very thick, fill=blue!12, rounded corners=14pt]
    (0.2,0.4) ellipse (1.5 and 1.0);
  \node at (1.6,0.4) {$X$};
  \node[etiket, align=center] at (1.6,-1.05)
    {Tek parça: boş olmayan,\\$X$'ten farklı clopen küme \emph{yok}};

  % --- Orta: ayirici ok ---
  \draw[-{Stealth[length=2.6mm]}, thick] (3.5,0.4) -- (4.6,0.4)
    node[midway, above, etiket] {kopar};

  % --- Sag: baglantisiz uzay (iki ayri acik parca) ---
  \node[etiket] at (7.4,2.55) {\textbf{Bağlantısız}};
  \draw[teal!70!black, very thick, fill=teal!12, rounded corners=10pt]
    (5.6,0.4) ellipse (0.95 and 0.85);
  \node at (6.05,0.4) {$U$};
  \draw[orange!80!black, very thick, fill=orange!14, rounded corners=10pt]
    (8.4,0.4) ellipse (0.95 and 0.85);
  \node at (8.75,0.4) {$V$};
  \node[etiket, align=center] at (7.4,-1.05)
    {$X = U \sqcup V$,\quad $U,V$ açık ve boş değil\\
     ikisi de clopen $\Rightarrow$ bağlantısız};
\end{tikzpicture}

  \caption{Bağlantılı (tek parça) ile bağlantısız ($X = U \sqcup V$, iki ayrı
    clopen parça) arasındaki fark.}
\end{figure}

Bu ``kesilememe'' fikrinin teknik karşılığı \textbf{clopen ayrılma}dır: uzay
bağlantısızdır ancak ve ancak boş olmayan, uzayın tamamından farklı, hem açık hem
kapalı (\emph{clopen}) bir küme varsa.

\begin{sezgi}
``Bağlantılı'' uzayın tek parça olmasıdır; ``yol-bağlantılı'' ise herhangi iki
noktayı sürekli bir yol ile birleştirebilmektir. İkinci kavram birincisinden
\emph{daha güçlü}: yol kurabiliyorsanız zaten tek parçasınızdır, ama tek parça
olmak her zaman yol kurabildiğiniz anlamına gelmez (topolojist sinüs eğrisi).
\end{sezgi}

\subsection{Bağlantılılık Kavramları}

\begin{definition}[Bağlantılı Uzay]
$(X, \tau)$ uzayı \textbf{bağlantılı} ise: $X = U \cup V$, $U \cap V = \emptyset$,
$U,V$ açık $\Rightarrow U = \emptyset$ veya $V = \emptyset$.
\end{definition}

\begin{definition}[Yol-Bağlantılı]
$\forall x,y \in X$: $\exists$ sürekli $f:[0,1] \to X$, $f(0) = x$, $f(1) = y$.
\end{definition}

\begin{table}[h]
\centering
\begin{tabular}{ll}
\toprule
\textbf{Kavram} & \textbf{Özellik} \\
\midrule
Bağlantılı & Clopen bölme yok \\
Yol-bağlantılı & Her iki nokta arasında sürekli yol var \\
Ark-bağlantılı & Yol-bağlantılı; yollar enjektif \\
Lokal bağlantılı & Her noktanın bağlantılı komşuluğu \\
Tamamen bağlantısız & Her bağlantılı alt küme tek noktadır \\
\bottomrule
\end{tabular}
\caption{Bağlantılılık hiyerarşisi}
\end{table}

\textbf{Sıralama:} Ark-bağlantılı $\Rightarrow$ Yol-bağlantılı $\Rightarrow$ Bağlantılı.

\subsection{Clopen Ayrılma}

$X$ bağlantısız $\Leftrightarrow$ boş olmayan trivial olmayan bir clopen $A \subsetneq X$ vardır.

\subsection{Bağlantılı Bileşenler}

Bir uzayı, içerme bakımından \textbf{en büyük} bağlantılı alt kümelerine
ayırabiliriz; bunlara \emph{bağlantılı bileşenler} denir. Bileşenler $X$'i örten,
ayrık parçalardır. Uzay bağlantılıdır $\Leftrightarrow$ tek bir bileşeni vardır.

\begin{figure}[ht]
  \centering
  \begin{tikzpicture}[font=\small, etiket/.style={font=\scriptsize},
    nokta/.style={circle, fill, inner sep=1.6pt}]
  % ust cerceve: bir uzayin 3 ayri "ada"si = 3 baglantili bilesen
  \draw[gray!55, thick, rounded corners=8pt] (-0.5,-1.3) rectangle (9.5,1.6);
  \node[gray!60!black, etiket, anchor=north west] at (-0.4,1.5) {$X$};

  % Bilesen 1: ucgen (baglantili)
  \node[blue!70!black, etiket] at (1.2,1.2) {$C_1$};
  \draw[blue!70!black, thick, fill=blue!10, rounded corners=10pt]
    (0.0,-0.9) -- (2.4,-0.9) -- (1.2,0.9) -- cycle;
  \node[nokta, blue!70!black] (a1) at (0.4,-0.6) {};
  \node[nokta, blue!70!black] (a2) at (2.0,-0.6) {};
  \node[nokta, blue!70!black] (a3) at (1.2,0.6) {};
  \draw[blue!70!black] (a1)--(a2)--(a3)--(a1);

  % Bilesen 2: tek nokta (izole)
  \node[teal!70!black, etiket] at (4.5,1.2) {$C_2$};
  \draw[teal!70!black, thick, fill=teal!10] (4.5,0.0) circle (0.7);
  \node[nokta, teal!70!black] at (4.5,0.0) {};

  % Bilesen 3: yol (baglantili)
  \node[orange!80!black, etiket] at (7.6,1.2) {$C_3$};
  \draw[orange!80!black, thick, fill=orange!10, rounded corners=10pt]
    (6.3,-0.8) rectangle (9.0,0.7);
  \node[nokta, orange!80!black] (b1) at (6.7,-0.3) {};
  \node[nokta, orange!80!black] (b2) at (7.65,0.3) {};
  \node[nokta, orange!80!black] (b3) at (8.6,-0.3) {};
  \draw[orange!80!black] (b1)--(b2)--(b3);

  \node[etiket, align=center] at (4.5,-2.0)
    {Bağlantılı \emph{bileşenler}: $X$'i örten, en büyük bağlantılı, ayrık parçalar.\\
     $X = C_1 \sqcup C_2 \sqcup C_3$ \quad — \quad $X$ bağlantılı $\iff$ tek bileşen ($k=1$).};
\end{tikzpicture}

  \caption{Bağlantılı bileşenler: $X = C_1 \sqcup C_2 \sqcup C_3$; her parça en
    büyük bağlantılı alt kümedir. $X$ bağlantılı $\iff$ tek bileşen.}
\end{figure}

\begin{karsiornek}
\textbf{Topolojist sinüs eğrisi — bağlantılı ama yol-bağlantılı değil.}
$S = \{(x, \sin\tfrac{1}{x}) : 0 < x \le 1\} \cup (\{0\}\times[-1,1])$ kümesini
ele alın. Sağdaki eğri $x \to 0^+$ giderken sonsuz hızla salınır ve soldaki dikey
segmente keyfi yakın olur; bu yüzden $S$ \textbf{bağlantılıdır}. Ama dikey
segmentteki bir noktayı eğri üzerindeki bir noktaya bağlayan \emph{sürekli} bir
yol \textbf{yoktur} --- yol $x=0$'a yaklaşırken sinüs salınımı yakınsamayı
engeller. Demek ki $S$ bağlantılı, fakat yol-bağlantılı değildir.
\end{karsiornek}

\begin{figure}[ht]
  \centering
  \begin{tikzpicture}[font=\small, etiket/.style={font=\scriptsize}]
  % topolojist sinus egrisi: y = sin(1/x), x in (0, ~0.9], + dikey segment {0}x[-1,1]
  % dikey limit segmenti
  \draw[blue!70!black, very thick] (0,-1) -- (0,1);
  \node[blue!60!black, etiket, left=2pt] at (0,0) {$\{0\}\times[-1,1]$};

  % sin(1/x) egrisi: gercek x in (0, 0.9], yatay olcek 3.3 -> grafik x ekseni
  % yakin x=0 icin yogun ornekleme (sonsuz salinim gorunur olsun)
  \draw[orange!80!black, very thick, domain=0.028:0.9, samples=900, smooth, variable=\x]
    plot ({\x*3.3},{sin((1/\x) r)});
  \node[orange!75!black, etiket] at (2.05,1.45) {$y=\sin\!\left(\tfrac{1}{x}\right),\; x>0$};

  % salinim oku
  \draw[-{Stealth[length=2.4mm]}, red!70!black, thick] (0.95,-1.35) -- (0.2,-1.35);
  \node[red!70!black, etiket, align=center] at (1.7,-1.45)
    {$x\to 0^+$ iken sonsuz salınım};

  \node[etiket, align=center, fill=blue!5, rounded corners, inner sep=4pt]
    at (1.5,-2.35)
    {\textbf{Bağlantılı} ama \textbf{yol-bağlantılı değil:} sağdaki eğri ile soldaki
     dikey segmenti\\ birleştiren \emph{sürekli} bir yol yoktur (salınım yakınsamayı engeller).};
\end{tikzpicture}

  \caption{Topolojist sinüs eğrisi: $x\to 0^+$ iken sonsuz salınım. Eğri ile
    dikey limit segmenti arasında sürekli yol kurulamaz.}
\end{figure}

%-------------------------------------------------------------------
\section{Teoremler}

\begin{theorem}
Bağlantılı küme süreklilik altında bağlantılıdır:
$f: X \to Y$ sürekli, $X$ bağlantılı $\Rightarrow f(X)$ bağlantılı.
\end{theorem}

\begin{proof}[İspat eskizi]
$f(X)$ bağlantısız olsaydı, $f(X) = U \sqcup V$ biçiminde boş olmayan, ayrık, açık
iki kümeye ayrılırdı. Süreklilikten $f^{-1}(U)$ ve $f^{-1}(V)$ açıktır; ayrıktır;
birleşimleri $X$'tir ve ikisi de boş değildir. Bu, $X$'in bir clopen ayrılmasıdır
--- $X$'in bağlantılılığıyla çelişir. Demek ki $f(X)$ bağlantılıdır.
\end{proof}

\begin{theorem}
Yol-bağlantılı $\Rightarrow$ bağlantılı. Tersi genel olarak yanlış.
(Karşı örnek: Topolojist sinüs eğrisi.)
\end{theorem}

\begin{proof}[İspat eskizi]
$X$ yol-bağlantılı ama bağlantısız olsaydı, $X = U \sqcup V$ clopen ayrılması
olurdu. $x \in U$, $y \in V$ alıp $f:[0,1]\to X$, $f(0)=x$, $f(1)=y$ sürekli
yolunu kuralım. $[0,1]$ bağlantılı olduğundan görüntü $f([0,1])$ de bağlantılıdır;
ama bu görüntü $U$ ve $V$ tarafından gerçek bir clopen ayrılmaya uğrar --- çelişki.
Ters yön yanlıştır: topolojist sinüs eğrisi bağlantılı, ama yol-bağlantılı değildir.
\end{proof}

\begin{theorem}[Ara Değer Teoremi]
$f: X \to \mathbb{R}$ sürekli, $X$ bağlantılı $\Rightarrow f(X)$ bir aralıktır.
\end{theorem}

\begin{proof}[İspat eskizi]
$f(X) \subseteq \mathbb{R}$ bağlantılıdır (önceki teorem). $\mathbb{R}$'de
bağlantılı kümeler tam olarak aralıklardır; demek ki $f(X)$ bir aralıktır.
$a, b \in f(X)$ ise aralık özelliğinden $[a,b] \subseteq f(X)$; yani her
$c \in [a,b]$ için $f^{-1}(c) \neq \emptyset$.
\end{proof}

\begin{theorem}
Gerçek doğru $\mathbb{R}$'de bağlantılı alt kümeler tam olarak aralıklardır.
\end{theorem}

\begin{proof}[İspat eskizi]
(Aralık $\Rightarrow$ bağlantılı) Bir $I$ aralığı $U \sqcup V$ clopen ayrılmasına
uğrasaydı, $a\in U$, $b\in V$ ($a<b$) alıp $c=\sup\{x\in U : x\le b\}$ tanımlanır;
$c$'nin hangi parçada olduğu hem açıklık hem kapalılıkla çelişir. (Bağlantılı
$\Rightarrow$ aralık) $S$ aralık değilse $a<c<b$, $a,b\in S$, $c\notin S$ olan bir
$c$ vardır; o zaman $(-\infty,c)\cap S$ ve $(c,\infty)\cap S$ bir clopen ayrılma verir.
\end{proof}

%-------------------------------------------------------------------
\section{Algoritmalar}

\subsection{Sonlu Bağlantılılık — $O(|\tau|^2)$}

\begin{algorithm}
\caption{BagliMi$(X, \tau)$}
\begin{algorithmic}[1]
\For{each non-empty $A \subsetneq X$}
    \If{$A \in \tau$ \textbf{and} $X \setminus A \in \tau$}
        \State \Return False \Comment{Clopen bölme bulundu}
    \EndIf
\EndFor
\State \Return True
\end{algorithmic}
\end{algorithm}

Karmaşıklık: $O(|\tau|^2)$ — clopen küme taraması.

%-------------------------------------------------------------------
\section{pytop API}

\begin{lstlisting}[language=Python]
from pytop import (
    is_connected,
    is_path_connected,
    is_locally_connected,
    is_totally_disconnected,
)
\end{lstlisting}

%-------------------------------------------------------------------
\section{Örnekler}

\subsection*{Örnek 5.1 — Sierpiński: Bağlantılı}

\begin{lstlisting}[language=Python]
s = sierpinski_space()
print("connected?", is_connected(s).status)
print("path_connected?", is_path_connected(s).status)
\end{lstlisting}

\begin{verbatim}
connected?       true
path_connected?  unknown
\end{verbatim}

\subsection*{Örnek 5.2 — Ayrık Topoloji: Tamamen Bağlantısız}

\begin{lstlisting}[language=Python]
d = discrete_topology(1, 2, 3)
print("connected?", is_connected(d).status)
print("totally_disconn?", is_totally_disconnected(d).status)
\end{lstlisting}

\begin{verbatim}
connected?          false
totally_disconn?    true
\end{verbatim}

\subsection*{Örnek 5.3 — Gerçek Doğru ve [0,1]}

\begin{lstlisting}[language=Python]
rl = real_line_metric()
ui = closed_unit_interval_metric()
print("R connected?", is_connected(rl).status)
print("[0,1] path_connected?", is_path_connected(ui).status)
\end{lstlisting}

\begin{verbatim}
R connected?          true
[0,1] path_connected? true
\end{verbatim}

\subsection*{Örnek 5.4 — Ark / Yol / Bağlantı Hiyerarşisi}

Hiyerarşi \textbf{Ark $\Rightarrow$ Yol $\Rightarrow$ Bağlantılı} idi. pytop bu üç
düzeyi ayrı ayrı sorgular; bazı uzaylarda bir düzey \texttt{true}, diğeri
\texttt{unknown} döner --- pytop yalnızca \emph{kanıtlayabildiğini} bildirir.

\begin{lstlisting}[language=Python]
from pytop import (indiscrete_topology, real_line_metric, discrete_topology,
                   is_arc_connected, is_path_connected, is_connected)

for name, sp in [("Indiscrete(2)", indiscrete_topology(1, 2)),
                 ("R", real_line_metric()),
                 ("Discrete(3)", discrete_topology(1, 2, 3))]:
    print(f"{name:14s} connected={is_connected(sp).status:6s} "
          f"path={is_path_connected(sp).status:8s} arc={is_arc_connected(sp).status}")
\end{lstlisting}

\begin{verbatim}
Indiscrete(2)  connected=true   path=unknown  arc=true
R              connected=true   path=true     arc=unknown
Discrete(3)    connected=false  path=unknown  arc=false
\end{verbatim}

\subsection*{Örnek 5.5 — Tek Çağrıyla Çoklu Sorgu: \texttt{analyze\_connectedness}}

\texttt{analyze\_connectedness(space, property\_name)} tek bir dağıtıcı sunar;
geçerli adlar: \texttt{connected}, \texttt{path\_connected},
\texttt{locally\_connected}, \texttt{totally\_disconnected}, \texttt{arc\_connected}.

\begin{lstlisting}[language=Python]
from pytop import analyze_connectedness, sierpinski_space

s = sierpinski_space()
for prop in ["connected", "path_connected", "locally_connected",
             "totally_disconnected", "arc_connected"]:
    print(f"  {prop:22s}: {analyze_connectedness(s, prop).status}")
\end{lstlisting}

\begin{verbatim}
  connected             : true
  path_connected        : unknown
  locally_connected     : unknown
  totally_disconnected  : false
  arc_connected         : false
\end{verbatim}

Sierpiński bağlantılıdır ve tamamen bağlantısız \emph{değildir} --- tutarlı,
çünkü (tek noktadan büyük) bağlantılı bir uzay tamamen bağlantısız olamaz.

\subsection*{Örnek 5.6 — Clopen Ayrılmayı Elle Kurmak}

$X=\{1,2,3,4\}$ üzerinde $U=\{1,2\}$, $V=\{3,4\}$ taban alınırsa ikisi de clopen
olur; $X = U \sqcup V$ bir ayrılmadır. İç içe açıklar ise bölme vermez.

\begin{lstlisting}[language=Python]
from pytop import make_topology, is_connected, is_totally_disconnected

X = make_topology({1, 2, 3, 4}, {1, 2}, {3, 4})       # clopen bölme
print("split connected?", is_connected(X).status)
print("split totally_disconn?", is_totally_disconnected(X).status)

Y = make_topology({1, 2, 3}, {1}, {1, 2}, {1, 2, 3})  # ic ice, bölme yok
print("nested connected?", is_connected(Y).status)
\end{lstlisting}

\begin{verbatim}
split connected?        false
split totally_disconn?  false
nested connected?       true
\end{verbatim}

İlk uzay bağlantısız (clopen bölme var) ama \emph{tamamen} bağlantısız değil:
$\{1,2\}$ içinde clopen bölme yoktur. İkinci uzay iç içe açıklardan oluşur;
gerçek clopen bölme kurulamaz, bu yüzden bağlantılıdır.

%-------------------------------------------------------------------
\section{Alıştırmalar}

\subsection*{Kodlama}

\begin{enumerate}[label=K\arabic*.]
\item \texttt{make\_topology(\{1,2,3\},\{1\},\{2,3\})} için \texttt{is\_connected} hesaplayın.
\item \texttt{finite\_chain\_space(4)} zincirinin bağlantılı olup olmadığını kontrol edin.
\item İki noktalı ayrık ve indiscrete uzaylar üzerinde \texttt{is\_connected},
      \texttt{is\_path\_connected} karşılaştırın.
\item \texttt{indiscrete\_topology(1,2)} ve \texttt{real\_line\_metric()} için
      \texttt{is\_arc\_connected} ve \texttt{is\_path\_connected} çıktılarını
      karşılaştırın. Hangi uzayda hangi düzey \texttt{unknown} döner?
      \ipucu{Örnek 5.4.}
\item \texttt{discrete\_topology(1,2,3)} için
      \texttt{analyze\_connectedness(d, "connected")} ve
      \texttt{analyze\_connectedness(d, "totally\_disconnected")} çağırın; iki
      sonucun neden tutarlı olduğunu açıklayın. \ipucu{Örnek 5.5.}
\end{enumerate}

\subsection*{Teori}

\begin{enumerate}[label=T\arabic*.]
\item Bağlantılı + sürekli $\Rightarrow$ görüntü bağlantılı teoremini ispatlayın.
\item Yol-bağlantılı $\Rightarrow$ bağlantılı implicasyonunu ispatlayın.
\item Topolojist sinüs eğrisinin neden bağlantılı \emph{ama} yol-bağlantılı
      olmadığını açıklayın: dikey segmentten eğriye sürekli bir yol varsayıp
      çelişkiye ulaşın. \ipucu{\S1 karşı-örnek kutusu + sinüs eğrisi figürü.}
\end{enumerate}

\chapter{Sayılabilirlik Aksiyomları}

Sayılabilirlik aksiyomları, bir topolojik uzaydaki ``bilgi''nin sayılabilir büyüklükte
yapılarla tanımlanıp tanımlanamayacağını araştırır.

%-------------------------------------------------------------------
\section{Konu}

\subsection{Sayılabilirlik Aksiyomları}

\begin{table}[h]
\centering
\begin{tabular}{ll}
\toprule
\textbf{Aksiyom} & \textbf{Tanım} \\
\midrule
Birinci sayılabilir & Her $x \in X$ için sayılabilir yerel baz \\
İkinci sayılabilir & Tüm topoloji için sayılabilir global baz \\
Ayrılabilir & Sayılabilir yoğun alt küme: $\overline{D} = X$ \\
Lindelöf & Her açık örtünün sayılabilir alt-örtüsü \\
\bottomrule
\end{tabular}
\caption{Sayılabilirlik aksiyomları}
\end{table}

\textbf{Sıralamalar:}
\begin{itemize}
  \item $2^{\text{nd}}$ sayılabilir $\Rightarrow$ $1^{\text{st}}$ sayılabilir
  \item $2^{\text{nd}}$ sayılabilir $\Rightarrow$ ayrılabilir $\wedge$ Lindelöf
  \item Metrik $\Rightarrow$ $1^{\text{st}}$ sayılabilir
  \item Ayrılabilir metrik $\Rightarrow$ $2^{\text{nd}}$ sayılabilir
\end{itemize}

\begin{figure}[ht]
  \centering
  \begin{tikzpicture}[font=\small, node distance=10mm,
    aks/.style={draw, rounded corners, fill=blue!6, inner sep=5pt, align=center}]
  \node[aks] (sc) {2.\ sayılabilir\\ (sayılabilir baz)};
  \node[aks, above right=10mm and 16mm of sc] (fc) {1.\ sayılabilir\\ (sayılabilir yerel baz)};
  \node[aks, right=24mm of sc] (sep) {ayrılabilir\\ (yoğun say.\ alt küme)};
  \node[aks, below right=10mm and 16mm of sc] (lin) {Lindelöf\\ (say.\ alt-örtü)};
  \draw[-{Stealth[length=2.6mm]}, thick] (sc) -- (fc);
  \draw[-{Stealth[length=2.6mm]}, thick] (sc) -- (sep);
  \draw[-{Stealth[length=2.6mm]}, thick] (sc) -- (lin);
  \node[aks, fill=green!7, left=22mm of sc, align=center] (met) {metrik\\ uzay};
  \draw[-{Stealth[length=2.6mm]}, thick] (met) -- (fc.west);
  \draw[red!70!black, -{Stealth[length=2.2mm]}, dashed]
    (fc.east) to[bend left=22] node[right=2pt, font=\scriptsize, align=center]
    {tersi geçmez:\\ Sorgenfrey doğrusu\\ (1., ayrıl., Lindelöf\\ ama 2.\ sayılabilir değil)} (sep.east);
\end{tikzpicture}

  \caption{Sayılabilirlik aksiyomları arasındaki implikasyonlar: düz oklar
    $2^{\text{nd}}$ sayılabilirden $1^{\text{st}}$ sayılabilir, ayrılabilir ve
    Lindelöf'e gider; metrik uzaylar $1^{\text{st}}$ sayılabilirdir. Kesikli
    kırmızı ok terslerin geçmediği kanonik karşı-örneği (Sorgenfrey doğrusu)
    işaretler.}
  \label{fig:ch09:implikasyon}
\end{figure}

\begin{sezgi}
Sayılabilirlik aksiyomları, bir uzayı ``sayılabilir bir sözlükle'' anlatabilme
sorusudur. $2^{\text{nd}}$ sayılabilirlik uzayın \emph{tüm} açıklarını sayılabilir
bir baz listesinden kurar --- tek bir küresel sözlük. $1^{\text{st}}$ sayılabilirlik
ise yalnızca \emph{her nokta için ayrı} sayılabilir bir komşuluk listesi ister ---
yerel sözlükler. Küresel liste her noktaya yerel liste verdiğinden
$2^{\text{nd}}\Rightarrow 1^{\text{st}}$; tersi her zaman doğru değildir.
\end{sezgi}

\begin{figure}[ht]
  \centering
  \begin{tikzpicture}[font=\small]
  % carrier box
  \draw[rounded corners=8pt, thick] (-0.4,-1.4) rectangle (6.6,2.2);
  \node at (6.15,1.9) {$X$};
  % nested neighborhoods around x  (countable local base)
  \fill (2.1,0.4) circle (1.8pt) node[below=2pt, font=\scriptsize] {$x$};
  \draw[blue!70!black, thick, fill=blue!12, fill opacity=0.45] (2.1,0.4) circle (1.55);
  \draw[blue!75!black, thick, fill=blue!18, fill opacity=0.55] (2.1,0.4) circle (1.05);
  \draw[blue!80!black, thick, fill=blue!26, fill opacity=0.65] (2.1,0.4) circle (0.6);
  \node[blue!70!black, font=\scriptsize] at (2.1,1.72) {$B_1$};
  \node[blue!75!black, font=\scriptsize] at (2.1,1.2)  {$B_2$};
  \node[blue!80!black, font=\scriptsize] at (2.1,0.85) {$B_3$};
  \node[align=left, anchor=west, font=\scriptsize] at (4.0,0.4)
    {$\{B_n\}_{n\in\mathbb{N}}$\\ sayılabilir\\ \emph{yerel baz}:\\ $x$'in her\\ komşuluğu bir\\ $B_n$ içerir.};
\end{tikzpicture}

  \caption{Bir $x$ noktası çevresinde iç içe küçülen $B_1\supseteq B_2\supseteq B_3$
    komşulukları: sayılabilir bir \emph{yerel baz}. $x$'in her komşuluğu bir $B_n$
    içerir.}
  \label{fig:ch09:yerelbaz}
\end{figure}

\begin{karsiornek}
\textbf{Sorgenfrey doğrusu.} $[a,b)$ biçimli yarı-açık aralıklarla üretilen
Sorgenfrey doğrusu $1^{\text{st}}$ sayılabilir, ayrılabilir \emph{ve} Lindelöf'tür
ama $2^{\text{nd}}$ sayılabilir \emph{değildir}. Her $x$ için
$\{[x, x+1/n) : n\in\mathbb{N}\}$ sayılabilir yerel baz verir; $\mathbb{Q}$
yoğundur. Ama herhangi bir baz, her $x$ için tam $x$'te başlayan bir eleman
içermek zorundadır; farklı $x$'ler farklı eleman gerektirir, dolayısıyla baz en
az continuum büyüklükte --- sayılamaz. Yani
``$2^{\text{nd}}\Leftarrow 1^{\text{st}}\wedge$ ayrılabilir $\wedge$ Lindelöf''
çıkarımı geçersizdir. \texttt{sorgenfrey\_line\_like()} ile test edilir.
\end{karsiornek}

%-------------------------------------------------------------------
\section{Teoremler}

\begin{theorem}
$2^{\text{nd}}$ sayılabilir $\Rightarrow$ ayrılabilir $\wedge$ Lindelöf.
\begin{proof}[İspat eskizi]
$\mathcal{B}=\{B_n\}$ sayılabilir baz olsun. \textbf{Ayrılabilir:} her boş olmayan
$B_n$'den bir $d_n$ seç; $D=\{d_n\}$ sayılabilirdir ve her boş olmayan açık $U$ bir
$B_n\subseteq U$ içerdiğinden $d_n\in U\cap D$ --- yani $\overline{D}=X$.
\textbf{Lindelöf:} keyfi açık örtü $\{U_\alpha\}$'da kullanılan her noktayı içeren
bir $B_n$'i, onu kapsayan bir $U_\alpha$ ile eşle; kullanılan $B_n$'ler sayılabilir
olduğundan seçilen $U_\alpha$'lar sayılabilir bir alt-örtü verir.
\end{proof}
\end{theorem}

\begin{theorem}
Metrik uzay $\Rightarrow$ $1^{\text{st}}$ sayılabilir.
\begin{proof}[Kanıt İskeleti]
Her $x$ için $\{B(x, 1/n) : n \in \mathbb{N}\}$ sayılabilir yerel bazdır.
\end{proof}
\end{theorem}

\begin{theorem}
Ayrılabilir + metrik $\Rightarrow$ $2^{\text{nd}}$ sayılabilir.
\begin{proof}[Kanıt İskeleti]
$D \subseteq X$ yoğun sayılabilir; $\{B(d, 1/n) : d \in D,\, n \in \mathbb{N}\}$ sayılabilir bazdır.
\end{proof}
\end{theorem}

\begin{theorem}[Sorgenfrey Karşı Örneği]
Sorgenfrey doğrusu ayrılabilir ve Lindelöf'tür ama $2^{\text{nd}}$ sayılabilir
değildir. (Yukarıdaki karşı-örnek kutusuna bakın.)
\end{theorem}

\begin{theorem}[Sayılamaz Ayrık Uzay --- İkinci Karşı Örnek]
Sayılamaz taşıyıcı üzerindeki ayrık topoloji $1^{\text{st}}$ sayılabilirdir ama
ayrılabilir, Lindelöf ve $2^{\text{nd}}$ sayılabilir \emph{değildir}.
\begin{proof}[İspat eskizi]
Her $\{x\}$ açık olduğundan $\{\{x\}\}$ tek elemanlı yerel bazdır
($1^{\text{st}}$ sayılabilir). Yoğun bir küme her $\{x\}$'i kesmek zorunda olduğundan
taşıyıcının tamamını içermelidir; taşıyıcı sayılamaz olduğundan ayrılabilir değildir.
$\{\{x\}\}_{x\in X}$ örtüsünün sayılabilir alt-örtüsü yoktur (Lindelöf değil) ve her
baz tüm $\{x\}$'leri içermek zorunda olduğundan sayılamaz ($2^{\text{nd}}$ sayılabilir
değil). \texttt{uncountable\_discrete\_space()} ile doğrulanır.
\end{proof}
\end{theorem}

\begin{figure}[ht]
  \centering
  \begin{tikzpicture}[font=\small]
  % real line
  \draw[thick, -{Stealth[length=2.4mm]}] (-0.4,0) -- (8.6,0) node[right] {$\mathbb{R}$};
  \foreach \x in {0,1,2,3,4,5,6,7,8}
    \draw[thick] (\x,-0.08) -- (\x,0.08);
  \node[font=\scriptsize] at (0,-0.32) {$0$};
  \node[font=\scriptsize] at (4,-0.32) {$4$};
  \node[font=\scriptsize] at (8,-0.32) {$8$};
  % sample rational points (dense countable subset Q)
  \foreach \x in {0.3,0.7,1.25,1.6,2.1,2.55,3.0,3.4,3.85,4.3,4.7,5.2,5.65,6.1,6.5,6.95,7.4,7.8}
    \fill[red!75!black] (\x,0) circle (1.3pt);
  % a target open interval
  \draw[blue!70!black, thick] (4.95,0.45) -- (5.95,0.45);
  \draw[blue!70!black, thick] (4.95,0.35) -- (4.95,0.55);
  \draw[blue!70!black, thick] (5.95,0.35) -- (5.95,0.55);
  \node[blue!70!black, font=\scriptsize] at (4.55,0.78) {$(a,b)$};
  \draw[red!75!black, -{Stealth}, thick] (5.65,1.45) -- (5.45,0.55);
  \node[red!75!black, font=\scriptsize, align=center] at (6.35,1.45) {içinde bir\\ $q\in\mathbb{Q}$ var};
  \node[align=center, font=\scriptsize] at (4.0,-1.0)
    {$\mathbb{Q}$ (kırmızı noktalar) sayılabilir ve $\overline{\mathbb{Q}}=\mathbb{R}$:\\
     her boş olmayan açık aralık bir rasyonel içerir $\Rightarrow$ $\mathbb{R}$ ayrılabilirdir.};
\end{tikzpicture}

  \caption{Ayrılabilirlik: $\mathbb{Q}$ (kırmızı noktalar) sayılabilirdir ve
    $\overline{\mathbb{Q}}=\mathbb{R}$. Her boş olmayan açık aralık $(a,b)$ bir
    rasyonel içerdiğinden $\mathbb{R}$ ayrılabilirdir.}
  \label{fig:ch09:ayrilabilirlik}
\end{figure}

%-------------------------------------------------------------------
\section{Algoritmalar}

\subsection{Sonlu Uzayda Sayılabilirlik}

Sonlu uzayda her aksiyom trivially sağlanır. $O(1)$.

\subsection{Sonsuz Uzayda: Tag-Tabanlı Çıkarım}

pytop sembolik uzaylarda tag koridorlarından sayılabilirlik çıkarır:
\begin{itemize}
  \item \texttt{'second\_countable'} $\in$ tags $\Rightarrow$ tüm aksiyomlar
  \item \texttt{'metric'} $\wedge$ \texttt{'separable'} $\Rightarrow$ \texttt{'second\_countable'}
\end{itemize}

%-------------------------------------------------------------------
\section{pytop API}

\begin{lstlisting}[language=Python]
from pytop import (
    is_first_countable,
    is_second_countable,
    is_separable,
    is_lindelof,
)
# Yardimcilar:
from pytop import countability_report, is_dense_subset
\end{lstlisting}

\texttt{countability\_report(space)} dört aksiyomu bir \texttt{dict} içinde
toplar; \texttt{is\_dense\_subset(space, subset)} ayrılabilirliği somut bir
tanıkla gösterir.

%-------------------------------------------------------------------
\section{Örnekler}

\subsection*{Örnek 5.1 — Gerçek Doğru: Tüm 4 Aksiyom}

\begin{lstlisting}[language=Python]
rl = real_line_metric()
print("1st countable?", is_first_countable(rl).status)
print("2nd countable?", is_second_countable(rl).status)
print("separable?", is_separable(rl).status)
print("lindelof?", is_lindelof(rl).status)
\end{lstlisting}

\begin{verbatim}
1st countable?  true
2nd countable?  true
separable?      true
lindelof?       true
\end{verbatim}

\subsection*{Örnek 5.5 — Sorgenfrey Doğrusu: 2nd Countable $\times$}

\begin{lstlisting}[language=Python]
from pytop import sorgenfrey_line_like
sl = sorgenfrey_line_like()
print("1st countable?", is_first_countable(sl).status)
print("2nd countable?", is_second_countable(sl).status)
print("separable?", is_separable(sl).status)
\end{lstlisting}

\begin{verbatim}
1st countable?   true
2nd countable?   false
separable?       true
\end{verbatim}

\subsection*{Örnek 5.6 — Toplu Rapor + Yoğunluk Tanığı}

\begin{lstlisting}[language=Python]
from pytop import countability_report, is_dense_subset
from pytop import sorgenfrey_line_like, discrete_topology

rep = countability_report(sorgenfrey_line_like())
for key in ['first_countable','second_countable','separable','lindelof']:
    print(f"  {key:<17}: {rep[key].status}")

d = discrete_topology(1, 2, 3)
print("dense {1,2,3}?", is_dense_subset(d, {1, 2, 3}).status)
print("dense {1,2}?  ", is_dense_subset(d, {1, 2}).status)
\end{lstlisting}

\begin{verbatim}
  first_countable  : true
  second_countable : false
  separable        : true
  lindelof         : true
dense {1,2,3}? true
dense {1,2}?   false
\end{verbatim}

\subsection*{Örnek 5.7 — Sayılamaz Ayrık Uzay: İkinci Karşı Örnek}

\begin{lstlisting}[language=Python]
from pytop import uncountable_discrete_space
from pytop import (is_first_countable, is_second_countable,
                   is_separable, is_lindelof)

u = uncountable_discrete_space()
print("1st countable?", is_first_countable(u).status)
print("2nd countable?", is_second_countable(u).status)
print("separable?    ", is_separable(u).status)
print("lindelof?     ", is_lindelof(u).status)
\end{lstlisting}

\begin{verbatim}
1st countable? true
2nd countable? false
separable?     false
lindelof?      false
\end{verbatim}

Her tekil $\{x\}$ açık olduğundan yerel baz tek elemanlıdır ($1^{\text{st}}$
sayılabilir), ama yoğun sayılabilir alt küme ve sayılabilir alt-örtü yoktur.

%-------------------------------------------------------------------
\section{Alıştırmalar}

\subsection*{Kodlama}

\begin{enumerate}[label=K\arabic*.]
\item \texttt{finite\_chain\_space(5)} üzerinde \texttt{is\_first\_countable} ve
      \texttt{is\_second\_countable} hesaplayın.
\item İki noktalı ayrık ve indiscrete uzaylar üzerinde \texttt{is\_separable} karşılaştırın.
\item \texttt{integers\_discrete()} için $2^{\text{nd}}$ sayılabilir mi?
\item \texttt{uncountable\_discrete\_space()} üzerinde dört aksiyomu
      \texttt{countability\_report} ile hesaplayıp \texttt{integers\_discrete()}
      ile karşılaştırın. Hangi aksiyom(lar)da ayrışıyorlar?
      \ipucu{Fark, taşıyıcının sayılabilir olup olmamasındadır.}
\end{enumerate}

\subsection*{Teori}

\begin{enumerate}[label=T\arabic*.]
\item $2^{\text{nd}}$ sayılabilir $\Rightarrow$ ayrılabilir olduğunu ispatlayın.
\item Sorgenfrey doğrusunun $2^{\text{nd}}$ sayılabilir olmadığını gösterin.
\item ``$1^{\text{st}}$ sayılabilir $\Rightarrow$ ayrılabilir'' çıkarımının yanlış
      olduğunu bir karşı-örnekle gösterin.
      \ipucu{Sayılamaz ayrık uzay; yoğun küme her $\{x\}$'i kesmek zorundadır.}
\end{enumerate}

\chapter{Sürekli Fonksiyonlar ve Homeomorfizmalar}

Topolojik uzaylar arasındaki fonksiyonlarda süreklilik, açık kümelerin
geri çekiminin açık olması koşuluna dayanır.

%-------------------------------------------------------------------
\section{Konu}

\begin{sezgi}
Sürekliliği ``açıklığın geri taşınması'' gibi düşünün: hedef uzayda hangi
açık kümeye bakarsanız bakın, onun \emph{geri çekimi} (preimage) kaynak
uzayda yine açık olmalı. Analizden bildiğiniz $\varepsilon$--$\delta$ tanımı
bunun metrik özelidir; topolojide $\varepsilon$--$\delta$ yerine doğrudan
açık kümelerin dilini kullanırız. Homeomorfizma ise bu taşımanın \emph{iki
yönde} de bozulmadan işlemesi --- yani iki uzayın topolojik olarak ``aynı''
olmasıdır.
\end{sezgi}

\begin{definition}[Sürekli Fonksiyon]
$f: (X, \tau_X) \to (Y, \tau_Y)$ fonksiyonu \textbf{sürekli} ise:
\[
\forall V \in \tau_Y:\; f^{-1}(V) \in \tau_X.
\]
\end{definition}

\begin{figure}[ht]
  \centering
  \begin{tikzpicture}[font=\small]
  % --- domain X ---
  \draw[rounded corners=8pt, thick] (-0.4,-1.3) rectangle (3.0,1.5);
  \node at (-0.05,1.2) {$X$};
  \draw[blue!70!black, thick, fill=blue!10] (1.3,-0.1) ellipse [x radius=0.95, y radius=0.7];
  \node[blue!70!black] at (1.3,-0.1) {$f^{-1}(V)$};
  % --- codomain Y ---
  \draw[rounded corners=8pt, thick] (5.0,-1.3) rectangle (8.4,1.5);
  \node at (8.05,1.2) {$Y$};
  \draw[violet!80!black, thick, fill=violet!10] (6.7,-0.1) ellipse [x radius=0.85, y radius=0.6];
  \node[violet!80!black] at (6.7,-0.1) {$V$};
  % --- arrow f ---
  \draw[-{Stealth}, thick] (3.25,0.55) to[bend left=12]
    node[above, pos=0.5] {$f$} (4.75,0.55);
  % --- caption note ---
  \node[font=\scriptsize, align=center] at (4.0,-1.85)
    {$V$ açık $\Rightarrow f^{-1}(V)$ açık olmalı (her açık $V$ için)};
\end{tikzpicture}

  \caption{Süreklilik: hedef uzaydaki her açık $V$ için geri çekim
    $f^{-1}(V)$ kaynak uzayda açık olmalıdır.}
  \label{fig:ch10:sureklilik}
\end{figure}

Eşdeğer koşullar:
\begin{itemize}
  \item Her kapalı $F \subseteq Y$ için $f^{-1}(F)$ kapalıdır.
  \item Her $x \in X$ ve $f(x)$'in her komşuluğu $W$ için $f^{-1}(W)$, $x$'in komşuluğudur.
\end{itemize}

\begin{definition}[Homeomorfizma]
$f: X \to Y$ \textbf{homeomorfizma} ise: bijektif, $f$ sürekli, $f^{-1}$ sürekli.
$(X, \tau_X) \cong (Y, \tau_Y)$ ile gösterilir.
\end{definition}

\begin{figure}[ht]
  \centering
  \begin{tikzpicture}[font=\small]
  % --- space X ---
  \draw[blue!70!black, thick, fill=blue!8] (0,0) ellipse [x radius=1.2, y radius=0.95];
  \node[blue!70!black] at (0,1.3) {$(X,\tau_X)$};
  % --- space Y ---
  \draw[violet!80!black, thick, fill=violet!8] (5.0,0) ellipse [x radius=1.2, y radius=0.95];
  \node[violet!80!black] at (5.0,1.3) {$(Y,\tau_Y)$};
  % --- two arrows ---
  \draw[-{Stealth}, thick] (1.3,0.35) to[bend left=22]
    node[above, pos=0.5] {$f$ sürekli} (3.7,0.35);
  \draw[-{Stealth}, thick] (3.7,-0.35) to[bend left=22]
    node[below, pos=0.5] {$f^{-1}$ sürekli} (1.3,-0.35);
  % --- caption ---
  \node[font=\scriptsize, align=center] at (2.5,-1.7)
    {bijektif, iki yön de sürekli $\Rightarrow (X,\tau_X)\cong(Y,\tau_Y)$};
\end{tikzpicture}

  \caption{Homeomorfizma: bijektif bir eşleme hem $f$ hem $f^{-1}$ yönünde
    sürekliyse iki uzay topolojik olarak özdeştir, $(X,\tau_X)\cong(Y,\tau_Y)$.}
  \label{fig:ch10:homeo}
\end{figure}

\begin{karsiornek}
Süreklilik \emph{yöne duyarlıdır}. Aynı taşıyıcı küme $\{0,1\}$ üzerinde
özdeşlik fonksiyonunu düşünün: ayrık topolojiden (ince) Sierpiński'ye (kaba)
\textbf{sürekli}, ama Sierpiński'den ayrığa \textbf{sürekli değil} --- çünkü
ayrıktaki açık $\{0\}$'ın geri çekimi Sierpiński'de açık değildir. Yani
``kaba $\to$ ince'' yönünde özdeşlik genelde süreklilik kaybeder
(Örnek~5.8'de \texttt{is\_continuous\_finite\_map} ile doğrulanır).
\end{karsiornek}

\begin{table}[h]
\centering
\begin{tabular}{ll}
\toprule
\textbf{Tür} & \textbf{Koşul} \\
\midrule
Sabit fonksiyon & $f(x) = c$; her zaman sürekli \\
Özdeşlik & $f(x) = x$; her zaman sürekli \\
Bileşke & $f$, $g$ sürekli $\Rightarrow g \circ f$ sürekli \\
Homeomorfizma & Bijektif; $f$ ve $f^{-1}$ sürekli \\
\bottomrule
\end{tabular}
\caption{Süreklilik hiyerarşisi}
\end{table}

%-------------------------------------------------------------------
\section{Teoremler}

\begin{theorem}
Her sabit fonksiyon süreklidir.
\begin{proof}[Kanıt İskeleti]
$f^{-1}(V) = X$ eğer $c \in V$, $\emptyset$ eğer $c \notin V$; her ikisi de açıktır.
\end{proof}
\end{theorem}

\begin{theorem}
$f: X \to Y$ ve $g: Y \to Z$ sürekli $\Rightarrow g \circ f$ süreklidir.
\begin{proof}[İspat eskizi]
$W \subseteq Z$ açık olsun. $g$ sürekli olduğundan $g^{-1}(W) \subseteq Y$
açıktır. $f$ sürekli olduğundan $f^{-1}(g^{-1}(W)) \subseteq X$ açıktır.
$(g\circ f)^{-1}(W) = f^{-1}(g^{-1}(W))$ olduğundan $g\circ f$'in her açık
kümenin geri çekimi açıktır; yani $g\circ f$ süreklidir.
\end{proof}
\end{theorem}

\begin{figure}[ht]
  \centering
  \begin{tikzpicture}[font=\small]
  \draw[blue!70!black, thick, fill=blue!8] (0,0) ellipse [x radius=0.95, y radius=0.85];
  \node[blue!70!black] at (0,1.15) {$X$};
  \draw[teal!70!black, thick, fill=teal!8] (3.5,0) ellipse [x radius=0.95, y radius=0.85];
  \node[teal!70!black] at (3.5,1.15) {$Y$};
  \draw[violet!80!black, thick, fill=violet!8] (7.0,0) ellipse [x radius=0.95, y radius=0.85];
  \node[violet!80!black] at (7.0,1.15) {$Z$};
  \draw[-{Stealth}, thick] (1.05,0.0) -- node[above] {$f$} (2.45,0.0);
  \draw[-{Stealth}, thick] (4.55,0.0) -- node[above] {$g$} (5.95,0.0);
  \draw[-{Stealth}, thick, red!70!black] (0.3,-0.75) to[bend right=24]
    node[below, pos=0.5] {$g\circ f$} (6.7,-0.75);
  \node[font=\scriptsize, align=center] at (3.5,-2.0)
    {$f$ ve $g$ sürekli $\Rightarrow g\circ f$ sürekli};
\end{tikzpicture}

  \caption{Bileşke süreklilik: $f$ ve $g$ sürekliyse $g\circ f$ de süreklidir;
    geri çekim zincirlenir, $(g\circ f)^{-1}(W) = f^{-1}(g^{-1}(W))$.}
  \label{fig:ch10:bileske}
\end{figure}

\begin{theorem}
$f: X \to Y$ sürekli, $X$ kompakt $\Rightarrow f(X)$ kompakttır.
\begin{proof}[İspat eskizi]
$f(X)$'in bir açık örtüsü $\{V_\alpha\}$ verilsin. $f$ sürekli olduğundan her
$f^{-1}(V_\alpha)$ $X$'te açıktır ve bunlar $X$'i örter. $X$ kompakt olduğundan
sonlu bir alt-örtü $f^{-1}(V_{\alpha_1}), \dots, f^{-1}(V_{\alpha_n})$ vardır.
O zaman $V_{\alpha_1}, \dots, V_{\alpha_n}$ $f(X)$'i örter: süreklilik açık
örtüyü geri çeker, kompaktlık sonlu alt-örtü verir, görüntüye geri itilir.
Demek ki $f(X)$ kompakttır.
\end{proof}
\end{theorem}

\begin{theorem}
$f: X \to Y$ sürekli, $X$ bağlantılı $\Rightarrow f(X)$ bağlantılıdır.
\end{theorem}

\begin{theorem}
Kompakt $X$'ten Hausdorff $Y$'ye sürekli bijeksiyon $\Rightarrow$ homeomorfizma.
\end{theorem}

%-------------------------------------------------------------------
\section{Algoritmalar}

\subsection{is\_continuous\_finite\_map — $O(|\tau_Y| \cdot |X|)$}

\begin{algorithm}
\caption{IsContinuous$(X, \tau_X, Y, \tau_Y, f)$}
\begin{algorithmic}[1]
\For{each $V \in \tau_Y$}
    \State preimage $\leftarrow \{x \in X : f(x) \in V\}$
    \If{preimage $\notin \tau_X$}
        \State \Return False
    \EndIf
\EndFor
\State \Return True
\end{algorithmic}
\end{algorithm}

%-------------------------------------------------------------------
\section{pytop API}

\begin{lstlisting}[language=Python]
from pytop import (
    is_continuous_finite_map,
    finite_homeomorphism_result,
    sierpinski_space,
    discrete_topology,
    indiscrete_topology,
    make_topology,
    make_set,
    empty_set,
)
\end{lstlisting}

\texttt{is\_continuous\_finite\_map(domain, domain\_topology, codomain, codomain\_topology, mapping)} $\to$ \texttt{bool}

\texttt{finite\_homeomorphism\_result(left, right)} $\to$ \texttt{Result}

%-------------------------------------------------------------------
\section{Örnekler}

\subsection*{Örnek 5.2 — Ayrık Topoloji: Her Fonksiyon Sürekli}

\begin{lstlisting}[language=Python]
d = discrete_topology(0, 1)
d_pts = list(d.carrier)
d_topo = [set(u) for u in d.topology]

f_id   = {0: 0, 1: 1}
f_swap = {0: 1, 1: 0}
print("ozdeslik:", is_continuous_finite_map(d_pts, d_topo, d_pts, d_topo, f_id))
print("swap:    ", is_continuous_finite_map(d_pts, d_topo, d_pts, d_topo, f_swap))
\end{lstlisting}

\begin{verbatim}
ozdeslik continuous: True
swap     continuous: True
\end{verbatim}

\subsection*{Örnek 5.3 — Sierpiński $\to$ Ayrık: Sürekli Değil}

\begin{lstlisting}[language=Python]
s = sierpinski_space()
s_pts = list(s.carrier)
s_topo = [set(u) for u in s.topology]
print("id: S->D:", is_continuous_finite_map(s_pts, s_topo, d_pts, d_topo, f_id))
print("id: D->S:", is_continuous_finite_map(d_pts, d_topo, s_pts, s_topo, f_id))
\end{lstlisting}

\begin{verbatim}
id: Sierpinski -> Disc continuous: False
id: Disc -> Sierpinski continuous: True
\end{verbatim}

Sierpiński $\tau = \{\emptyset, \{0,1\}, \{1\}\}$: $f^{-1}(\{0\}) = \{0\}$ Sierpiński'de açık değil.

\subsection*{Örnek 5.5 — Homeomorfizma Kontrolü}

\begin{lstlisting}[language=Python]
d2a = discrete_topology(1, 2)
d2b = discrete_topology('a', 'b')
ind2 = indiscrete_topology(1, 2)
print("D(1,2) ~ D(a,b):", finite_homeomorphism_result(d2a, d2b).status)
print("D(1,2) ~ S(0,1):", finite_homeomorphism_result(d2a, s).status)
print("D(1,2) ~ Ind(1,2):", finite_homeomorphism_result(d2a, ind2).status)
\end{lstlisting}

\begin{verbatim}
D(1,2) ~ D(a,b): true
D(1,2) ~ S(0,1): false
D(1,2) ~ Ind(1,2): false
\end{verbatim}

\subsection*{Örnek 5.7 — Bileşke Süreklilik: $g \circ f$}

\begin{lstlisting}[language=Python]
TX = make_topology(make_set(1, 2, 3), make_set(1), make_set(1, 2))
TY = make_topology(make_set('a', 'b', 'c'), make_set('a'), make_set('a', 'b'))
TZ = make_topology(make_set('p', 'q', 'r'), make_set('p'), make_set('p', 'q'))
X_pts, X_topo = list(TX.carrier), list(TX.topology)
Y_pts, Y_topo = list(TY.carrier), list(TY.topology)
Z_pts, Z_topo = list(TZ.carrier), list(TZ.topology)

f = {1: 'a', 2: 'b', 3: 'c'}
g = {'a': 'p', 'b': 'q', 'c': 'r'}
gof = {x: g[f[x]] for x in X_pts}
print("f continuous:    ", is_continuous_finite_map(X_pts, X_topo, Y_pts, Y_topo, f))
print("g continuous:    ", is_continuous_finite_map(Y_pts, Y_topo, Z_pts, Z_topo, g))
print("g of f continuous:", is_continuous_finite_map(X_pts, X_topo, Z_pts, Z_topo, gof))
\end{lstlisting}

\begin{verbatim}
f continuous:     True
g continuous:     True
g of f continuous: True
\end{verbatim}

İki sürekli fonksiyonun geri çekimleri zincirlenir:
$(g \circ f)^{-1}(W) = f^{-1}(g^{-1}(W))$ (Teorem~2.2'nin sayısal doğrulaması).

\subsection*{Örnek 5.8 — Süreklilik Yöne Duyarlı}

\begin{lstlisting}[language=Python]
fine = discrete_topology(0, 1)
coarse = sierpinski_space()
fine_pts, fine_topo = list(fine.carrier), list(fine.topology)
co_pts, co_topo = list(coarse.carrier), list(coarse.topology)
idmap = {0: 0, 1: 1}
print("id: discrete -> sierpinski:", is_continuous_finite_map(fine_pts, fine_topo, co_pts, co_topo, idmap))
print("id: sierpinski -> discrete:", is_continuous_finite_map(co_pts, co_topo, fine_pts, fine_topo, idmap))
\end{lstlisting}

\begin{verbatim}
id: discrete -> sierpinski: True
id: sierpinski -> discrete: False
\end{verbatim}

Ayrıktaki açık $\{0\}$'ın geri çekimi yine $\{0\}$'dır; bu Sierpiński
$\tau = \{\emptyset, \{1\}, \{0,1\}\}$ içinde açık olmadığından ``kaba $\to$
ince'' yönü süreklilik kaybeder.

%-------------------------------------------------------------------
\section{Alıştırmalar}

\subsection*{Kodlama}

\begin{enumerate}[label=K\arabic*.]
\item Özel bir topolojide $f(0)=0, f(1)=0, f(2)=1$ fonksiyonunun sürekli
      olup olmadığını kontrol edin.
\item Sierpiński uzayından kendisine giden dört fonksiyonu test edin.
\item \texttt{finite\_homeomorphism\_result} ile iki ayrık uzayın homeomorf olduğunu doğrulayın.
\item \texttt{make\_topology} ile $\{1,2,3\}$ üzerinde $\{\alpha\}$, $\{\alpha,\beta\}$
      açıklarıyla zincir topolojileri kurun; zinciri koruyan $f$ ve $g$ tanımlayıp
      $f$, $g$ ve $g\circ f$'in üçünün de sürekli olduğunu
      \texttt{is\_continuous\_finite\_map} ile doğrulayın. \ipucu{Teorem~2.2'nin sayısal kontrolü.}
\end{enumerate}

\subsection*{Teori}

\begin{enumerate}[label=T\arabic*.]
\item Her sabit fonksiyonun sürekli olduğunu ispatlayın.
\item $f \circ g$ bileşkesinin sürekli olduğunu ispatlayın.
\item $f: X \to Y$ sürekli ve $X$ kompakt ise $f(X)$'in kompakt olduğunu
      ispatlayın. \ipucu{$f(X)$'in açık örtüsünü $f$ ile geri çekin, $X$'in
      kompaktlığıyla sonlu alt-örtü alın, görüntüye geri itin.}
\end{enumerate}

\chapter{Alt Uzay ve Çarpım Topolojisi}

Alt uzay topolojisi bir uzayın alt kümesine miras kalan topolojiyi tanımlar.
Çarpım topolojisi ise iki uzayı birleştirerek yeni bir uzay kurar.

%-------------------------------------------------------------------
\section{Konu}

\begin{definition}[Alt Uzay Topolojisi]
$(X, \tau)$ uzayı ve $A \subseteq X$ verilsin. \textbf{Alt uzay topolojisi:}
\[
\tau_A = \{U \cap A : U \in \tau\}.
\]
\end{definition}

\begin{sezgi}
$A$'yı $X$'ten kesilen bir ``pencere'' gibi düşünün. $A$ üzerindeki açıkları
sıfırdan tanımlamayız; $X$'in hazır açıklarını $A$ ile keser, ne kadarı $A$'ya
düşüyorsa onu açık sayarız. Bu yüzden $\tau_A$, $\tau$'dan \emph{miras} alınır.
\end{sezgi}

\begin{figure}[ht]
\centering
\begin{tikzpicture}[scale=1.0, font=\small]
  % Ambient space X
  \draw[rounded corners=10pt, thick] (-0.5,-1.6) rectangle (6.7,2.2);
  \node at (6.25,1.85) {$X$};
  % Open set U of X
  \draw[blue!70!black, thick, fill=blue!8]
    plot[smooth cycle, tension=0.8] coordinates
    {(0.5,-0.4) (2.6,-1.0) (4.9,-0.2) (4.6,1.5) (2.5,1.9) (0.6,1.2)};
  \node[blue!70!black] at (4.25,1.45) {$U \in \tau$};
  % Subset A (a horizontal strip)
  \draw[red!70!black, thick, fill=red!7]
    (-0.2,-0.9) rectangle (6.4,0.1);
  \node[red!70!black] at (-0.05,-1.25) {$A$};
  % Intersection U cap A highlighted
  \draw[violet!70!black, very thick, fill=violet!22, fill opacity=0.7]
    (0.55,-0.9) -- (4.75,-0.9) -- (4.7,0.1) -- (0.62,0.1) -- cycle;
  \node[violet!75!black, align=center] at (2.65,-0.4)
    {$U \cap A$};
  \node[align=left, anchor=west, font=\scriptsize] at (7.0,0.3)
    {Alt uzay topolojisi:\\ $\tau_A = \{\,U \cap A : U \in \tau\,\}$.\\[2pt]
     $A$'nın açıkları, $X$'in\\ açıklarını $A$ ile kesmekle doğar.};
\end{tikzpicture}

\caption{Alt uzay topolojisi: $A \subseteq X$'te açıklar $U \cap A$ biçimindedir.}
\end{figure}

\begin{definition}[Çarpım Topolojisi]
$(X, \tau_X)$ ve $(Y, \tau_Y)$ verilsin. Çarpım topolojisi $X \times Y$ üzerinde
$\{U \times V : U \in \tau_X,\, V \in \tau_Y\}$ bazından üretilir.
\end{definition}

Projeksiyon $\pi_X: X \times Y \to X$ ve $\pi_Y: X \times Y \to Y$ süreklidir.

\begin{sezgi}
Çarpım topolojisinin bazı, düzlemdeki açık dikdörtgenlere benzer: bir ``yatay''
açık $U$ ile bir ``dikey'' açık $V$'yi çarparak $U \times V$ kutusunu kurarız.
Bu kutular tüm açıkların \emph{bazıdır} --- açıklar bu kutuların birleşimleridir.
\end{sezgi}

\begin{figure}[ht]
\centering
\begin{tikzpicture}[scale=1.0, font=\small]
  % Axes for X (horizontal) and Y (vertical)
  \draw[-{Stealth}] (-0.3,0) -- (5.6,0) node[right] {$X$};
  \draw[-{Stealth}] (0,-0.3) -- (0,4.4) node[above] {$Y$};
  % U interval on X axis
  \draw[blue!70!black, very thick] (1.4,0) -- (3.8,0);
  \node[blue!70!black, below=2pt] at (2.6,0) {$U \in \tau_X$};
  % V interval on Y axis
  \draw[red!70!black, very thick] (0,1.3) -- (0,3.5);
  \node[red!70!black, left=2pt, rotate=90, anchor=south] at (0,2.4) {$V \in \tau_Y$};
  % Basic box U x V
  \draw[violet!70!black, thick, fill=violet!15]
    (1.4,1.3) rectangle (3.8,3.5);
  \node[violet!75!black] at (2.6,2.4) {$U \times V$};
  % Dashed guide lines
  \draw[blue!60!black, dashed] (1.4,0) -- (1.4,1.3);
  \draw[blue!60!black, dashed] (3.8,0) -- (3.8,1.3);
  \draw[red!60!black, dashed] (0,1.3) -- (1.4,1.3);
  \draw[red!60!black, dashed] (0,3.5) -- (1.4,3.5);
  \node[align=left, anchor=west, font=\scriptsize] at (4.4,3.4)
    {Çarpım topolojisinin\\ baz elemanı: bir\\ dikdörtgen kutu\\ $U \times V$.};
\end{tikzpicture}

\caption{Çarpım topolojisinin baz elemanı: $U \times V$ dikdörtgen kutusu.}
\end{figure}

\begin{figure}[ht]
\centering
\begin{tikzpicture}[scale=1.0, font=\small]
  % Product box X x Y in the middle-top
  \draw[violet!70!black, thick, fill=violet!10] (2.0,1.6) rectangle (5.2,3.6);
  \node[violet!75!black] at (3.6,3.9) {$X \times Y$};
  \fill (3.6,2.6) circle (1.8pt);
  \node[font=\scriptsize, above right=0pt] at (3.6,2.6) {$(x,y)$};
  % X line bottom-left
  \draw[blue!70!black, thick, fill=blue!8] (-0.4,-0.4) rectangle (2.6,0.4);
  \node[blue!70!black] at (2.3,0.55) {$X$};
  \fill (1.1,0.0) circle (1.8pt) node[below=2pt, font=\scriptsize] {$x$};
  % Y line bottom-right
  \draw[red!70!black, thick, fill=red!8] (4.6,-0.4) rectangle (7.6,0.4);
  \node[red!70!black] at (4.9,0.55) {$Y$};
  \fill (6.1,0.0) circle (1.8pt) node[below=2pt, font=\scriptsize] {$y$};
  % projection arrows
  \draw[blue!70!black, -{Stealth}, thick]
    (3.2,1.5) to[bend right=18] node[left, font=\scriptsize, pos=0.55] {$\pi_X$} (1.0,0.55);
  \draw[red!70!black, -{Stealth}, thick]
    (4.2,1.5) to[bend left=18] node[right, font=\scriptsize, pos=0.55] {$\pi_Y$} (6.2,0.55);
  \node[font=\scriptsize, align=center] at (3.6,-1.15)
    {$\pi_X(x,y)=x,\quad \pi_Y(x,y)=y$ \; (her ikisi sürekli)};
\end{tikzpicture}

\caption{$\pi_X$ ve $\pi_Y$ projeksiyonları $(x,y)$ noktasını bileşenlerine düşürür.}
\end{figure}

\begin{karsiornek}
``$U \times V$ biçiminde olmayan açık yok'' demek yanlıştır. Birim diskte
$D(0,1) \times D(0,1)$ içindeki dairesel bir bölge, hiçbir tek $U \times V$
kutusuna eşit olmasa da açıktır; çünkü açıklar kutuların \emph{birleşimidir}.
$\{U \times V\}$ yalnızca \textbf{baz}'dır, topolojinin tamamı değil.
\end{karsiornek}

\begin{table}[h]
\centering
\begin{tabular}{lll}
\toprule
\textbf{Özellik} & \textbf{Alt uzay} & \textbf{Çarpım} \\
\midrule
Hausdorff & \checkmark & \checkmark \\
Kompaktlık & Yalnız kapalı alt uzay & \checkmark (Tychonoff) \\
Bağlantılılık & Hayır (genel) & \checkmark \\
T0/T1/T2 & \checkmark & \checkmark \\
\bottomrule
\end{tabular}
\caption{Topolojik özelliklerin kalıtımı}
\end{table}

%-------------------------------------------------------------------
\section{Teoremler}

\begin{theorem}
Hausdorff uzayın her alt uzayı Hausdorff'tur.
\begin{proof}[İspat eskizi]
$x \ne y$ noktaları $A \subseteq X$ içinde olsun. $X$ Hausdorff olduğundan
$X$'te ayrık açıklar $U \ni x$, $V \ni y$ vardır. O hâlde $U \cap A$ ve
$V \cap A$, $A$'da açık, $x$ ile $y$'yi ayıran komşuluklardır ve
$(U \cap A) \cap (V \cap A) = (U \cap V) \cap A = \emptyset$.
\end{proof}
\end{theorem}

\begin{theorem}
Kompakt uzayın kapalı alt uzayı kompakttır.
\begin{proof}[İspat eskizi]
$A \subseteq X$ kapalı, $X$ kompakt olsun. $A$'nın bir açık örtüsü
$\{U_\alpha \cap A\}$ verilsin. Her $U_\alpha$'yı açık olan $X \setminus A$ ile
birlikte alırsak $X$'in açık örtüsünü elde ederiz; $X$ kompakt olduğundan
sonlu alt-örtü vardır. $X \setminus A$'yı atınca $A$'yı örten sonlu sayıda
$U_\alpha \cap A$ kalır.
\end{proof}
\end{theorem}

\begin{theorem}
$X$ ve $Y$ bağlantılı $\Rightarrow X \times Y$ bağlantılıdır.
\begin{proof}[Kanıt İskeleti]
$\{x_0\} \times Y$ ve $X \times \{y_0\}$ bağlantılı dilimler; kesişimleri üzerinden birleşerek tüm çarpımı kapsar.
\end{proof}
\end{theorem}

\begin{theorem}[Tychonoff, $n=2$]
$X$ ve $Y$ kompakt $\Rightarrow X \times Y$ kompakttır.
\end{theorem}

\begin{theorem}
$X$ ve $Y$ Hausdorff $\Rightarrow X \times Y$ Hausdorff'tur.
\begin{proof}[İspat eskizi]
$(x_1,y_1) \ne (x_2,y_2)$ ise $x_1 \ne x_2$ veya $y_1 \ne y_2$. Diyelim
$x_1 \ne x_2$; $X$'te ayrık $U_1 \ni x_1$, $U_2 \ni x_2$ seçilir. O zaman
$U_1 \times Y$ ve $U_2 \times Y$ çarpımda açık, ayrık ve iki noktayı ayırır.
\end{proof}
\end{theorem}

\begin{theorem}[Evrensel Özellik — Çarpım]
$f: Z \to X \times Y$ sürekli $\Leftrightarrow \pi_X \circ f$ ve $\pi_Y \circ f$ süreklidir.
\begin{proof}[İspat eskizi]
($\Rightarrow$) Projeksiyonlar sürekli, sürekli fonksiyonların bileşkesi
süreklidir. ($\Leftarrow$) Bir baz elemanı $U \times V$'nin ön-görüntüsü
$f^{-1}(U \times V) = (\pi_X \circ f)^{-1}(U) \cap (\pi_Y \circ f)^{-1}(V)$
olup iki açık ön-görüntünün kesişimi olarak $Z$'de açıktır; baz üzerinde
süreklilik tüm topolojide sürekliliği verir.
\end{proof}
\end{theorem}

%-------------------------------------------------------------------
\section{Algoritmalar}

\subsection{finite\_subspace — $O(|\tau| \cdot |A|)$}

\begin{algorithm}
\caption{Subspace$(X, \tau, A)$}
\begin{algorithmic}[1]
\State $\tau_A \leftarrow \{\}$
\For{each $U \in \tau$}
    \State $\tau_A$.add$(U \cap A)$
\EndFor
\State \Return $(A, \tau_A)$
\end{algorithmic}
\end{algorithm}

\subsection{binary\_product — $O(|\tau_X| \cdot |\tau_Y|)$}

\begin{algorithm}
\caption{Product$(X, \tau_X, Y, \tau_Y)$}
\begin{algorithmic}[1]
\State basis $\leftarrow \{U \times V : U \in \tau_X,\, V \in \tau_Y\}$
\State \Return topology\_from\_basis$(X \times Y,$ basis$)$
\end{algorithmic}
\end{algorithm}

%-------------------------------------------------------------------
\section{pytop API}

\begin{lstlisting}[language=Python]
from pytop import (
    finite_subspace,
    binary_product,
    sierpinski_space,
    discrete_topology,
    indiscrete_topology,
    make_topology,
    make_set,
    empty_set,
    is_compact,
    is_connected,
    is_t2,
    separation_inherited_by_subspace,
)
\end{lstlisting}

\texttt{finite\_subspace(space, subset)} $\to$ \texttt{FiniteTopologicalSpace}

\texttt{binary\_product(left, right)} $\to$ \texttt{FiniteTopologicalSpace} (carrier: çift-demetler)

\texttt{separation\_inherited\_by\_subspace(subspace\_obj, feature="hausdorff")} $\to$ \texttt{Result}

%-------------------------------------------------------------------
\section{Örnekler}

\subsection*{Örnek 5.1 — Alt Uzay: Ayrık Topoloji}

\begin{lstlisting}[language=Python]
d4 = discrete_topology(1, 2, 3, 4)
sub12 = finite_subspace(d4, [1, 2])
print("carrier:", sub12.carrier)
print("topology:", sorted([sorted(list(u)) for u in sub12.topology],
                          key=lambda x: (len(x), x)))
\end{lstlisting}

\begin{verbatim}
carrier: (1, 2)
topology: [[], [1], [2], [1, 2]]
\end{verbatim}

\subsection*{Örnek 5.4 — Çarpım: $D(0,1) \times D(0,1)$}

\begin{lstlisting}[language=Python]
d2 = discrete_topology(0, 1)
prod_dd = binary_product(d2, d2)
print("carrier:", sorted(prod_dd.carrier))
print("topology size:", len(list(prod_dd.topology)))
print("compact:", is_compact(prod_dd).status)
print("connected:", is_connected(prod_dd).status)
\end{lstlisting}

\begin{verbatim}
carrier: [(0, 0), (0, 1), (1, 0), (1, 1)]
topology size: 16
compact: true
connected: false
\end{verbatim}

\subsection*{Örnek 5.5 — Çarpım: Sierpiński $\times$ Sierpiński}

\begin{lstlisting}[language=Python]
ss = binary_product(s, s)
print("topology size:", len(list(ss.topology)))
print("compact:", is_compact(ss).status)
print("connected:", is_connected(ss).status)
\end{lstlisting}

\begin{verbatim}
topology size: 6
compact: true
connected: true
\end{verbatim}

\subsection*{Örnek 5.7 — İndiscrete'in Alt Uzayı İndiscrete'tir}

\begin{lstlisting}[language=Python]
ind4 = indiscrete_topology(1, 2, 3, 4)
sub_ind = finite_subspace(ind4, [1, 2])
print("carrier:", sub_ind.carrier)
print("topology:", sorted([sorted(list(u)) for u in sub_ind.topology],
                          key=lambda x: (len(x), x)))
print("connected:", is_connected(sub_ind).status)
\end{lstlisting}

\begin{verbatim}
carrier: (1, 2)
topology: [[], [1, 2]]
connected: true
\end{verbatim}

\subsection*{Örnek 5.8 — Hausdorff Kalıtımı ve Çarpıma Etkisi}

\begin{lstlisting}[language=Python]
d4 = discrete_topology(1, 2, 3, 4)
sub_d = finite_subspace(d4, [1, 2])
inh = separation_inherited_by_subspace(sub_d, "hausdorff")
print("subspace Hausdorff inherited:", inh.status)

prod_T2 = binary_product(d2, d2)
print("D(0,1) x D(0,1) T2:", is_t2(prod_T2).status)
prod_notT2 = binary_product(d2, indiscrete_topology(0, 1))
print("D(0,1) x Ind(0,1) T2:", is_t2(prod_notT2).status)
\end{lstlisting}

\begin{verbatim}
subspace Hausdorff inherited: true
D(0,1) x D(0,1) T2: true
D(0,1) x Ind(0,1) T2: false
\end{verbatim}

%-------------------------------------------------------------------
\section{Alıştırmalar}

\subsection*{Kodlama}

\begin{enumerate}[label=K\arabic*.]
\item Sierpiński uzayının $\{0\}$ ve $\{1\}$ alt uzaylarını oluşturun ve
      hangi topolojiyi taşıdığını belirleyin.
\item 4-noktalı özel bir uzayda $\{2,3\}$ alt uzayını bulun.
\item \texttt{binary\_product(sierpinski\_space(), discrete\_topology(0,1))}
      için \texttt{is\_compact} ve \texttt{is\_connected} kontrol edin.
\item \texttt{indiscrete\_topology(1,2,3,4)} uzayının \texttt{[1,2]} alt uzayını
      oluşturup açık küme sayısı ile \texttt{is\_connected} çıktısını bulun;
      Örnek~5.7 ile karşılaştırın.
\item \texttt{separation\_inherited\_by\_subspace} ile
      \texttt{discrete\_topology(1,2,3)}'ten alınan \texttt{[1,2]} alt uzayının
      Hausdorff özelliğini miras alıp almadığını kontrol edin.
\end{enumerate}

\subsection*{Teori}

\begin{enumerate}[label=T\arabic*.]
\item Alt uzay topolojisinde dahil etme $i: A \to X$'in sürekli olduğunu ispatlayın.
\item $X$ ve $Y$ bağlantılı $\Rightarrow X \times Y$ bağlantılı olduğunu ispatlayın.
\item Çarpımın evrensel özelliğini ispatlayın: $f: Z \to X \times Y$ sürekli
      $\Leftrightarrow \pi_X \circ f$ ve $\pi_Y \circ f$ süreklidir.
\end{enumerate}

\chapter{Bölüm Topolojisi}

Bölüm topolojisi, bir topolojik uzayın noktalarını denklik sınıflarına göre
\emph{yapıştırarak} yeni bir uzay kurar.
Alt uzay ve çarpım topolojisiyle birlikte temel üç inşaat yönteminden birini oluşturur.

%-------------------------------------------------------------------
\section{Konu}

\begin{definition}[Bölüm Kümesi]
$(X, \tau)$ topolojik uzayı ve $X$ üzerinde bir denklik bağıntısı $\sim$ verilsin.
Denklik sınıfı $[x] = \{y \in X : y \sim x\}$.
\textbf{Bölüm kümesi:} $X/{\sim} = \{[x] : x \in X\}$.
\end{definition}

\begin{definition}[Bölüm Topolojisi]
Bölüm haritası $q: X \to X/{\sim}$,\ $q(x) = [x]$.
\textbf{Bölüm topolojisi} $\tau_{X/\sim}$ üzerinde:
\[
U \subseteq X/{\sim} \text{ açıktır} \iff q^{-1}(U) \in \tau.
\]
Bu, $q$'yu sürekli kılan en ince topolojidir.
\end{definition}

\begin{sezgi}
Bölüm topolojisini bir ``katlama'' gibi düşünün. Elinizde bir kâğıt şerit
($[0,1]$) varsa ve iki ucunu birbirine yapıştırırsanız bir çember ($S^1$) elde
edersiniz. Yapıştırma kuralı bir denklik bağıntısıdır ($0 \sim 1$); bölüm
topolojisi ise ``kâğıdı yırtmadan'' yapılan bu katlamanın doğal topolojisidir.
$q$ haritası her noktayı yapıştırma sonrası bulunduğu konuma gönderir; bir
kümenin yapıştırılmış uzayda açık olması, ön-görüntüsünün orijinal kâğıtta açık
olmasına bağlıdır.
\end{sezgi}

\begin{definition}[Doymuş Küme]
Bir $A \subseteq X$ kümesi \textbf{doymuştur} (saturated), eğer her $x \in A$
için $x$'in tüm denklik sınıfı $[x]$ yine $A$ içinde kalıyorsa, yani
$A = q^{-1}(q(A))$ ise. Bölüm topolojisinin açık kümeleri tam olarak $X$'in
doymuş açık kümelerine karşılık gelir:
\[
q(A)\ X/{\sim}\ \text{içinde açık} \iff A\ \text{doymuş ve}\ X\ \text{içinde açık.}
\]
\end{definition}

\begin{figure}[ht]
\centering
\begin{tikzpicture}[font=\small, etiket/.style={font=\scriptsize}]
  % SOL: [0,1] serit -- uclar 0 ve 1 isaretli
  \draw[blue!70!black, very thick] (0,0) -- (4,0);
  \fill[red!75!black] (0,0) circle (2.2pt);
  \fill[red!75!black] (4,0) circle (2.2pt);
  \node[red!70!black, etiket, below=3pt] at (0,0) {$0$};
  \node[red!70!black, etiket, below=3pt] at (4,0) {$1$};
  \node[etiket, above=4pt] at (2,0) {$[0,1]$ şerit};

  % yapistirma oku
  \draw[-{Stealth[length=2.6mm]}, orange!85!black, very thick]
    (4.6,0) -- (6.2,0);
  \node[orange!80!black, etiket, align=center, above=2pt] at (5.4,0.1)
    {$0 \sim 1$\\yapıştır};

  % SAG: cember S^1, yapistirma noktasi tek
  \draw[blue!70!black, very thick] (8.4,0) circle (1.1);
  \fill[red!75!black] (8.4,1.1) circle (2.4pt);
  \node[red!70!black, etiket, above=3pt] at (8.4,1.1) {$[0]=[1]$};
  \node[etiket, below=4pt] at (8.4,-1.15) {$S^1 = [0,1]/{\sim}$};

  % q haritasi etiketi (cemberin altindaki etikete carpmasin diye sola hizali)
  \node[etiket, align=center, fill=blue!5, rounded corners, inner sep=4pt,
    text width=6.2cm, anchor=north west] at (-0.2,-2.3)
    {Bölüm haritası $q$ her noktayı sınıfına gönderir; \emph{tek} fark, iki ucun
     tek bir noktada (\,$[0]=[1]$\,) çakışmasıdır.};
\end{tikzpicture}

\caption{$[0,1]$ şeridinin uçları yapıştırılarak çember $S^1$ elde edilir;
$q$ haritası her noktayı denklik sınıfına gönderir.}
\end{figure}

\begin{table}[h]
\centering
\begin{tabular}{ll}
\toprule
\textbf{Özellik} & \textbf{Bölüm uzayı için durum} \\
\midrule
Kompaktlık & $X$ kompakt $\Rightarrow X/{\sim}$ kompakt \\
Bağlantılılık & $X$ bağlantılı $\Rightarrow X/{\sim}$ bağlantılı \\
Hausdorff & Genel olarak korunmaz \\
T1 & $\sim$ kapalı bağıntı ise korunur \\
\bottomrule
\end{tabular}
\caption{Topolojik özelliklerin bölüm altında korunumu}
\end{table}

\begin{figure}[ht]
\centering
\begin{tikzpicture}[font=\small, etiket/.style={font=\scriptsize},
  nokta/.style={circle, fill=blue!70!black, inner sep=1.8pt}]
  % SOL: X = {a,b,c,d}, sinif kutulariyla gruplanmis
  \node[etiket, above] at (1.3,2.2) {$X = \{a,b,c,d\}$};

  % {a,b} kutusu
  \draw[violet!60!black, thick, rounded corners, fill=violet!7]
    (-0.1,1.0) rectangle (1.2,2.0);
  \node[nokta, label={[etiket]left:$a$}] (a) at (0.25,1.5) {};
  \node[nokta, label={[etiket]right:$b$}] (b) at (0.85,1.5) {};
  \node[violet!60!black, etiket, below=1pt] at (0.55,1.0) {$\{a,b\}$};

  % {c}
  \draw[teal!55!black, thick, rounded corners, fill=teal!7]
    (1.7,1.15) rectangle (2.5,1.85);
  \node[nokta, label={[etiket]above:$c$}] (c) at (2.1,1.5) {};

  % {d}
  \draw[orange!75!black, thick, rounded corners, fill=orange!7]
    (3.0,1.15) rectangle (3.8,1.85);
  \node[nokta, label={[etiket]above:$d$}] (d) at (3.4,1.5) {};

  % q oklari asagi
  \draw[-{Stealth[length=2.2mm]}, gray!70] (0.55,0.95) -- (0.55,-0.05);
  \draw[-{Stealth[length=2.2mm]}, gray!70] (2.1,1.1) -- (2.1,-0.05);
  \draw[-{Stealth[length=2.2mm]}, gray!70] (3.4,1.1) -- (3.4,-0.05);
  \node[etiket, right] at (3.55,0.55) {$q$};

  % SAG/ALT: X/~ = uc nokta
  \node[etiket, below] at (1.95,-0.9) {$X/{\sim} = \{[a],[c],[d]\}$ (üç nokta)};
  \node[circle, draw=violet!60!black, fill=violet!18, inner sep=2.6pt,
    label={[etiket]below:$[a]$}] at (0.55,-0.45) {};
  \node[circle, draw=teal!55!black, fill=teal!18, inner sep=2.6pt,
    label={[etiket]below:$[c]$}] at (2.1,-0.45) {};
  \node[circle, draw=orange!75!black, fill=orange!18, inner sep=2.6pt,
    label={[etiket]below:$[d]$}] at (3.4,-0.45) {};
\end{tikzpicture}

\caption{Denklik sınıfları bölüm uzayında birer noktaya çöker: dört nokta
$\{a,b\}, \{c\}, \{d\}$ sınıflarıyla üç noktaya iner.}
\end{figure}

\begin{karsiornek}
Hausdorff özelliği bölüm altında korunmaz. $\mathbb{R}$ üzerinde
``$x \sim y \iff x - y \in \mathbb{Q}$'' bağıntısını alın. Bölüm uzayı
$\mathbb{R}/\mathbb{Q}$ kaba (indiscrete) bir uzaydır: tek açık kümeler
$\emptyset$ ve tüm uzaydır, çünkü $\mathbb{Q}$ ile öteleme altında doymuş tek
açık küme boş küme ve $\mathbb{R}$'dir. Dolayısıyla iki farklı nokta hiçbir
zaman ayrık komşuluklarla ayrılamaz --- $\mathbb{R}$ Hausdorff olmasına rağmen
$\mathbb{R}/\mathbb{Q}$ T1 bile değildir. Sonlu modelde de aynı tuzak geçerlidir:
doymamış açık kümeler yapıştırma sonrası ``kaybolur'' ve ayırma aksiyomlarını
bozabilir.
\end{karsiornek}

%-------------------------------------------------------------------
\section{Teoremler}

\begin{theorem}[Bölüm haritası sürekliliği]
$q: X \to X/{\sim}$ her zaman süreklidir.
\end{theorem}

\begin{theorem}[Evrensel Özellik]
$g: X/{\sim} \to Z$ olsun.
$g$ süreklidir $\Longleftrightarrow g \circ q: X \to Z$ süreklidir.
\begin{proof}[İspat eskizi]
($\Rightarrow$) $g$ sürekli ve $q$ sürekli (Teorem~1) olduğundan, iki sürekli
haritanın bileşkesi $g \circ q$ süreklidir.
($\Leftarrow$) $g \circ q$'nun sürekli olduğunu varsayalım. $Z$'de bir $V$ açık
kümesi alın; $g^{-1}(V) \subseteq X/{\sim}$ kümesinin bölüm topolojisinde açık
olduğunu göstermeliyiz. Tanım gereği bu, $q^{-1}(g^{-1}(V))$ kümesinin $X$'te
açık olmasına denktir. Fakat
$q^{-1}(g^{-1}(V)) = (g \circ q)^{-1}(V)$ ve $g \circ q$ sürekli olduğundan bu
küme $X$'te açıktır. Dolayısıyla $g^{-1}(V)$ açıktır ve $g$ süreklidir. Bu
özellik, bölüm uzayından çıkan sürekli haritaların kontrolünü $X$ üzerindeki
haritaların kontrolüne indirger ve bölüm topolojisini evrensel bir nesne olarak
karakterize eder.
\end{proof}
\end{theorem}

\begin{theorem}[Kompaktlık Korunumu]
$X$ kompakt $\Rightarrow X/{\sim}$ kompakttır.
\begin{proof}[İspat eskizi]
$q$ sürekli ve örtendir; kompakt bir kümenin sürekli görüntüsü kompakt olduğundan
$X/{\sim} = q(X)$ kompakttır.
\end{proof}
\end{theorem}

\begin{theorem}[Bağlantılılık Korunumu]
$X$ bağlantılı $\Rightarrow X/{\sim}$ bağlantılıdır.
\end{theorem}

\begin{theorem}[Hausdorff Kriteri]
$X/{\sim}$ Hausdorff $\Longleftrightarrow$ $\sim$ grafiği $X \times X$'te kapalıdır.
\end{theorem}

%-------------------------------------------------------------------
\section{Algoritmalar}

\subsection{quotient\_set — $O(|X|^2 \cdot \alpha(|X|))$}

\begin{algorithm}
\caption{QuotientSet$(X, \sim)$}
\begin{algorithmic}[1]
\State Her $x$ için sınıf $\leftarrow \{x\}$ (union-find başlat)
\For{each $(x, y) \in {\sim}$}
    \State \textbf{Union}$(x, y)$
\EndFor
\State \Return kök temsilcilere göre sınıfları tuple olarak
\end{algorithmic}
\end{algorithm}

\subsection{finite\_quotient\_contract — $O(|X| + |P|)$}

\begin{algorithm}
\caption{FiniteQuotientContract$(X, P)$}
\begin{algorithmic}[1]
\State \textit{// $P$: $X$'in örtüşmeyen bölüntüsü}
\State Doğrula: $P$ her $x$'i kapsar ve bloklar ayrık
\State \Return FiniteConstructionContract(status=\texttt{true}, carrier\_size=$|X|$, block\_count=$|P|$)
\end{algorithmic}
\end{algorithm}

%-------------------------------------------------------------------
\section{pytop API}

\begin{lstlisting}[language=Python]
from pytop import (
    quotient_set,               # denklik bagintisindan bolum kumesi
    finite_quotient_contract,   # boluntuden insaat sozlesmesi
    finite_quotient_summary,    # kisa ozet dizesi
    make_quotient_map,          # QuotientMap nesnesi
    is_quotient_map,            # Result: harita bolum haritasi mi?
    equivalence_class,          # bir noktanin denklik sinifi
    partition_from_equivalence, # denklik bagintisi -> boluntu
    is_equivalence_relation,    # baginti denklik bagintisi mi?
    equivalence_from_partition, # boluntu -> denklik bagintisi
    equivalence_from_classes,   # sinif bloklarindan denklik bagintisi
    canonical_projection_from_equivalence,  # q haritasini sozluk olarak
    discrete_topology,          # ayrik topoloji
    sierpinski_space,           # Sierpinski uzayi
    make_topology,              # tasiyici + aciklardan FiniteTopologicalSpace
)
\end{lstlisting}

\texttt{quotient\_set(carrier, relation)} $\to$ \texttt{tuple[frozenset, ...]}

\texttt{finite\_quotient\_contract(carrier, partition)} $\to$ \texttt{FiniteConstructionContract}

\texttt{finite\_quotient\_summary(carrier, partition)} $\to$ \texttt{str}

\texttt{make\_quotient\_map(domain, codomain)} $\to$ \texttt{QuotientMap}

\texttt{is\_quotient\_map(map\_obj)} $\to$ \texttt{Result}

\texttt{equivalence\_from\_partition(universe, partition)} $\to$ \texttt{set[tuple]}

\texttt{partition\_from\_equivalence(carrier, relation)} $\to$ \texttt{set[frozenset]}

\texttt{equivalence\_class(carrier, relation, point)} $\to$ \texttt{set}

\texttt{canonical\_projection\_from\_equivalence(carrier, relation)} $\to$ \texttt{dict}

%-------------------------------------------------------------------
\section{Örnekler}

\subsection*{Örnek 5.1 — İki Noktayı Özdeşleştirme}

\begin{lstlisting}[language=Python]
d4 = discrete_topology(0, 1, 2, 3)
partition_1 = [[0, 1], [2], [3]]
qs1 = quotient_set(
    [0, 1, 2, 3],
    [(0,0),(1,1),(2,2),(3,3),(0,1),(1,0)]
)
print("Bolum kumesi:", qs1)
contract1 = finite_quotient_contract([0, 1, 2, 3], partition_1)
print("Blok sayisi:", contract1.block_count)
print("Durum:", contract1.status)
\end{lstlisting}

\begin{verbatim}
Bolum kumesi: (frozenset({2}), frozenset({3}), frozenset({0, 1}))
Blok sayisi: 3
Durum: true
\end{verbatim}

\subsection*{Örnek 5.2 — Tüm Noktaları Tek Noktaya Çökertme}

\begin{lstlisting}[language=Python]
carrier = [0, 1, 2, 3]
qs2 = quotient_set(
    carrier,
    [(i, j) for i in carrier for j in carrier]
)
print("Tek blok:", qs2)
contract2 = finite_quotient_contract(carrier, [carrier])
print("Blok sayisi:", contract2.block_count)
r2 = contract2.to_result()
print("Metadata:", r2.metadata['block_sizes'])
\end{lstlisting}

\begin{verbatim}
Tek blok: (frozenset({0, 1, 2, 3}),)
Blok sayisi: 1
Metadata: [4]
\end{verbatim}

\subsection*{Örnek 5.3 — Özel Topolojide Bölüm Sözleşmesi}

\begin{lstlisting}[language=Python]
X5 = [1, 2, 3, 4, 5]
tau5 = [set(), {1,2}, {3,4}, {1,2,3,4}, {1,2,3,4,5}]
fts5 = make_topology(X5, *tau5)
partition5 = [[1, 2], [3, 4], [5]]
contract5 = finite_quotient_contract(X5, partition5)
r5 = contract5.to_result()
print("Status:", r5.status)
print("Bloklar:", r5.metadata['block_count'])
print("Blok boyutlari:", r5.metadata['block_sizes'])
\end{lstlisting}

\begin{verbatim}
Status: true
Bloklar: 3
Blok boyutlari: [2, 2, 1]
\end{verbatim}

\subsection*{Örnek 5.4 — Bölüm Haritası Nesnesi}

\begin{lstlisting}[language=Python]
d3 = discrete_topology(0, 1, 2)
d2 = discrete_topology('a', 'b')
qmap = make_quotient_map(d3, d2)
print("Etiketler:", sorted(qmap.tags))
r_qmap = is_quotient_map(qmap)
print("is_quotient_map status:", r_qmap.status)
print("Justification:", r_qmap.justification[0])
\end{lstlisting}

\begin{verbatim}
Etiketler: ['continuous', 'quotient', 'surjective']
is_quotient_map status: true
Justification: Map tag 'quotient' is explicitly present.
\end{verbatim}

\subsection*{Örnek 5.5 — [0,1] Uçlarını Yapıştırarak Çember $S^1$}

$[0,1]$ aralığını $5$ noktayla modelliyoruz. İki ucu ($0$ ve $4$)
özdeşleştirmek, şeridi bir çembere yapıştırmaya karşılık gelir.

\begin{lstlisting}[language=Python]
ring = [0, 1, 2, 3, 4]
ring_classes = [[0, 4], [1], [2], [3]]
ring_rel = equivalence_from_partition(ring, ring_classes)
print("Denklik mi?", is_equivalence_relation(ring, ring_rel))
print("Bolum kumesi boyutu:", len(quotient_set(ring, ring_rel)))
print("Ozet:", finite_quotient_summary(ring, ring_classes))
print("0'in sinifi:", sorted(equivalence_class(ring, ring_rel, 0)))
print("4'un sinifi:", sorted(equivalence_class(ring, ring_rel, 4)))
\end{lstlisting}

\begin{verbatim}
Denklik mi? True
Bolum kumesi boyutu: 4
Ozet: quotient: status=true, carrier=5, blocks=4
0'in sinifi: [0, 4]
4'un sinifi: [0, 4]
\end{verbatim}

Uçlar tek bir sınıfta birleşti: $5$ noktalı şerit, $4$ noktalı çembere indi.

\subsection*{Örnek 5.6 — Doymuş Kümeler ve Kanonik İzdüşüm}

Kanonik izdüşüm $q$ her noktayı denklik sınıfına gönderir. Bir kümenin
\textbf{doymuş} olması, içerdiği her noktanın sınıfını da tamamen içermesi
demektir.

\begin{lstlisting}[language=Python]
sat_carrier = ['a', 'b', 'c', 'd']
sat_part = [['a', 'b'], ['c'], ['d']]
sat_rel = equivalence_from_partition(sat_carrier, sat_part)
proj = canonical_projection_from_equivalence(sat_carrier, sat_rel)
for pt in sat_carrier:
    print(f"q({pt}) =", sorted(proj[pt]))
cls_a = equivalence_class(sat_carrier, sat_rel, 'a')
A = {'a', 'b'}
print("{a,b} doymus mu?",
      all(equivalence_class(sat_carrier, sat_rel, p) <= A for p in A))
print("{a} doymus mu?", cls_a <= {'a'})
\end{lstlisting}

\begin{verbatim}
q(a) = ['a', 'b']
q(b) = ['a', 'b']
q(c) = ['c']
q(d) = ['d']
{a,b} doymus mu? True
{a} doymus mu? False
\end{verbatim}

$\{a, b\}$ doymuştur, ama $\{a\}$ doymuş değildir: sınıf arkadaşı $b$'yi
içermez. Bölüm topolojisinin açık kümeleri ancak doymuş açık kümelerden gelir.

\subsection*{Örnek 5.7 — Sınıflardan Noktalara: Round-Trip}

\texttt{equivalence\_from\_classes} sınıf bloklarından denklik bağıntısı kurar;
\texttt{partition\_from\_equivalence} bağıntıyı tekrar bölüntüye çevirir.

\begin{lstlisting}[language=Python]
rt_blocks = equivalence_from_classes([1, 2], [3], [4, 5])
print("Denklik mi?", is_equivalence_relation([1, 2, 3, 4, 5], rt_blocks))
rt_part = partition_from_equivalence([1, 2, 3, 4, 5], rt_blocks)
print("Blok sayisi:", len(rt_part))
print("Bloklar:", sorted(sorted(b) for b in rt_part))
\end{lstlisting}

\begin{verbatim}
Denklik mi? True
Blok sayisi: 3
Bloklar: [[1, 2], [3], [4, 5]]
\end{verbatim}

Sınıflardan kurulan bağıntı, bölüntüye çevrilip geri okunduğunda aynı üç bloğu
verir: bölüm uzayının $3$ noktası, tam olarak $3$ denklik sınıfına karşılık gelir.

%-------------------------------------------------------------------
\section{Alıştırmalar}

\subsection*{Kodlama}

\begin{enumerate}[label=K\arabic*.]
\item \texttt{discrete\_topology(0,1,2,3,4)} üzerinde \texttt{[[0,4],[1,3],[2]]}
      bölüntüsünü kullanarak \texttt{finite\_quotient\_contract} çağırın;
      blok sayısını ve boyutlarını yazdırın.
\item \texttt{sierpinski\_space()} üzerinde trivial bölüntü ile
      \texttt{finite\_quotient\_summary} çağırın ve çıktıyı yorumlayın.
\item \texttt{make\_topology([1,2,3,4], set(), \{1,2\}, \{3,4\}, \{1,2,3,4\})} üzerinde
      \texttt{[[1,2],[3,4]]} ve \texttt{[[1],[2],[3],[4]]} bölüntülerini karşılaştırın.
\item \texttt{[0,1,2,3,4,5]} taşıyıcısı üzerinde \texttt{[[0,5],[1],[2],[3],[4]]}
      bölüntüsünü \texttt{equivalence\_from\_partition} ile denklik bağıntısına
      çevirin; \texttt{is\_equivalence\_relation} ile doğrulayın ve $0$ ile $5$'in
      \texttt{equivalence\_class} çıktılarının aynı olduğunu gösterin.
\item \texttt{['x','y','z']} taşıyıcısı ve \texttt{[['x','y'],['z']]} bölüntüsü için
      \texttt{canonical\_projection\_from\_equivalence} çağırıp her noktanın
      görüntüsünü yazdırın; \texttt{\{'x','y'\}} kümesinin doymuş,
      \texttt{\{'x'\}} kümesinin doymuş olmadığını gösterin.
\end{enumerate}

\subsection*{Teori}

\begin{enumerate}[label=T\arabic*.]
\item $q: X \to X/{\sim}$ bölüm haritasının $\tau_{X/\sim}$ tanımı gereği
      her zaman sürekli olduğunu ispatlayın.
\item $X$ bağlantılı $\Rightarrow X/{\sim}$ bağlantılı olduğunu ispatlayın.
      ($q$ sürekli ve surjektif; bağlantılı kümenin sürekli görüntüsü bağlantılıdır.)
\item Bölüm haritasının \textbf{evrensel özelliğini} ispatlayın:
      $g: X/{\sim} \to Z$ için $g$ sürekli $\iff g \circ q$ sürekli.
      \ipucu{$(g \circ q)^{-1}(V) = q^{-1}(g^{-1}(V))$ eşitliğini bölüm
      topolojisinin tanımıyla birleştirin.}
\item $A \subseteq X$ doymuş ($A = q^{-1}(q(A))$) ise $q(A)$'nın bölüm uzayında
      açık olması için $A$'nın $X$'te açık olmasının yeterli ve gerekli olduğunu
      gösterin. Doymamış bir $A$ için bu yazışmanın neden bozulduğunu açıklayın.
\end{enumerate}

\chapter{Başlangıç ve Son Topoloji}
\label{ch:initial_final}

\textbf{Başlangıç topolojisi} (initial topology), verilen haritaları sürekli kılan en
kaba (coarsest) topolojidir. \textbf{Son topoloji} (final topology), verilen haritaları
sürekli kılan en ince (finest) topolojidir. Altuzay ve çarpım topolojileri başlangıç
topolojisinin; bölüm topolojisi ise son topolojisinin özel halleridir.

\begin{sezgi}
Başlangıç topolojisini, $X$'e ``tam yeterince'' açık küme koyan bir ayar düğmesi
gibi düşünün. Çok az açık koyarsanız $f_\alpha$ haritaları sürekli olmaz; çok fazla
koyarsanız gereksizdir. Geri çekme $f_\alpha^{-1}(U)$ ile üretilen alt-baz,
haritaları sürekli kılan \emph{en küçük} (en kaba) açık-küme koleksiyonudur. Son
topoloji ise tam tersi yönde çalışır: $Y$'ye en çok açık kümeyi koyar, ama her
$g_i$'yi sürekli bozmadan.
\end{sezgi}

\begin{figure}[ht]
\centering
\begin{tikzpicture}[font=\small]
  % Domain X (left)
  \draw[rounded corners=6pt, thick] (-0.6,-1.5) rectangle (1.9,1.6);
  \node at (0.65,1.85) {$X$ (taşıyıcı, topolojisiz)};
  \draw[blue!70!black, thick, fill=blue!10]
    (0.65,0.55) ellipse [x radius=0.78, y radius=0.42];
  \node[blue!70!black, font=\scriptsize] at (0.65,0.55) {$f_1^{-1}(U_1)$};
  \draw[red!70!black, thick, fill=red!10]
    (0.65,-0.7) ellipse [x radius=0.78, y radius=0.42];
  \node[red!70!black, font=\scriptsize] at (0.65,-0.7) {$f_2^{-1}(U_2)$};

  % Codomain Y_1 (top right)
  \draw[rounded corners=6pt, thick] (5.4,0.5) rectangle (7.6,1.9);
  \node at (6.5,2.15) {$(Y_1,\tau_1)$};
  \draw[blue!70!black, thick, fill=blue!12]
    (6.5,1.2) ellipse [x radius=0.55, y radius=0.32];
  \node[blue!70!black, font=\scriptsize] at (6.5,1.2) {$U_1$};

  % Codomain Y_2 (bottom right)
  \draw[rounded corners=6pt, thick] (5.4,-1.9) rectangle (7.6,-0.5);
  \node at (6.5,-2.2) {$(Y_2,\tau_2)$};
  \draw[red!70!black, thick, fill=red!12]
    (6.5,-1.2) ellipse [x radius=0.55, y radius=0.32];
  \node[red!70!black, font=\scriptsize] at (6.5,-1.2) {$U_2$};

  % Maps
  \draw[blue!70!black, -{Stealth[length=2.4mm]}, thick]
    (1.5,0.55) to[bend left=12] node[above, font=\scriptsize, pos=0.55] {$f_1$} (5.95,1.2);
  \draw[red!70!black, -{Stealth[length=2.4mm]}, thick]
    (1.5,-0.7) to[bend right=12] node[below, font=\scriptsize, pos=0.55] {$f_2$} (5.95,-1.2);

  \node[font=\scriptsize, align=center] at (3.0,-3.0)
    {Alt-baz $\mathcal{S}=\{f_\alpha^{-1}(U)\}$ geri çekme ile üretilir;\\
     $\tau_{\mathrm{ini}}=\langle\mathcal{S}\rangle$ her $f_\alpha$'yı sürekli kılan \emph{en kaba} topoloji.};
\end{tikzpicture}

\caption{Başlangıç topolojisi: bir harita ailesini sürekli kılan en kaba topoloji,
geri çekme $f_\alpha^{-1}(U)$ ile üretilen alt-bazdan doğar.}
\label{fig:ch13_baslangic}
\end{figure}

% -----------------------------------------------------------------------
\section{Başlangıç Topolojisi}
\label{sec:initial}

$X$ kümesi ve haritalar ailesi $\{f_\alpha : X \to (Y_\alpha, \tau_\alpha)\}_{\alpha \in A}$
verilsin.

\begin{definition}[Başlangıç Topolojisi]
\label{def:initial_topology}
\textbf{Başlangıç topolojisi} $\tau_{\mathrm{ini}}$, her $f_\alpha$'yı $(X, \tau_{\mathrm{ini}}) \to (Y_\alpha, \tau_\alpha)$
sürekli kılan topolojilerin kesişimidir. Alt-baz:
\[
\mathcal{S} = \{f_\alpha^{-1}(U) \mid U \in \tau_\alpha,\ \alpha \in A\},
\quad
\tau_{\mathrm{ini}} = \langle \mathcal{S} \rangle.
\]
\end{definition}

\begin{theorem}[Varlık ve Teklik]
\label{thm:initial_exists}
$\tau_{\mathrm{ini}}$ tekil ve vardır; her $f_\alpha$'yı sürekli kılan tüm topolojilerin
en kaba olanıdır.
\end{theorem}

\begin{proof}[Kanıt İskeleti]
$\mathcal{S}$'nin ürettiği topoloji $\langle\mathcal{S}\rangle$ her $f_\alpha$'yı
sürekli kılar (çünkü $f_\alpha^{-1}(U) \in \langle\mathcal{S}\rangle$ her $U \in \tau_\alpha$
için). Eğer $\tau'$ de her $f_\alpha$'yı sürekli kılıyorsa, $\mathcal{S} \subseteq \tau'$
olduğundan $\langle\mathcal{S}\rangle \subseteq \tau'$; yani $\tau_{\mathrm{ini}}$ en kaba.
\end{proof}

\begin{theorem}[Evrensel Özellik]
\label{thm:initial_universal}
$h: Z \to (X, \tau_{\mathrm{ini}})$ süreklidir ancak ve ancak $f_\alpha \circ h$ her
$\alpha$ için sürekli olduğunda.
\end{theorem}

\begin{proof}[İspat eskizi]
($\Rightarrow$) $h$ sürekli ve $f_\alpha$ sürekli ise bileşke $f_\alpha \circ h$
süreklidir. ($\Leftarrow$) $\tau_{\mathrm{ini}}$'nin bir alt-baz elemanı
$f_\alpha^{-1}(U)$ biçimindedir ve ön-görüntüsü
$h^{-1}(f_\alpha^{-1}(U)) = (f_\alpha \circ h)^{-1}(U)$, ki bu $f_\alpha \circ h$
sürekli olduğundan $Z$'de açıktır. Süreklilik alt-bazda denetlenebildiğinden $h$
süreklidir.
\end{proof}

\begin{figure}[ht]
\centering
\begin{tikzpicture}[font=\small, node distance=14mm,
    uz/.style={draw, rounded corners, fill=blue!6, inner sep=6pt, minimum width=10mm}]
  \node[uz] (Z) at (0,0) {$Z$};
  \node[uz] (X) at (3.4,0) {$(X,\tau_{\mathrm{ini}})$};
  \node[uz] (Y) at (3.4,-2.4) {$(Y_\alpha,\tau_\alpha)$};

  \draw[-{Stealth[length=2.6mm]}, thick, dashed]
    (Z) -- node[above, font=\scriptsize] {$h$ (tek, sürekli)} (X);
  \draw[-{Stealth[length=2.6mm]}, thick]
    (X) -- node[right, font=\scriptsize] {$f_\alpha$} (Y);
  \draw[-{Stealth[length=2.6mm]}, thick]
    (Z) -- node[below left, font=\scriptsize, pos=0.55] {$f_\alpha\circ h$} (Y);

  \node[font=\scriptsize, align=center] at (1.7,-3.5)
    {Evrensel özellik: $h$ sürekli $\iff$ her $f_\alpha\circ h$ sürekli.\\
     $\tau_{\mathrm{ini}}$ bu faktörizasyonu mümkün kılan \emph{tek} topolojidir.};
\end{tikzpicture}

\caption{Evrensel özellik: $h$ sürekli $\iff$ her $f_\alpha\circ h$ sürekli;
$\tau_{\mathrm{ini}}$ bu faktörizasyonu mümkün kılan tek topolojidir.}
\label{fig:ch13_evrensel}
\end{figure}

\begin{theorem}[Özel Haller]
\label{thm:initial_special}
\begin{enumerate}
  \item Altuzay topolojisi, ekleme haritası $i: A \hookrightarrow X$'in başlangıç
        topolojisidir.
  \item Çarpım topolojisi, projeksiyon haritaları $\{\pi_\alpha\}$'nın başlangıç
        topolojisidir.
\end{enumerate}
\end{theorem}

% -----------------------------------------------------------------------
\section{Algoritmalar}
\label{sec:initial_algorithms}

\subsection{Başlangıç Topolojisi İnşası}

\begin{algorithm}
\caption{Başlangıç Topolojisi (Sonlu Durum)}
\label{alg:initial_topology}
\begin{algorithmic}[1]
\Require Sonlu taşıyıcı $X$, haritalar $\{(f_\alpha, \tau_\alpha)\}$
\Ensure $\tau_{\mathrm{ini}}$ — en kaba süreklilik topolojisi
\State $\mathcal{S} \leftarrow \emptyset$
\For{her $f_\alpha$}
  \For{her $U \in \tau_\alpha$}
    \State $\mathcal{S} \leftarrow \mathcal{S} \cup \{f_\alpha^{-1}(U)\}$
    \State \textit{// $f_\alpha^{-1}(U) = \{x \in X \mid f_\alpha(x) \in U\}$}
  \EndFor
\EndFor
\State \Return $\langle \mathcal{S} \rangle$ \textit{// alt-bazdan topoloji üret}
\end{algorithmic}
\end{algorithm}

\textbf{Karmaşıklık.} Ön-görüntü hesabı $O(|\tau_\alpha| \cdot |X|)$; topoloji üretimi
$O(|\mathcal{S}| \cdot 2^{|X|})$ en kötü durumda.

\subsection{Son Topoloji İnşası}

\begin{algorithm}
\caption{Son Topoloji (Sonlu Durum)}
\label{alg:final_topology}
\begin{algorithmic}[1]
\Require Sonlu $Y$, kaynak uzaylar $\{(X_i, \sigma_i)\}$, haritalar $\{g_i : X_i \to Y\}$
\Ensure $\tau_{\mathrm{fin}}$ — en ince süreklilik topolojisi
\State $\tau_{\mathrm{fin}} \leftarrow \emptyset$
\For{her $V \subseteq Y$}
  \If{$g_i^{-1}(V) \in \sigma_i$ her $i$ için}
    \State $\tau_{\mathrm{fin}} \leftarrow \tau_{\mathrm{fin}} \cup \{V\}$
  \EndIf
\EndFor
\State \Return $\tau_{\mathrm{fin}}$
\end{algorithmic}
\end{algorithm}

\textbf{Karmaşıklık.} $O(2^{|Y|} \cdot \sum_i |\sigma_i|)$.

% -----------------------------------------------------------------------
\section{pytop API}
\label{sec:initial_api}

\begin{lstlisting}[language=Python, caption={Başlangıç Topolojisi API}]
from pytop.finite_map_analysis import FiniteMap
from pytop import (
    initial_topology_from_maps,
    generic_quotient_space_from_map,
    discrete_topology,
    indiscrete_topology,
    make_topology,
    sierpinski_space,
    finite_subspace,
    binary_product,
    is_continuous_finite_map,
    is_t0,
    is_t2,
)

# FiniteMap dogrudan veri sinifi olarak kullanilir (make_function degil)
# tau_ini = initial_topology_from_maps(carrier, [f1, f2, ...])
# quot    = generic_quotient_space_from_map(q_map)   # son topoloji, bolum ozel hali
\end{lstlisting}

\textbf{Önemli:} \texttt{pytop.quotient\_space\_from\_map} sembolik motoru kullanır;
sonlu-duyarlı sürüm \texttt{generic\_quotient\_space\_from\_map}'tir.

% -----------------------------------------------------------------------
\section{Örnekler}
\label{sec:initial_examples}

\subsection{Tek Harita ile Başlangıç Topolojisi}

\begin{lstlisting}[language=Python]
X_carrier = [0, 1, 2, 3]
Y_d = discrete_topology(0, 1)
X_ph = make_topology(X_carrier, set(), set(X_carrier))

f = FiniteMap(domain=X_ph, codomain=Y_d, name="f",
              mapping={0: 0, 1: 0, 2: 1, 3: 1})
tau_f = initial_topology_from_maps(X_carrier, [f])

print("Baslangic topolojisi:")
for u in sorted([sorted(u) for u in tau_f.topology], key=lambda x: (len(x), x)):
    print("  ", u if u else "empty")
print("Eleman sayisi:", len(tau_f.topology))
\end{lstlisting}

\begin{lstlisting}[style=output]
Baslangic topolojisi:
   empty
   [0, 1]
   [2, 3]
   [0, 1, 2, 3]
Eleman sayisi: 4
\end{lstlisting}

$f^{-1}(\{0\}) = \{0,1\}$, $f^{-1}(\{1\}) = \{2,3\}$ — alt-baz $\{\{0,1\},\{2,3\}\}$,
üretilen topoloji 4 elemanlı.

\subsection{İki Harita — Topoloji İncelir}

\begin{lstlisting}[language=Python]
S_sier = sierpinski_space()
g = FiniteMap(domain=X_ph, codomain=S_sier, name="g",
              mapping={0: 0, 1: 1, 2: 1, 3: 1})
tau_fg = initial_topology_from_maps(X_carrier, [f, g])

print("f tek basina:", len(tau_f.topology), "eleman  |  f+g:", len(tau_fg.topology), "eleman")
\end{lstlisting}

\begin{lstlisting}[style=output]
f tek basina: 4 eleman  |  f+g: 6 eleman
\end{lstlisting}

$g^{-1}(\{1\}) = \{1,2,3\}$ yeni bir alt-baz elemanı ekler; kesişim $\{0,1\} \cap \{1,2,3\} = \{1\}$
yeni bir açık küme üretir.

\subsection{Altuzay = Başlangıç}

\begin{lstlisting}[language=Python]
fts_X = make_topology([0, 1, 2, 3], set(), {0, 1}, {2, 3}, {0, 1, 2, 3})
sub_A  = finite_subspace(fts_X, [1, 2])

A_dom  = discrete_topology(1, 2)
inc    = FiniteMap(domain=A_dom, codomain=fts_X, name="i", mapping={1: 1, 2: 2})
init_A = initial_topology_from_maps([1, 2], [inc])

sub_tau  = sorted([sorted(u) for u in sub_A.topology],  key=lambda x: (len(x), x))
init_tau = sorted([sorted(u) for u in init_A.topology], key=lambda x: (len(x), x))
print("Esit mi:", sub_tau == init_tau)
\end{lstlisting}

\begin{lstlisting}[style=output]
Esit mi: True
\end{lstlisting}

\subsection{Çarpım = Başlangıç}

\begin{lstlisting}[language=Python]
d2   = discrete_topology(0, 1)
prod = binary_product(d2, d2)

pi_x = FiniteMap(domain=prod, codomain=d2, name="pi_x",
                 mapping={(0,0):0, (0,1):0, (1,0):1, (1,1):1})
pi_y = FiniteMap(domain=prod, codomain=d2, name="pi_y",
                 mapping={(0,0):0, (0,1):1, (1,0):0, (1,1):1})

init_prod = initial_topology_from_maps(list(prod.carrier), [pi_x, pi_y])
prod_fs   = {frozenset(u) for u in prod.topology}
init_fs   = {frozenset(u) for u in init_prod.topology}
print("Carpim == Baslangic:", prod_fs == init_fs)
\end{lstlisting}

\begin{lstlisting}[style=output]
Carpim == Baslangic: True
\end{lstlisting}

\subsection{Son Topoloji — Manuel Hesap}

\begin{lstlisting}[language=Python]
X1     = make_topology([0, 1, 2], set(), {0, 1}, {0, 1, 2})
Y_car  = [0, 1]
g1_map = {0: 0, 1: 0, 2: 1}
tau_X1 = {frozenset(u) for u in X1.topology}

open_Y = []
for mask in range(1 << len(Y_car)):
    V        = frozenset(Y_car[i] for i in range(len(Y_car)) if mask & (1 << i))
    preimage = frozenset(x for x in X1.carrier if g1_map[x] in V)
    if preimage in tau_X1:
        open_Y.append(set(V))

final_Y   = make_topology(Y_car, *open_Y)
final_tau = sorted([sorted(u) for u in final_Y.topology], key=lambda x: (len(x), x))
print("Son topoloji:", final_tau)
\end{lstlisting}

\begin{lstlisting}[style=output]
Son topoloji: [[], [0], [0, 1]]
\end{lstlisting}

\subsection{Bölüm = Son Topoloji}

\begin{lstlisting}[language=Python]
Y_ind    = indiscrete_topology(0, 1)
q        = FiniteMap(domain=X1, codomain=Y_ind, name="q",
                     mapping={0: 0, 1: 0, 2: 1})
quot_Y   = generic_quotient_space_from_map(q)
quot_tau = sorted([sorted(u) for u in quot_Y.topology], key=lambda x: (len(x), x))

print("Bolum topolojisi:", quot_tau)
print("Son topoloji ile esit:", quot_tau == final_tau)
\end{lstlisting}

\begin{lstlisting}[style=output]
Bolum topolojisi: [[], [0], [0, 1]]
Son topoloji ile esit: True
\end{lstlisting}

\begin{figure}[ht]
\centering
\begin{tikzpicture}[font=\small]
  % Source X_1 (top left)
  \draw[rounded corners=6pt, thick] (-0.4,0.5) rectangle (1.8,1.9);
  \node at (0.7,2.15) {$(X_1,\sigma_1)$};
  \draw[blue!70!black, thick, fill=blue!12]
    (0.7,1.2) ellipse [x radius=0.55, y radius=0.32];

  % Source X_2 (bottom left)
  \draw[rounded corners=6pt, thick] (-0.4,-1.9) rectangle (1.8,-0.5);
  \node at (0.7,-2.2) {$(X_2,\sigma_2)$};
  \draw[red!70!black, thick, fill=red!12]
    (0.7,-1.2) ellipse [x radius=0.55, y radius=0.32];

  % Target Y (right)
  \draw[rounded corners=6pt, thick] (5.2,-1.5) rectangle (7.7,1.6);
  \node at (6.45,1.85) {$Y$ (taşıyıcı)};
  \draw[violet!75!black, thick, fill=violet!10]
    (6.45,0.0) ellipse [x radius=0.85, y radius=0.55];
  \node[violet!75!black, font=\scriptsize] at (6.45,0.0) {$V$};
  \node[violet!75!black, font=\scriptsize, align=center] at (6.45,-1.05)
    {$V$ açık $\iff$\\$g_i^{-1}(V)\in\sigma_i$};

  % Maps
  \draw[blue!70!black, -{Stealth[length=2.4mm]}, thick]
    (1.4,1.2) to[bend left=12] node[above, font=\scriptsize, pos=0.5] {$g_1$} (5.55,0.45);
  \draw[red!70!black, -{Stealth[length=2.4mm]}, thick]
    (1.4,-1.2) to[bend right=12] node[below, font=\scriptsize, pos=0.5] {$g_2$} (5.55,-0.45);

  \node[font=\scriptsize, align=center] at (3.3,-2.9)
    {İleri itme: $\tau_{\mathrm{fin}}$ her $g_i$'yi sürekli kılan \emph{en ince} topoloji;\\
     bölüm topolojisi tek surjektif $q$'nun son topolojisidir.};
\end{tikzpicture}

\caption{Son topoloji: kaynak uzayları ortak bir hedefe iten en ince topoloji;
bölüm topolojisi tek surjektif haritanın son topolojisidir.}
\label{fig:ch13_son}
\end{figure}

\subsection{Yeni Örnek — ``En Kaba'' Olmanın Doğrudan Doğrulanması}

Başlangıç topolojisinin tanımı \emph{en kaba} olmaktır: $f$'yi sürekli kılar, ama
ondan daha kaba hiçbir topoloji kılamaz. Bunu süreklilik yüklemiyle doğrulayalım.

\begin{lstlisting}[language=Python]
ini      = [set(u) for u in tau_f.topology]
codom    = [set(u) for u in Y_d.topology]
coarser  = [set(), set(X_carrier)]   # indiscrete: tau_ini'den daha kaba

print("f, tau_ini'ye gore surekli mi:",
      is_continuous_finite_map(X_carrier, ini, list(Y_d.carrier), codom, f.mapping))
print("f, daha kaba (indiscrete) topolojiye gore surekli mi:",
      is_continuous_finite_map(X_carrier, coarser, list(Y_d.carrier), codom, f.mapping))
print("tau_ini eleman sayisi:", len(tau_f.topology))
\end{lstlisting}

\begin{lstlisting}[style=output]
f, tau_ini'ye gore surekli mi: True
f, daha kaba (indiscrete) topolojiye gore surekli mi: False
tau_ini eleman sayisi: 4
\end{lstlisting}

$\tau_{\mathrm{ini}}$ ile $f$ sürekli; aynı $f$ indiscrete topolojiyle sürekli
\textbf{değil} — çünkü $f^{-1}(\{0\})=\{0,1\}$ orada açık değildir. $\tau_{\mathrm{ini}}$
tam da bu açık kümeyi ekleyecek kadar incedir, fazlası değil.

\begin{karsiornek}
``Daha çok açık küme her zaman daha iyidir'' sanmayın. Discrete topoloji de $f$'yi
sürekli kılar (16 açık küme), ama \emph{en kaba} değildir — başlangıç topolojisi
yalnızca süreklilik için gereken 4 açık kümeyi koyar. Gereğinden ince bir topoloji
evrensel özelliği bozar: $\tau_{\mathrm{ini}}$'den daha ince bir $\tau'$ seçerseniz,
$f_\alpha\circ h$ sürekli olduğu hâlde $h$ sürekli olmayan bir $h$ bulunabilir.
\end{karsiornek}

\subsection{Yeni Örnek — Evrensel Özelliğin Doğrulanması}

İki harita $f,g: X \to \mathrm{discrete}(\{0,1\})$ ile $X=\{a,b\}$ üzerinde başlangıç
topolojisini kurar, ardından $h: Z \to X$ sürekliliğinin her bileşenin sürekliliğine
denk olduğunu gösteririz.

\begin{lstlisting}[language=Python]
Xc   = ["a", "b"]
Xph  = make_topology(Xc, set(), set(Xc))
Yd2  = discrete_topology(0, 1)
fA   = FiniteMap(domain=Xph, codomain=Yd2, name="fA", mapping={"a": 0, "b": 1})
gA   = FiniteMap(domain=Xph, codomain=Yd2, name="gA", mapping={"a": 1, "b": 0})

tau_X  = initial_topology_from_maps(Xc, [fA, gA])
Xopens = [set(u) for u in tau_X.topology]
Yopens = [set(u) for u in Yd2.topology]

Zc, Zopens = ["p", "q"], [set(), {"p"}, {"q"}, {"p", "q"}]
h  = {"p": "a", "q": "b"}
fh = {z: fA.mapping[h[z]] for z in Zc}     # fA . h
gh = {z: gA.mapping[h[z]] for z in Zc}     # gA . h

c_fh = is_continuous_finite_map(Zc, Zopens, list(Yd2.carrier), Yopens, fh)
c_gh = is_continuous_finite_map(Zc, Zopens, list(Yd2.carrier), Yopens, gh)
c_h  = is_continuous_finite_map(Zc, Zopens, Xc, Xopens, h)

print("tau_ini(X) eleman sayisi:", len(tau_X.topology))
print("fA.h surekli:", c_fh, "| gA.h surekli:", c_gh, "| h surekli:", c_h)
print("Evrensel ozellik saglaniyor mu:", c_h == (c_fh and c_gh))
\end{lstlisting}

\begin{lstlisting}[style=output]
tau_ini(X) eleman sayisi: 4
fA.h surekli: True | gA.h surekli: True | h surekli: True
Evrensel ozellik saglaniyor mu: True
\end{lstlisting}

$\tau_{\mathrm{ini}}(X)$ discrete$(\{a,b\})$'ye eşittir; $h$ süreklidir ancak ve ancak
hem $f_A\circ h$ hem $g_A\circ h$ sürekli olduğunda. Yüklem bu denkliği
\texttt{c\_h == (c\_fh and c\_gh)} ile doğrular.

% -----------------------------------------------------------------------
\section{Alıştırmalar}
\label{sec:initial_exercises}

\begin{enumerate}
  \item $f: \{a,b,c,d\} \to \mathrm{discrete}(\{0,1\})$, $f(a)=f(b)=0$, $f(c)=f(d)=1$
        için \texttt{initial\_topology\_from\_maps} kullanın ve sonucu
        \texttt{finite\_subspace} ile karşılaştırın.

  \item Çarpım özel halinde yalnızca $\pi_x$ kullanıldığında başlangıç topolojisi
        ne olur? Discrete'den daha kaba mı?

  \item $X=\{0,1,2,3,4\}$, $\sigma_X = $ discrete, $g: X \to \{0,1,2\}$ haritası
        için son topolojiyi manuel hesaplayın.

  \item \textit{(Teori)} $h: Z \to (X, \tau_{\mathrm{ini}})$ sürekliliğinin
        $f_\alpha \circ h$ sürekliliğine denk olduğunu ispatlayın.

  \item \textit{(Teori)} $\tau_{\mathrm{fin}}$'den daha ince hiçbir topoloji her
        $g_i$'yi sürekli kılamaz; bunu tanımdaki en-ince koşulundan türetin.

  \item \texttt{initial\_topology\_from\_maps} ile bir
        $f: \{0,1,2,3\} \to \mathrm{discrete}(\{0,1\})$, $f(\{0,1\})=0$,
        $f(\{2,3\})=1$ haritası için $\tau_{\mathrm{ini}}$ kurun. Sonra
        \texttt{is\_continuous\_finite\_map} ile $f$'nin $\tau_{\mathrm{ini}}$'ye göre
        sürekli ama indiscrete topolojiye göre sürekli olmadığını doğrulayın. Aynı
        $f$ discrete topolojiyle de sürekli midir? Bu, ``en kaba'' iddiasını nasıl
        destekler?

  \item \textit{(Evrensel özellik)} $X=\{a,b\}$ üzerinde
        $f(a)=0,f(b)=1$ ve $g(a)=1,g(b)=0$ (kodomen \texttt{discrete(\{0,1\})}) için
        $\tau_{\mathrm{ini}}$'yi hesaplayın. $Z=\mathrm{discrete}(\{p,q\})$,
        $h(p)=a$, $h(q)=b$ alın. \texttt{is\_continuous\_finite\_map} ile
        $f\circ h$, $g\circ h$ ve $h$ sürekliliklerini hesaplayıp $h$'nin
        sürekliliğinin bileşenlerin sürekliliğine denk olduğunu
        (\texttt{c\_h == (c\_fh and c\_gh)}) gösterin.
\end{enumerate}


% ── Part II: Metrik Uzaylar ────────────────────────────────────────────────────
\part{Metrik Uzaylar}

\chapter{Metrik Uzaylar}

Metrik uzaylar, mesafe fonksiyonuna dayanan topolojik yapılardır. Her metrik uzay
doğal olarak bir topolojik uzay oluşturur; bu yapı ``indüklenen topoloji'' olarak adlandırılır.

%-------------------------------------------------------------------
\section{Konu}

\subsection{Sezgisel Giriş — ``Mesafe Topolojiyi Doğurur''}

Metrik uzayı tek cümlede özetlemek gerekirse: \emph{noktalar arasındaki mesafe
fikrinden bir topoloji üretmektir.} Elimizde yalnız ``iki nokta ne kadar uzak?''
sorusunu yanıtlayan bir $d(x,y)$ fonksiyonu vardır; bundan ``bir noktanın yakınında
olmak'' kavramını --- yani açık kümeleri --- kurarız. Köprü tek bir nesnedir:
\textbf{açık yuvar.}

\begin{figure}[ht]
  \centering
  \begin{tikzpicture}[font=\small, etiket/.style={font=\scriptsize}]
  % --- Acik yuvar B(x,r): merkez x, yaricap r, kesikli sinir (sinir dahil degil) ---
  \def\R{2.4}
  % yuvar ici
  \fill[blue!12] (0,0) circle (\R);
  % kesikli cember: sinir B(x,r)'ye dahil DEGIL
  \draw[blue!70!black, very thick, dashed] (0,0) circle (\R);
  % merkez
  \fill[black] (0,0) circle (1.6pt);
  \node[below left=1pt] at (0,0) {$x$};
  % yaricap oku
  \draw[-{Stealth[length=2.4mm]}, thick] (0,0) -- (35:\R)
    node[midway, above, sloped, etiket] {$r$};
  % ic nokta y (dahil)
  \fill[teal!70!black] (-1.0,0.7) circle (1.6pt);
  \node[teal!60!black, etiket, below=1pt] at (-1.0,0.7) {$y$};
  \node[teal!60!black, etiket, align=center] at (-1.05,1.35) {$d(x,y)<r$\\(dahil)};
  % dis nokta z (haric)
  \fill[red!70!black] (62:3.05) circle (1.6pt);
  \node[red!60!black, etiket, right=1pt] at (62:3.05) {$z$};
  \node[red!60!black, etiket, align=center] at (3.0,2.65) {$d(x,z)\ge r$\\(haric)};

  \node[etiket, align=center, fill=blue!6, rounded corners, inner sep=4pt]
    at (0,-3.15)
    {$B(x,r)=\{y\in M : d(x,y)<r\}$ --- kesikli sınır açık yuvara \textbf{dahil değildir}.};
\end{tikzpicture}

  \caption{Açık yuvar $B(x,r)$: merkez $x$'ten $r$'den daha yakın olan tüm
    noktalar; kesikli sınır yuvara dahil değildir.}
\end{figure}

Bu yuvarları temel taş yaparak indüklenen topolojiyi kurarız: bir küme açıktır
$\Leftrightarrow$ her noktasının etrafına tümüyle içeride kalan bir açık yuvar sığar.

\begin{sezgi}
Metrik, sayısal bir ``uzaklık cetveli''dir; topoloji ise yalnız ``yakın mı, uzak
mı?'' sorusunu önemser. Aynı topolojiyi farklı cetveller (metrikler) üretebilir ---
önemli olan hangi noktaların hangi yuvarın \emph{içinde} kaldığıdır, mesafenin tam
sayısal değeri değil.
\end{sezgi}

Aynı düzlem üzerinde farklı metrikler farklı biçimde yuvarlar üretir ama çoğu zaman
\textbf{aynı topolojiyi} indükler:

\begin{figure}[ht]
  \centering
  \begin{tikzpicture}[font=\small, etiket/.style={font=\scriptsize}]
  % --- Uc metrigin birim yuvari B(0,1) ---
  \def\S{1.6}  % olcek

  % ---- 1) Oklid metrigi: cember ----
  \begin{scope}[shift={(0,0)}]
    \draw[->] (-\S-0.4,0) -- (\S+0.4,0);
    \draw[->] (0,-\S-0.4) -- (0,\S+0.4);
    \fill[blue!15] (0,0) circle (\S);
    \draw[blue!70!black, very thick] (0,0) circle (\S);
    \fill[black] (0,0) circle (1.2pt);
    \node[etiket, align=center] at (0,-\S-1.0)
      {\textbf{Öklid} $\ell^2$\\$\sqrt{x^2+y^2}<1$};
  \end{scope}

  % ---- 2) Taksi (Manhattan) metrigi: elmas / kare ----
  \begin{scope}[shift={(2*\S+1.8,0)}]
    \draw[->] (-\S-0.4,0) -- (\S+0.4,0);
    \draw[->] (0,-\S-0.4) -- (0,\S+0.4);
    \fill[teal!15] (\S,0) -- (0,\S) -- (-\S,0) -- (0,-\S) -- cycle;
    \draw[teal!70!black, very thick] (\S,0) -- (0,\S) -- (-\S,0) -- (0,-\S) -- cycle;
    \fill[black] (0,0) circle (1.2pt);
    \node[etiket, align=center] at (0,-\S-1.0)
      {\textbf{Taksi} $\ell^1$\\$|x|+|y|<1$};
  \end{scope}

  % ---- 3) Maksimum (Cebysev) metrigi: eksen-hizali kare ----
  \begin{scope}[shift={(4*\S+3.6,0)}]
    \draw[->] (-\S-0.4,0) -- (\S+0.4,0);
    \draw[->] (0,-\S-0.4) -- (0,\S+0.4);
    \fill[orange!18] (-\S,-\S) rectangle (\S,\S);
    \draw[orange!80!black, very thick] (-\S,-\S) rectangle (\S,\S);
    \fill[black] (0,0) circle (1.2pt);
    \node[etiket, align=center] at (0,-\S-1.0)
      {\textbf{Maksimum} $\ell^\infty$\\$\max(|x|,|y|)<1$};
  \end{scope}

  \node[etiket, align=center, fill=gray!8, rounded corners, inner sep=4pt]
    at (2*\S+1.8,\S+1.15)
    {Aynı düzlemde üç farklı metriğin birim yuvarı $B(0,1)$:
     biçimleri farklı, indükledikleri \emph{topoloji} aynıdır.};
\end{tikzpicture}

  \caption{Üç farklı metriğin birim yuvarı $B(0,1)$: Öklid çember, taksi
    (Manhattan) elması, maksimum metrik karesi --- biçimleri farklı, indükledikleri
    topoloji aynı.}
\end{figure}

\subsection{Metrik Aksiyomları}

\begin{definition}[Metrik Uzay]
$M$ kümesi ve $d: M \times M \to [0,\infty)$ verilsin.
$(M, d)$ bir \textbf{metrik uzay}, $d$ ise \textbf{metrik} ise:
\begin{enumerate}
  \item[(M1)] \textbf{Özdeşlik:} $d(x,y) = 0 \Leftrightarrow x = y$
  \item[(M2)] \textbf{Simetri:} $d(x,y) = d(y,x)$
  \item[(M3)] \textbf{Üçgen:} $d(x,z) \leq d(x,y) + d(y,z)$
\end{enumerate}
\end{definition}

Üçgen eşitsizliği, üç aksiyom içinde en derin olanıdır: ``$x$'ten $z$'ye doğrudan
gitmek, $y$'ye uğrayıp gitmekten kısadır'' der.

\begin{figure}[ht]
  \centering
  \begin{tikzpicture}[font=\small, etiket/.style={font=\scriptsize}]
  % --- Ucgen esitsizligi: d(x,z) <= d(x,y) + d(y,z) ---
  \coordinate (x) at (0,0);
  \coordinate (y) at (5.2,1.9);
  \coordinate (z) at (6.6,-0.6);

  % dolayli yol x -> y -> z
  \draw[teal!70!black, very thick] (x) -- (y)
    node[midway, above left, etiket, teal!60!black] {$d(x,y)$};
  \draw[teal!70!black, very thick] (y) -- (z)
    node[midway, right, etiket, teal!60!black] {$d(y,z)$};
  % dogrudan yol x -> z (kestirme)
  \draw[red!75!black, very thick, dashed] (x) -- (z)
    node[midway, below, etiket, red!60!black] {$d(x,z)$};

  % noktalar
  \foreach \p/\pos in {x/{below left}, y/{above}, z/{right}} {
    \fill[black] (\p) circle (1.8pt);
    \node[\pos=2pt] at (\p) {$\p$};
  }

  \node[etiket, align=center, fill=red!6, rounded corners, inner sep=4pt]
    at (3.2,-2.0)
    {\textbf{Üçgen eşitsizliği:}\;\;$d(x,z)\;\le\;d(x,y)+d(y,z)$\\[2pt]
     Doğrudan yol (kesikli kırmızı), $y$ üzerinden geçen dolaylı yoldan kısadır.};
\end{tikzpicture}

  \caption{Üçgen eşitsizliği: $x$'ten $z$'ye doğrudan mesafe, $y$ üzerinden geçen
    dolaylı yoldan kısadır.}
\end{figure}

\begin{karsiornek}
\textbf{Üçgen eşitsizliği şart.} $\{A,B,C\}$ üzerinde $d(A,B)=1$, $d(B,C)=1$ ama
$d(A,C)=5$ koyalım (özdeşlik ve simetri korunsun). M1 ve M2 sağlanır, ama
$d(A,C)=5 > d(A,B)+d(B,C)=2$ olduğundan M3 ihlal edilir. Bu fonksiyon bir
\emph{metrik değildir}; ürettiği ``yuvarlar'' tutarlı bir topoloji vermez.
\texttt{validate\_metric} bu durumu yakalar (Alıştırma K4).
\end{karsiornek}

\subsection{Temel Kavramlar}

\begin{itemize}
  \item \textbf{Açık top:} $B(x, r) = \{y \in M : d(x,y) < r\}$
  \item \textbf{Kapalı top:} $\bar{B}(x, r) = \{y \in M : d(x,y) \leq r\}$
  \item \textbf{İndüklenen topoloji:} $\tau_d = \{U : \forall x \in U,\,\exists\varepsilon>0,\,B(x,\varepsilon) \subseteq U\}$
  \item \textbf{Çap:} $\mathrm{diam}(A) = \sup\{d(x,y) : x,y \in A\}$
\end{itemize}

\subsection{Standart Metrikler}

\begin{table}[h]
\centering
\begin{tabular}{lll}
\toprule
\textbf{Metrik} & \textbf{Tanım} & \textbf{Uzay} \\
\midrule
Öklid & $d(x,y) = |x-y|$ & $\mathbb{R}$ \\
Ayrık & $d(x,y) = 0$ veya $1$ & Herhangi $X$ \\
$\ell^\infty$ (max) & $\max_i d_i(x_i,y_i)$ & Çarpım uzayı \\
Sınırlı (capped) & $d'(x,y) = \min(d(x,y),1)$ & Eşdeğer topoloji \\
\bottomrule
\end{tabular}
\caption{Standart metrikler}
\end{table}

%-------------------------------------------------------------------
\section{Teoremler}

\begin{theorem}
Her metrik uzay Hausdorff (T2) ve $1^{\text{st}}$ sayılabilirdir.
\end{theorem}

\begin{proof}[İspat eskizi]
(T2) $x \neq y$ olsun; o hâlde $r = d(x,y) > 0$ (M1). $B(x, r/2)$ ve $B(y, r/2)$
yuvarları ayrıktır: ortak bir $z$ noktası olsaydı, üçgen eşitsizliğinden (M3)
$r = d(x,y) \leq d(x,z) + d(z,y) < r/2 + r/2 = r$ çelişkisi çıkardı. Demek ki $x$ ve
$y$ açık komşuluklarla ayrılır. (1.\ sayılabilir) Her $x$ için sayılabilir
$\{B(x, 1/n) : n \in \mathbb{N}\}$ ailesi yerel bazdır: $x \in U$ açık ise
$B(x,\varepsilon) \subseteq U$ olacak $\varepsilon > 0$ vardır; $1/n < \varepsilon$
seçince $B(x, 1/n) \subseteq U$.
\end{proof}

\begin{theorem}[Sınırlı Metrik Eşdeğerliği]
$d'(x,y) = \min(d(x,y),1)$ ile $d$ aynı topolojiyi indükler.
\end{theorem}

\begin{proof}[İspat eskizi]
İki metrik aynı topolojiyi indükler $\Leftrightarrow$ her $d$-açık yuvarın içine bir
$d'$-açık yuvar sığar ve tersi. Yarıçapı $r \leq 1$ olan yuvarlarda
$B_d(x, r) = B_{d'}(x, r)$ tam çakışır, çünkü $d$ ile $d'$ yalnız mesafesi $\geq 1$
olan çiftlerde ayrışır. Topoloji ``yeterince küçük yuvarlar'' tarafından
belirlendiğinden büyük yarıçaplı farklılık topolojiyi değiştirmez; $\tau_d = \tau_{d'}$.
\end{proof}

\begin{theorem}
$\mathbb{R}^n$ üzerindeki max, sum ve Öklid metrikleri topolojik olarak eşdeğerdir.
\end{theorem}

\begin{proof}[İspat eskizi]
Üç metrik karşılıklı sabit çarpanlarla sınırlanır:
$d_\infty \leq d_2 \leq d_1 \leq n\,d_\infty$ (max $\leq$ Öklid $\leq$ taksi
$\leq n\cdot$max). Böyle bir ``sandviç'' her metriğin yuvarının içine ötekinin
yuvarını sığdırır (ör.\ $B_\infty(x, r/n) \subseteq B_1(x, r) \subseteq B_\infty(x, r)$).
Karşılıklı içerme $\Rightarrow$ aynı açık kümeler $\Rightarrow$ aynı topoloji.
Bölüm 1.0'daki birim-yuvar figürü bu üç biçimi yan yana gösterir.
\end{proof}

%-------------------------------------------------------------------
\section{Algoritmalar}

\subsection{validate\_metric — $O(|X|^3)$}

\begin{algorithm}
\caption{ValidateMetric$(X, d)$}
\begin{algorithmic}[1]
\For{each $x \in X$}
    \If{$d(x,x) \neq 0$} \Return False \EndIf
    \For{each $y \neq x$}
        \If{$d(x,y) = 0$} \Return False \EndIf
    \EndFor
\EndFor
\For{each $x,y \in X$}
    \If{$d(x,y) \neq d(y,x)$} \Return False \EndIf
\EndFor
\For{each $x,y,z \in X$}
    \If{$d(x,z) > d(x,y) + d(y,z)$} \Return False \EndIf
\EndFor
\State \Return True
\end{algorithmic}
\end{algorithm}

Karmaşıklık: $O(|X|^3)$ — üçgen eşitsizliği dominates.

\subsection{open\_ball — $O(|X|)$}

\begin{algorithm}
\caption{OpenBall$(x, r, X, d)$}
\begin{algorithmic}[1]
\State \Return $\{y \in X : d(x,y) < r\}$
\end{algorithmic}
\end{algorithm}

%-------------------------------------------------------------------
\section{pytop API}

\begin{lstlisting}[language=Python]
from pytop.metric_spaces import (
    FiniteMetricSpace,
    SymbolicMetricSpace,
    validate_metric,
)
from pytop import (
    open_ball,
    closed_ball,
    distance_to_subset,
    diameter_of_subset,
    is_bounded_subset,
    capped_metric,
)
\end{lstlisting}

\texttt{FiniteMetricSpace(carrier=tuple, distance=dict)} — mesafe sözlüğü \texttt{(a,b): float}.

\texttt{open\_ball(fms, point, radius)} $\to$ \texttt{set} döner (Result değil).

\texttt{capped\_metric(distance\_fn, cap=1.0)} $\to$ \texttt{Callable} döner.

\textbf{Metrik $\to$ TDA köprüsü:} Bir \texttt{FiniteMetricSpace}
(\texttt{carrier} + \texttt{distance}) doğrudan
\texttt{pytop.persistent\_homology(space, max\_dimension=, max\_scale=)} fonksiyonuna
beslenebilir; bu, ham nokta bulutunu Vietoris–Rips filtrasyonuna çevirip kalıcı
homolojiyi hesaplar. Dönen \texttt{PersistencePair} alanları: \texttt{dimension},
\texttt{birth}, \texttt{death}, \texttt{persistence}, \texttt{is\_essential}
(Örnek 5.8).

%-------------------------------------------------------------------
\section{Örnekler}

\subsection*{Örnek 5.1 — Ayrık Metrik Uzay}

\begin{lstlisting}[language=Python]
points = [1, 2, 3, 4]
dist_discrete = {(a, b): (0 if a == b else 1)
                 for a in points for b in points}
fms = FiniteMetricSpace(carrier=tuple(points), distance=dist_discrete)
print("carrier:", fms.carrier)
print("d(1,2):", fms.distance[(1,2)])
\end{lstlisting}

\begin{verbatim}
carrier: (1, 2, 3, 4)
d(1,2): 1
\end{verbatim}

\subsection*{Örnek 5.3 — open\_ball}

\begin{lstlisting}[language=Python]
print("B(1, 0.5) =", open_ball(fms, 1, 0.5))
print("B(1, 1.5) =", open_ball(fms, 1, 1.5))
\end{lstlisting}

\begin{verbatim}
B(1, 0.5) = {1}
B(1, 1.5) = {1, 2, 3, 4}
\end{verbatim}

\subsection*{Örnek 5.6 — capped\_metric}

\begin{lstlisting}[language=Python]
cap_fn = capped_metric(fms_line.distance, cap=1.0)
print("d(0,3) capped:", cap_fn(0, 3))
print("d(0,1) capped:", cap_fn(0, 1))
\end{lstlisting}

\begin{verbatim}
d(0,3) capped: 1.0
d(0,1) capped: 1.0
\end{verbatim}

\subsection*{Örnek 5.7 — Üç Metrik, Üç Yuvar, Aynı Geçerlilik}

Düzlemdeki aynı beş nokta üzerinde Öklid, taksi ($\ell^1$) ve maksimum
($\ell^\infty$) metriklerini kurarız: üçü de geçerli metriktir ama aynı çift için
farklı mesafe verir (Bölüm 1.0 birim-yuvar figürünün sayısal karşılığı).

\begin{lstlisting}[language=Python]
import math
from pytop.metric_spaces import FiniteMetricSpace, validate_metric

grid = [(0, 0), (1, 0), (0, 1), (1, 1), (2, 1)]
carrier = tuple(range(len(grid)))

def euclid(i, j):    return math.dist(grid[i], grid[j])
def taxicab(i, j):   return abs(grid[i][0]-grid[j][0]) + abs(grid[i][1]-grid[j][1])
def chebyshev(i, j): return max(abs(grid[i][0]-grid[j][0]), abs(grid[i][1]-grid[j][1]))

for name, d in [("Oklid", euclid), ("Taksi", taxicab), ("Maksimum", chebyshev)]:
    dist = {(i, j): d(i, j) for i in carrier for j in carrier}
    fms = FiniteMetricSpace(carrier=carrier, distance=dist)
    print(name, "d(0,4)=", round(dist[(0,4)], 4), "metrik?", validate_metric(fms).status)
\end{lstlisting}

\begin{verbatim}
Oklid    d(0,4)= 2.2361 metrik? true
Taksi    d(0,4)= 3.0    metrik? true
Maksimum d(0,4)= 2.0    metrik? true
\end{verbatim}

Mesafeler farklı (Öklid $\sqrt5$, taksi $3$, maksimum $2$), ama üçü de aynı
topolojiyi indükler (Teorem 2.3).

\subsection*{Örnek 5.8 — Metrik $\to$ Kalıcı Homoloji (TDA Köprüsü)}

Metrik uzayın hesaplama gücüne köprü: bir nokta bulutundan \texttt{FiniteMetricSpace}
kurup doğrudan \texttt{persistent\_homology}'ye veririz. Çember üzerinde örneklenmiş
8 nokta --- ``delik'' (1-boyutlu döngü) topolojik imza olarak ortaya çıkmalıdır.

\begin{lstlisting}[language=Python]
import math, pytop
from pytop.metric_spaces import FiniteMetricSpace

n = 8
pts = [(math.cos(2*math.pi*k/n), math.sin(2*math.pi*k/n)) for k in range(n)]
carrier = tuple(range(n))
dist = {(i, j): math.dist(pts[i], pts[j]) for i in carrier for j in carrier}
circle = FiniteMetricSpace(carrier=carrier, distance=dist)

pairs = pytop.persistent_homology(circle, max_dimension=2, max_scale=2.5)
components = [p for p in pairs if p.dimension == 0 and p.is_essential]
loops = [p for p in pairs if p.dimension == 1]
loop = loops[0]
print("bilesen (H0):", len(components), " dongu (H1):", len(loops))
print("dongu: dogum=", round(loop.birth,4), "olum=", round(loop.death,4))
\end{lstlisting}

\begin{verbatim}
bilesen (H0): 1  dongu (H1): 1
dongu: dogum= 0.7654 olum= 1.8478
\end{verbatim}

Tam beklenen sonuç: çember bağlantılı olduğundan tek kalıcı bileşen ($H_0$), bir
deliği olduğundan tek kalıcı döngü ($H_1$). Döngü, komşu noktalar bağlandığında
doğar ($\approx 0.765$) ve yarıçap deliği dolduracak kadar büyüyünce ölür
($\approx 1.848$). Ham koordinatlardan (\texttt{math.dist}) pytop topolojik bir
özelliği --- deliğin varlığını --- sayısal olarak çıkardı.

%-------------------------------------------------------------------
\section{Alıştırmalar}

\subsection*{Kodlama}

\begin{enumerate}[label=K\arabic*.]
\item $\{A,B,C,D\}$ üzerinde bir metrik tanımlayın ve \texttt{validate\_metric} ile
      doğrulayın. Üçgen eşitsizliğini ihlal eden bir örnek de deneyin.
\item Öklid metrikli $\{0,\ldots,9\}$ üzerinde $B(5,3.0)$ ve $\bar{B}(5,3.0)$
      hesaplayın.
\item \texttt{normalized\_metric} kullanarak bir metriği $[0,1]$'e normalleştirin.
\item \S 1.0'daki karşı-örneği koda dökün: $\{A,B,C,D\}$ üzerinde $d(A,C)=5$,
      $d(A,B)=d(B,C)=1$ olan bir mesafe sözlüğü kurun (geri kalanını ayrık metrikle
      doldurun) ve \texttt{validate\_metric} ile üçgen eşitsizliğinin ihlal
      edildiğini gösterin; \texttt{justification} alanını yazdırın.
\item \emph{(TDA köprüsü)} Birim kareyi örnekleyen \textbf{dolu} bir $3\times 3$
      ızgara (9 nokta) üzerinde Öklid \texttt{FiniteMetricSpace} kurun ve
      \texttt{persistent\_homology(space, max\_dimension=2, max\_scale=2.5)}
      çalıştırın. $H_1$ döngülerinin \texttt{persistence} değerlerine bakın:
      Örnek 5.8'deki çemberin uzun ömürlü ($\approx 1.08$) döngüsüne benzer
      \textbf{kalıcı} bir döngü var mı? Dolu karenin neden gerçek bir deliği yoktur?
\end{enumerate}

\subsection*{Teori}

\begin{enumerate}[label=T\arabic*.]
\item Her metrik uzayın Hausdorff olduğunu ispatlayın.
\item $d' = \min(d,1)$ ve $d$ aynı topolojiyi indüklediğini gösterin.
\item Üçgen eşitsizliğinin (M3) Teorem 2.1'deki Hausdorff ispatında tam olarak
      nerede kullanıldığını açıklayın. (İpucu: $B(x, r/2)$ ve $B(y, r/2)$'nin ayrık
      olduğunu gösteren adım hangisidir?)
\end{enumerate}

\chapter{Metrik Tamlık ve Kompaktlık}

Metrik tamlık (completeness), bir metrik uzayda Cauchy dizilerinin yakınsayıp
yakınsamadığını araştıran temel bir kavramdır.

%-------------------------------------------------------------------
\section{Konu}

\subsection{Cauchy Dizisi ve Tamlık}

\begin{definition}[Cauchy Dizisi]
$(M, d)$ metrik uzayında $\{x_n\}$ dizisi \textbf{Cauchy} ise:
\[
\forall\varepsilon>0,\;\exists N \in \mathbb{N}:\;
m,n \geq N \Rightarrow d(x_m, x_n) < \varepsilon.
\]
\end{definition}

\begin{definition}[Tam Metrik Uzay]
$(M, d)$ \textbf{tam} ise: Her Cauchy dizisi $M$'de bir limite yakınsır.
\end{definition}

\begin{figure}[ht]
\centering
\begin{tikzpicture}[font=\small]
  % eksen
  \draw[-{Stealth[length=2.4mm]}] (-0.4,0) -- (9.2,0) node[right] {$n$};
  % limit cizgisi
  \draw[gray!55, dashed] (0,2.1) -- (8.6,2.1) node[right, black] {$L$};
  % Cauchy dizisinin terimleri: L'ye yaklasan, salinim azalan
  \foreach \n/\y in {0/3.5, 1/0.8, 2/3.0, 3/1.2, 4/2.6, 5/1.7, 6/2.35, 7/1.9, 8/2.18}
    \fill[blue!70!black] (\n,\y) circle (1.6pt);
  % kuyrugun capi -> 0 vurgusu (saga dogru daralan bant)
  \fill[orange!18] (5.5,1.55) -- (8.5,2.0) -- (8.5,2.2) -- (5.5,1.85) -- cycle;
  \node[orange!75!black, font=\scriptsize] at (7.0,1.15) {kuyruğun çapı $\to 0$};
  % ilk terimler arasi buyuk mesafe
  \draw[<->, red!70!black] (0,3.5) -- node[right, font=\scriptsize] {büyük} (0,0.8);
  \node[font=\scriptsize] at (1.6,-0.5) {$x_n$ terimleri};
  \node[blue!70!black] at (7.3,3.4) {$\{x_n\}$ Cauchy};
\end{tikzpicture}

\caption{Cauchy dizisinin terimleri birbirine yaklaşır: kuyruğun çapı sıfıra gider.}
\end{figure}

\begin{sezgi}
Bir Cauchy dizisini, ilerledikçe birbirine sokulan bir terim kalabalığı gibi
düşünün. İlk terimler birbirinden uzaktayken, yeterince ileri gidildiğinde tüm
terimler bir noktada toplanır: ``kuyruğun çapı''
$\bigl(\sup_{m,n\geq N} d(x_m, x_n)\bigr)$ sıfıra iner. \textbf{Tam} uzayda bu
kalabalığın sıkıştığı yerde gerçekten bir nokta (limit) durur; tam olmayan
uzayda ise o yerde bir \emph{boşluk} olabilir.
\end{sezgi}

\subsection{Tamamen Sınırlılık}

$A \subseteq M$ \textbf{tamamen sınırlı} ise: her $\varepsilon>0$ için $A$'nın
sonlu bir $\varepsilon$-net kümesi vardır.
\[
S = \{s_1,\ldots,s_n\} \subseteq M,\quad
\forall x \in A\;\exists i:\; d(x, s_i) < \varepsilon.
\]

\subsection{Metrik Kompaktlık Karakterizasyonu}

\[
(M, d) \text{ tam} \wedge \text{tamamen sınırlı}
\Longleftrightarrow
(M, d) \text{ kompakt.}
\]

\begin{table}[h]
\centering
\begin{tabular}{llll}
\toprule
\textbf{Uzay} & \textbf{Tam?} & \textbf{Tam. sınırlı?} & \textbf{Kompakt?} \\
\midrule
$\mathbb{R}$ & \checkmark & $\times$ & $\times$ \\
$[0,1]$ & \checkmark & \checkmark & \checkmark \\
$(0,1)$ & $\times$ & \checkmark & $\times$ \\
$\mathbb{Q}$ & $\times$ & $\times$ & $\times$ \\
Sonlu metrik & \checkmark & \checkmark & \checkmark \\
\bottomrule
\end{tabular}
\caption{Tamlık ve kompaktlık örnekleri}
\end{table}

\begin{figure}[ht]
\centering
\begin{tikzpicture}[font=\small]
  % --- Q dogrusu (ustte): limit yok ---
  \draw[-{Stealth[length=2.4mm]}] (-0.3,2.4) -- (6.6,2.4) node[right] {$\mathbb{Q}$};
  \draw[red!70!black, dashed] (4.2,2.05) -- (4.2,2.75);
  \node[red!70!black, font=\scriptsize] at (4.2,3.05) {$\sqrt{2}\notin\mathbb{Q}$};
  \foreach \x in {1.0, 2.0, 2.8, 3.4, 3.8, 4.0}
    \fill[blue!70!black] (\x,2.4) circle (1.5pt);
  \node[blue!70!black, font=\scriptsize] at (1.7,2.0) {$x_n\to\sqrt{2}$};
  \node[red!70!black] at (5.7,2.95) {limit \emph{yok}};

  % --- R dogrusu (altta): limit var ---
  \draw[-{Stealth[length=2.4mm]}] (-0.3,0.4) -- (6.6,0.4) node[right] {$\mathbb{R}$};
  \fill[green!45!black] (4.2,0.4) circle (2.2pt);
  \node[green!45!black, font=\scriptsize] at (4.2,0.05) {$\sqrt{2}$};
  \foreach \x in {1.0, 2.0, 2.8, 3.4, 3.8, 4.0}
    \fill[blue!70!black] (\x,0.4) circle (1.5pt);
  \node[green!45!black] at (5.7,0.95) {limit \emph{var}};

  % ayrim
  \node[font=\scriptsize] at (3.1,1.4) {aynı Cauchy dizisi};
\end{tikzpicture}

\caption{$\mathbb{Q}$'da $\sqrt{2}$'ye yakınsayan Cauchy dizisinin limiti yoktur;
aynı dizi $\mathbb{R}$'de yakınsar.}
\end{figure}

\begin{karsiornek}
Rasyonel sayılar $\mathbb{Q}$, Öklid metriğiyle \textbf{tam değildir}.
$x_1 = 1,\ x_2 = 1.4,\ x_3 = 1.41,\ x_4 = 1.414,\ \ldots$ dizisi $\sqrt{2}$'nin
ondalık açılımıdır: her terim rasyoneldir ve dizi Cauchy'dir
($d(x_m, x_n) \to 0$). Ama limit $\sqrt{2} \notin \mathbb{Q}$ olduğundan dizi
$\mathbb{Q}$ içinde \emph{hiçbir} noktaya yakınsamaz. Aynı dizi $\mathbb{R}$'de
$\sqrt{2}$'ye yakınsar --- fark uzayda ``boşluk'' olup olmamasıdır.
\end{karsiornek}

\paragraph{Tamlama (completion).}
Her metrik uzay $(M, d)$ bir \textbf{tam} uzayın $(\widehat{M}, \widehat{d})$
yoğun alt-uzayı olarak gömülebilir; $\widehat{M}$'ye $M$'nin \emph{tamlaması}
denir. Sezgisel olarak tamlama, eksik limitleri (``boşlukları'') ekleyerek uzayı
kapatır: $\mathbb{R} = \widehat{\mathbb{Q}}$ bu inşanın klasik örneğidir.

\begin{figure}[ht]
\centering
\begin{tikzpicture}[font=\small,
    kutu/.style={draw, rounded corners, align=center, inner sep=8pt}]
  % delikli uzay
  \node[kutu, fill=orange!8] (q) at (0,0)
    {$\mathbb{Q}$ (tam değil)\\[3pt] \emph{boşluklar} var};
  % delikleri temsil eden noktalar + bosluklar
  \foreach \x in {-0.9, -0.3, 0.4} \fill[blue!70!black] (\x,-0.95) circle (1.3pt);
  \draw[red!70!black, thick] (-0.6,-0.95) ++(-0.05,0.18) -- ++(0.1,-0.36);
  \draw[red!70!black, thick] (0.05,-0.95) ++(-0.05,0.18) -- ++(0.1,-0.36);
  \node[red!70!black, font=\scriptsize] at (0,-1.45) {eksik limitler};

  % tam uzay
  \node[kutu, fill=green!8] (r) at (7.6,0)
    {$\mathbb{R}=\widehat{\mathbb{Q}}$ (tam)\\[3pt] her Cauchy yakınsar};

  % tamlama oku (iki kutu arasi)
  \draw[-{Stealth[length=3mm]}, thick]
    (q.east) -- node[above] {tamlama $\widehat{\;\cdot\;}$}
                node[below, font=\scriptsize] {boşlukları doldur} (r.west);

  \foreach \x in {6.7, 7.3, 7.6, 7.9, 8.5} \fill[green!45!black] (\x,-0.95) circle (1.3pt);
  \node[green!45!black, font=\scriptsize] at (7.6,-1.45) {süreklilik};
\end{tikzpicture}

\caption{Tamlama: $\mathbb{Q}$'nun boşlukları doldurularak tam uzay $\mathbb{R}$
elde edilir.}
\end{figure}

%-------------------------------------------------------------------
\section{Teoremler}

\begin{theorem}[Heine-Borel Metrik Genelleştirmesi]
Metrik uzayda kompaktlık $\Leftrightarrow$ tamlık $\wedge$ tamamen sınırlılık.
\end{theorem}

\begin{theorem}[Banach Sabit-Nokta Teoremi]
$(M, d)$ tam metrik uzay; $T: M \to M$ büzülme ($K < 1$)
$\Rightarrow$ $T$'nin eşsiz sabit noktası $x^*$ vardır: $T(x^*) = x^*$.
\end{theorem}

\begin{proof}[İspat eskizi]
$K < 1$ büzülme sabiti olsun. Herhangi bir $x_0$ alıp $x_{n+1} = T(x_n)$
iterasyonunu kurun. Büzülmeden
$d(x_{n+1}, x_n) \leq K\, d(x_n, x_{n-1}) \leq K^n\, d(x_1, x_0)$. Üçgen
eşitsizliği ve geometrik seri ile
$d(x_{n+m}, x_n) \leq \tfrac{K^n}{1-K}\, d(x_1, x_0)$; $K < 1$ olduğundan bu
ifade $n \to \infty$ iken sıfıra gider, yani $\{x_n\}$ Cauchy'dir.
\textbf{Tamlık} kullanılarak dizi bir $x^*$'ye yakınsar. $T$ sürekli (Lipschitz)
olduğundan $T(x^*) = T(\lim x_n) = \lim x_{n+1} = x^*$: $x^*$ sabit noktadır.
\emph{Eşsizlik:} $T(p) = p$, $T(q) = q$ olsaydı
$d(p,q) = d(Tp, Tq) \leq K\, d(p,q)$, $K < 1$ ile bu ancak $d(p,q) = 0$, yani
$p = q$ ise mümkündür.
\end{proof}

\begin{theorem}[Tam Uzayın Kapalı Alt-Uzayı Tamdır]
$(M, d)$ tam ve $A \subseteq M$ kapalı ise $(A, d)$ de tamdır.
\end{theorem}

\begin{proof}[İspat eskizi]
$\{a_n\} \subseteq A$ bir Cauchy dizisi olsun. $\{a_n\}$ aynı zamanda $M$ içinde
Cauchy'dir ve $M$ tam olduğundan bir $a \in M$ limitine yakınsar. $A$
\textbf{kapalı} olduğundan kendi limit noktalarını içerir; $a_n \to a$ ve
$a_n \in A$ olması $a \in \overline{A} = A$ demektir. Böylece her Cauchy dizisi
$A$ içinde bir limite yakınsar: $A$ tamdır.
\end{proof}

\begin{theorem}[Baire Kategorisi Teoremi]
Tam metrik uzay 1.\ kategoriden değildir.
\end{theorem}

%-------------------------------------------------------------------
\section{Algoritmalar}

\subsection{Sonlu Metrikte is\_complete ve is\_totally\_bounded}

Sonlu metrik uzayda her Cauchy dizisi eninde sonunda sabit: trivially tam.
Taşıyıcı kendisi $\varepsilon$-net: trivially tamamen sınırlı.
Her ikisi: $O(1)$.

\subsection{metric\_compactness\_check}

\begin{algorithm}
\caption{MetricCompactnessCheck$(M, d)$}
\begin{algorithmic}[1]
\State complete $\leftarrow$ is\_complete$(M, d)$
\State totally\_bounded $\leftarrow$ is\_totally\_bounded$(M, d)$
\State \Return complete $\wedge$ totally\_bounded
\end{algorithmic}
\end{algorithm}

%-------------------------------------------------------------------
\section{pytop API}

\begin{lstlisting}[language=Python]
from pytop.metric_completeness import (
    is_complete,
    is_totally_bounded,
    metric_compactness_check,
    analyze_metric_completeness,
)
\end{lstlisting}

%-------------------------------------------------------------------
\section{Örnekler}

\subsection*{Örnek 5.1 — Sonlu Metrik: Tam ve Kompakt}

\begin{lstlisting}[language=Python]
points = ['A', 'B', 'C', 'D']
dist = {(a, b): (0 if a == b else 1)
        for a in points for b in points}
fms = FiniteMetricSpace(carrier=tuple(points), distance=dist)
print("is_complete:", is_complete(fms).status)
print("is_totally_bounded:", is_totally_bounded(fms).status)
mc = metric_compactness_check(fms)
print("metric_compactness:", mc.status, mc.value)
\end{lstlisting}

\begin{verbatim}
is_complete: true
is_totally_bounded: true
metric_compactness: true True
\end{verbatim}

\subsection*{Örnek 5.4 — analyze\_metric\_completeness}

\begin{lstlisting}[language=Python]
r = analyze_metric_completeness(fms)
print("status:", r.status)
print("value:", r.value)
\end{lstlisting}

\begin{verbatim}
status: true
value: {'is_complete': True, 'is_totally_bounded': True, 'metric_compact': True}
\end{verbatim}

\subsection*{Örnek 5.6 — Rasyoneller $\mathbb{Q}$: Tam Olmayan Uzay (Karşı-Örnek)}

\texttt{rationals\_metric()} sembolik bir uzaydır; \texttt{is\_complete} ondalık
taramayla karara bağlanamadığı için \texttt{unknown} döner --- fakat uzayın
\textbf{etiketleri} (tags) $\mathbb{Q}$'nun tam olmadığını kayıt altına alır
(\texttt{'not\_complete'}).

\begin{lstlisting}[language=Python]
from pytop import rationals_metric
from pytop.metric_completeness import is_complete

qm = rationals_metric()
res = is_complete(qm)
print("is_complete:", res.status)
print("'not_complete' tag:", 'not_complete' in qm.tags)
print("justification:", res.justification[0])
\end{lstlisting}

\begin{verbatim}
is_complete: unknown
'not_complete' tag: True
justification: Completeness has exact support only for explicit finite metric spaces.
\end{verbatim}

\subsection*{Örnek 5.7 — Banach ve Sabit-Nokta Profilleri}

\texttt{get\_fixed\_point\_profiles()} sabit-nokta teorisinin küratörlü
profillerini verir; \texttt{fixed\_point\_stability\_summary()} bunları
kararlılığa göre gruplar. Banach büzülme teoremi \textbf{çekici (attracting)}
sabit nokta profiline karşılık gelir.

\begin{lstlisting}[language=Python]
from pytop import get_fixed_point_profiles, fixed_point_stability_summary

profiles = get_fixed_point_profiles()
print("profil sayisi:", len(profiles))
for p in profiles:
    print(f"  {p.key:24s}: {p.stability}")

summary = fixed_point_stability_summary()
print("stable:", summary['stable'])
\end{lstlisting}

\begin{verbatim}
profil sayisi: 5
  attracting_fixed_point  : stable
  repelling_fixed_point   : unstable
  neutral_fixed_point     : neutral
  brouwer_fixed_point     : not_applicable
  periodic_point_n        : not_applicable
stable: ['attracting_fixed_point']
\end{verbatim}

%-------------------------------------------------------------------
\section{Alıştırmalar}

\subsection*{Kodlama}

\begin{enumerate}[label=K\arabic*.]
\item Rasyonel sayılar $\mathbb{Q}$ (sembolik) için \texttt{is\_complete} sonucunu
      kontrol edin.
\item Kendi 4-noktalı metrik uzayınızı oluşturun ve \texttt{metric\_compactness\_check}
      uygulayın.
\item \texttt{analyze\_metric\_completeness} çalıştırın ve \texttt{value} sözlüğünü inceleyin.
\item \texttt{rationals\_metric()} oluşturun; \texttt{is\_complete(...).status} ile
      \texttt{'not\_complete'} etiketinin uzayın \texttt{tags} kümesinde olup olmadığını
      yazdırın. Statü neden \texttt{unknown}, etiket neden \texttt{True}?
\item \texttt{get\_fixed\_point\_profiles()} listesini gezin ve
      \texttt{stability == 'stable'} olan profilin \texttt{key} alanını yazdırın.
      Bu profil hangi teoreme (Banach mı?) karşılık gelir?
\end{enumerate}

\subsection*{Teori}

\begin{enumerate}[label=T\arabic*.]
\item Banach sabit-nokta teoremini sözlü olarak açıklayın ve bir uygulama verin.
\item Kapalı $[0,1]$'in tam, açık $(0,1)$'in tam olmadığını gösterin.
      (Hint: $1/n$ dizisi Cauchy ama $0 \notin (0,1)$.)
\item ``Tam uzayın kapalı alt-uzayı tamdır'' teoremini ispatlayın. Ayrıca tam bir
      uzayda \textbf{açık} bir alt-uzayın tam olması gerekmediğini bir örnekle
      gösterin. (Hint: $\mathbb{R}$ tam, ama $(0,1) \subseteq \mathbb{R}$ açık ve
      tam değil.)
\item $\mathbb{Q}$'nun tam olmadığını $\sqrt{2}$'nin ondalık açılımı üzerinden
      ispatlayın: dizinin Cauchy olduğunu, fakat $\mathbb{Q}$ içinde limitinin
      bulunmadığını gösterin. Tamlama kavramıyla bu ``boşluğun'' nasıl
      kapatıldığını açıklayın.
\end{enumerate}

\chapter{Metrik Fonksiyonlar ve Sözleşmeler}

Bu bölümde metrik uzaylar arasındaki fonksiyon türleri (izometri, Lipschitz, sürekli)
ve bunların sınıflandırılması incelenmektedir.

%-------------------------------------------------------------------
\section{Konu}

\subsection{Metrik Fonksiyon Sınıfları}

$f: (M_1, d_1) \to (M_2, d_2)$ fonksiyonu:

\begin{table}[h]
\centering
\begin{tabular}{ll}
\toprule
\textbf{Sınıf} & \textbf{Koşul} \\
\midrule
Genişlemez (non-expansive) & $d_2(f(x),f(y)) \leq d_1(x,y)$ \\
Lipschitz & $\exists K>0:\; d_2(f(x),f(y)) \leq K \cdot d_1(x,y)$ \\
Üniform sürekli & $\forall\varepsilon>0\;\exists\delta>0:\; d_1(x,y)<\delta \Rightarrow d_2(f(x),f(y))<\varepsilon$ \\
Sürekli & Her noktada $\varepsilon$-$\delta$ sürekliliği \\
İzometri & $d_2(f(x),f(y)) = d_1(x,y)$ \\
Benzerlik & $\exists c>0:\; d_2(f(x),f(y)) = c \cdot d_1(x,y)$ \\
Homeomorfizma & Bijektif sürekli; ters de sürekli \\
\bottomrule
\end{tabular}
\caption{Metrik fonksiyon sınıfları}
\end{table}

\textbf{Sıralama:} İzometri $\Rightarrow$ Genişlemez $\Rightarrow$ Lipschitz
$\Rightarrow$ Üniform sürekli $\Rightarrow$ Sürekli.

\begin{figure}[ht]
\centering
\begin{tikzpicture}[font=\small]
  % --- domain (M1, d1) ---
  \draw[rounded corners=8pt, thick] (-0.4,-1.6) rectangle (3.2,1.8);
  \node at (-0.0,1.45) {$(M_1,d_1)$};
  \fill[blue!70!black] (0.6,0.9) circle (1.8pt) node[left, font=\scriptsize] {$x$};
  \fill[blue!70!black] (2.2,-0.9) circle (1.8pt) node[right, font=\scriptsize] {$y$};
  \draw[blue!70!black, thick] (0.6,0.9) -- node[above right, font=\scriptsize] {$d_1(x,y)$} (2.2,-0.9);
  % --- codomain (M2, d2) ---
  \draw[rounded corners=8pt, thick] (5.6,-1.6) rectangle (9.4,1.8);
  \node at (9.0,1.45) {$(M_2,d_2)$};
  \fill[violet!80!black] (6.6,0.9) circle (1.8pt) node[left, font=\scriptsize] {$f(x)$};
  \fill[violet!80!black] (8.2,-0.9) circle (1.8pt) node[right, font=\scriptsize] {$f(y)$};
  \draw[violet!80!black, thick] (6.6,0.9) -- node[above right, font=\scriptsize] {$d_2(f x,f y)$} (8.2,-0.9);
  % --- arrow f ---
  \draw[-{Stealth}, thick] (3.45,0.1) to[bend left=10] node[above, pos=0.5] {$f$} (5.35,0.1);
  % --- equality note ---
  \node[font=\scriptsize, align=center] at (4.4,-2.1)
    {izometri: $d_2(f x, f y) = d_1(x, y)$ (mesafe korunur)};
\end{tikzpicture}

\caption{İzometri: $f$, $x$ ile $y$ arasındaki mesafeyi korur ---
         $d_2(f(x),f(y)) = d_1(x,y)$.}
\end{figure}

\begin{sezgi}
Bir izometriyi, bir kâğıt parçasını \emph{buruşturmadan ve germeden} başka bir
yere taşıyan bir hareket gibi düşünün: döndürebilir, kaydırabilir,
yansıtabilirsiniz, ama iki nokta arasındaki mesafe asla değişmez. Düzlemde
öteleme, dönme ve yansıma izometridir. Benzerlik ise kâğıdı oranını koruyarak
büyüten/küçülten bir fotokopi makinesidir ($c=2 \Rightarrow$ her mesafe iki
katına çıkar); $c=1$ alındığında benzerlik tekrar izometriye iner.
\end{sezgi}

\subsection{Lipschitz Sabiti}

\[
K = \sup_{x \neq y} \frac{d_2(f(x),f(y))}{d_1(x,y)}.
\]

İzometri: $K=1$ ve eşitlik; benzerlik: $K=c$ sabit.

\begin{figure}[ht]
\centering
\begin{tikzpicture}[font=\small]
  % --- domain (M1, d1) ---
  \draw[rounded corners=8pt, thick] (-0.4,-1.7) rectangle (3.2,1.9);
  \node at (-0.0,1.55) {$(M_1,d_1)$};
  \fill[blue!70!black] (0.5,1.05) circle (1.8pt) node[left, font=\scriptsize] {$x$};
  \fill[blue!70!black] (2.4,-1.05) circle (1.8pt) node[right, font=\scriptsize] {$y$};
  \draw[blue!70!black, very thick] (0.5,1.05) -- node[above right, font=\scriptsize] {$d_1(x,y)$} (2.4,-1.05);
  % --- codomain (M2, d2): noktalar birbirine cok daha yakin ---
  \draw[rounded corners=8pt, thick] (5.6,-1.7) rectangle (9.4,1.9);
  \node at (9.0,1.55) {$(M_2,d_2)$};
  \fill[violet!80!black] (7.2,0.45) circle (1.8pt) node[above left, font=\scriptsize] {$f(x)$};
  \fill[violet!80!black] (7.9,-0.25) circle (1.8pt) node[below right, font=\scriptsize] {$f(y)$};
  \draw[violet!80!black, very thick] (7.2,0.45) -- node[above right, font=\scriptsize] {$\le k\,d_1(x,y)$} (7.9,-0.25);
  % --- arrow f ---
  \draw[-{Stealth}, thick] (3.45,0.05) to[bend left=10] node[above, pos=0.5] {$f$} (5.35,0.05);
  % --- contraction note ---
  \node[font=\scriptsize, align=center] at (4.4,-2.2)
    {büzülme: $d_2(f x, f y) \le k\, d_1(x,y),\; k<1$ (mesafe daralır)};
\end{tikzpicture}

\caption{Büzülme: $k<1$ olduğunda $f$ noktaları birbirine yaklaştırır ---
         $d_2(f(x),f(y)) \le k\,d_1(x,y)$.}
\end{figure}

\begin{sezgi}
Lipschitz sabiti $K$, fonksiyonun ``en dik eğimi'' gibidir: hiçbir nokta
çiftinde mesafeyi $K$ katından fazla büyütemez. $K<1$ olduğunda her adımda mesafe
kesin olarak küçülür --- buna \textbf{büzülme (contraction)} denir. Büzülmeyi
tekrar tekrar uyguladığınızda tüm noktalar tek bir noktaya (sabit noktaya) doğru
hunilenir; Banach sabit-nokta teoreminin kalbi budur.
\end{sezgi}

\begin{figure}[ht]
\centering
\begin{tikzpicture}[font=\small]
  % eksenler
  \draw[-{Stealth[length=2.2mm]}] (-0.3,0) -- (9.3,0) node[right] {$x$};
  \draw[-{Stealth[length=2.2mm]}] (0,-0.3) -- (0,3.4) node[above] {$f(x)$};
  % egim sinirli (Lipschitz) bir egri: tum noktalarda ayni delta is gorur
  \draw[blue!70!black, very thick] (0.2,0.4) .. controls (3.0,1.0) and (5.5,2.2) .. (9.0,2.9);
  % iki ornek nokta + ortak delta bandi
  \def\dx{0.9}
  % nokta A
  \draw[gray!55,dashed] (2.0,0) -- (2.0,1.05);
  \fill[red!75!black] (2.0,1.05) circle (1.6pt);
  \draw[orange!75!black,thick] (2.0-\dx,-0.18) -- (2.0+\dx,-0.18)
       node[midway,below,font=\scriptsize]{$\delta$};
  % nokta B (cok daha saglarda) -- AYNI delta
  \draw[gray!55,dashed] (7.0,0) -- (7.0,2.55);
  \fill[red!75!black] (7.0,2.55) circle (1.6pt);
  \draw[orange!75!black,thick] (7.0-\dx,-0.18) -- (7.0+\dx,-0.18)
       node[midway,below,font=\scriptsize]{$\delta$};
  % aciklama
  \node[font=\scriptsize, align=center, blue!70!black] at (5.0,3.25)
    {tüm noktalar için \emph{ortak} $\delta$ ($\varepsilon$ verildi)};
  \node[font=\scriptsize, align=center] at (4.6,-0.95)
    {düzgün süreklilik: $\delta$ noktadan bağımsız seçilebilir};
\end{tikzpicture}

\caption{Düzgün süreklilik: $\varepsilon$ verildiğinde aynı $\delta$ tüm
         noktalarda iş görür (noktadan bağımsız).}
\end{figure}

\begin{karsiornek}
Her sürekli fonksiyon düzgün (üniform) sürekli \emph{değildir}. $\mathbb{R}^+$
üzerinde $f(x) = 1/x$ süreklidir, ama düzgün sürekli değildir: $x$ sıfıra
yaklaştıkça eğim sınırsız diklenir, bu yüzden verilen bir $\varepsilon$ için
\emph{tek} bir $\delta$ bütün noktalarda iş göremez ($x \approx 0.001$ civarında
çok küçük bir $\delta$ gerekirken $x \approx 10$ civarında çok daha büyüğü
yeterlidir). Aynı şekilde $\mathbb{R}$ üzerinde $f(x)=x^2$ süreklidir ama düzgün
sürekli değildir. Düzgün süreklilik $\delta$'nın \emph{noktadan bağımsız}
seçilebilmesini ister; Lipschitz fonksiyonlar bunu $\delta=\varepsilon/K$ ile
otomatik sağlar.
\end{karsiornek}

%-------------------------------------------------------------------
\section{Teoremler}

\begin{theorem}
İzometri $\Rightarrow$ homeomorfizma.
\begin{proof}[İspat eskizi]
Önce \textbf{injektiflik}: $f$ bir izometri ve $f(x)=f(y)$ olsun. O zaman
$d_1(x,y) = d_2(f(x),f(y)) = 0$, dolayısıyla metriğin ayırma aksiyomundan $x=y$:
izometri daima injektiftir. Bijektif kabul edersek $f^{-1}$ vardır ve o da bir
izometridir. İzometri her mesafeyi koruduğu için $\varepsilon$-$\delta$
sürekliliğini $\delta=\varepsilon$ ile sağlar: hem $f$ hem $f^{-1}$ süreklidir.
Bijektif + sürekli + sürekli ters $=$ homeomorfizma.
\end{proof}
\end{theorem}

\begin{theorem}
Benzerlik, $c=1 \Rightarrow$ İzometri.
\end{theorem}

\begin{theorem}
Lipschitz $\Rightarrow$ Üniform sürekli $\Rightarrow$ Sürekli.
\begin{proof}[İspat eskizi]
$f$ sabiti $K>0$ olan Lipschitz olsun. Verilen $\varepsilon>0$ için
$\delta = \varepsilon/K$ seçin; $d_1(x,y)<\delta$ ise
$d_2(f(x),f(y)) \le K\,d_1(x,y) < K(\varepsilon/K) = \varepsilon$. Seçilen
$\delta$ noktaya bağlı değildir, dolayısıyla $f$ düzgün süreklidir; bu özel
olarak sürekliliği de içerir.
\end{proof}
\end{theorem}

\begin{theorem}
Kompakt metrik uzaydan metrik uzaya her sürekli fonksiyon üniform süreklidir.
\end{theorem}

%-------------------------------------------------------------------
\section{Algoritmalar}

\subsection{classify\_finite\_metric\_map — $O(|X|^2)$}

\begin{algorithm}
\caption{ClassifyMap$(X, d_1, Y, d_2, f)$}
\begin{algorithmic}[1]
\State $K \leftarrow 0$; isometry $\leftarrow$ True; similarity\_ratio $\leftarrow$ None
\For{each pair $(x_1, x_2)$ with $x_1 \neq x_2$}
    \State ratio $\leftarrow d_2(f(x_1),f(x_2)) / d_1(x_1,x_2)$
    \State $K \leftarrow \max(K, \text{ratio})$
    \If{$d_2(f(x_1),f(x_2)) \neq d_1(x_1,x_2)$}
        \State isometry $\leftarrow$ False
    \EndIf
    \If{similarity\_ratio is None}
        \State similarity\_ratio $\leftarrow$ ratio
    \ElsIf{ratio $\neq$ similarity\_ratio}
        \State similarity\_ratio $\leftarrow$ None
    \EndIf
\EndFor
\State \Return MetricMapProfile(lipschitz\_constant=$K$, isometry=isometry, $\ldots$)
\end{algorithmic}
\end{algorithm}

Karmaşıklık: $O(|X|^2)$.

%-------------------------------------------------------------------
\section{pytop API}

\begin{lstlisting}[language=Python]
from pytop.metric_spaces import FiniteMetricSpace
from pytop.metric_map_taxonomy import (
    classify_finite_metric_map,
    metric_map_profile,
    MetricMapProfile,
)
\end{lstlisting}

\texttt{classify\_finite\_metric\_map(fms1, fms2, f)} $\to$ \texttt{MetricMapProfile} döner (Result değil).

\texttt{MetricMapProfile} alanları: \texttt{isometry}, \texttt{similarity}, \texttt{similarity\_ratio},
\texttt{lipschitz\_constant}, \texttt{non\_expansive}, \texttt{bijective}, \texttt{lipschitz},
\texttt{continuous}, \texttt{uniformly\_continuous}, \texttt{homeomorphism}, \texttt{name}, \texttt{certification}.

%-------------------------------------------------------------------
\section{Örnekler}

\subsection*{Örnek 5.1 — Özdeşlik: İzometri}

\begin{lstlisting}[language=Python]
f_id = {1: 1, 2: 2, 3: 3, 4: 4}
prof_id = classify_finite_metric_map(fms, fms, f_id)
print("isometry:", prof_id.isometry)
print("lipschitz_constant:", prof_id.lipschitz_constant)
\end{lstlisting}

\begin{verbatim}
isometry: True
lipschitz_constant: 1.0
\end{verbatim}

\subsection*{Örnek 5.2 — 2× Ölçekleme: Benzerlik}

\begin{lstlisting}[language=Python]
f_scale = {1: 2, 2: 4, 3: 6, 4: 8}
prof_scale = classify_finite_metric_map(fms_line1, fms_line2, f_scale)
print("similarity:", prof_scale.similarity)
print("similarity_ratio:", prof_scale.similarity_ratio)
print("lipschitz_constant:", prof_scale.lipschitz_constant)
\end{lstlisting}

\begin{verbatim}
similarity: True
similarity_ratio: 2.0
lipschitz_constant: 2.0
\end{verbatim}

\subsection*{Örnek 5.3 — Sabit Fonksiyon: K=0}

\begin{lstlisting}[language=Python]
f_const = {1: 1, 2: 1, 3: 1, 4: 1}
prof_const = classify_finite_metric_map(fms, fms, f_const)
print("lipschitz_constant:", prof_const.lipschitz_constant)
print("bijective:", prof_const.bijective)
\end{lstlisting}

\begin{verbatim}
lipschitz_constant: 0.0
bijective: False
\end{verbatim}

\subsection*{Örnek 5.6 — Büzülme (Contraction): $K = 1/2$}

Öklid doğrusunda her mesafeyi yarıya indiren bir eşleme bir \textbf{büzülmedir}
($K<1$): Banach teoreminin sözleşme koşulunu sağlar.

\begin{lstlisting}[language=Python]
src = [0, 2, 4, 6]
img = [0, 1, 2, 3]
d_src = {(a, b): abs(a - b) for a in src for b in src}
d_img = {(a, b): abs(a - b) for a in img for b in img}
M_src = FiniteMetricSpace(carrier=tuple(src), distance=d_src)
M_img = FiniteMetricSpace(carrier=tuple(img), distance=d_img)
f_half = {0: 0, 2: 1, 4: 2, 6: 3}
prof_half = classify_finite_metric_map(M_src, M_img, f_half)
print("lipschitz_constant:", prof_half.lipschitz_constant)
print("similarity_ratio:", prof_half.similarity_ratio)
print("isometry:", prof_half.isometry)
\end{lstlisting}

\begin{verbatim}
lipschitz_constant: 0.5
similarity_ratio: 0.5
isometry: False
\end{verbatim}

\subsection*{Örnek 5.7 — Ayrık Metrikte Döngü: İzometri ve Homeomorfizma}

Ayrık metrikte her permütasyon mesafeyi korur: bir izometri ve (bijektif
olduğundan) homeomorfizmadır.

\begin{lstlisting}[language=Python]
dp = [1, 2, 3]
d_disc = {(a, b): (0 if a == b else 1) for a in dp for b in dp}
M_disc = FiniteMetricSpace(carrier=tuple(dp), distance=d_disc)
f_cycle = {1: 2, 2: 3, 3: 1}
prof_cycle = classify_finite_metric_map(M_disc, M_disc, f_cycle)
print("isometry:", prof_cycle.isometry)
print("homeomorphism:", prof_cycle.homeomorphism)
\end{lstlisting}

\begin{verbatim}
isometry: True
homeomorphism: True
\end{verbatim}

%-------------------------------------------------------------------
\section{Alıştırmalar}

\subsection*{Kodlama}

\begin{enumerate}[label=K\arabic*.]
\item $\{1,2,3\}$ üzerinde $f(1)=2, f(2)=3, f(3)=1$ (döngü) fonksiyonunu ayrık
      metrikte \texttt{classify\_finite\_metric\_map} ile sınıflandırın.
\item Öklid metrikli $\{0,1,2,3,4\}$ üzerinde $f(x) = x+1$ kaydırma fonksiyonunun
      Lipschitz sabitini hesaplayın.
\item Sabit fonksiyon $f(x) = 1$ için $K=0$ olduğunu doğrulayın.
\item Öklid doğrusunda $\{0,2,4,6\} \to \{0,1,2,3\}$ büzülmesini ($f(x)=x/2$) kurun;
      \texttt{lipschitz\_constant}, \texttt{similarity\_ratio}, \texttt{non\_expansive}
      alanlarını yazdırın ve $K<1$ olduğunu doğrulayın.
\item $\{1,2,3\}$ üzerinde özdeşlik olmayan bir döngüyü hem ayrık hem Öklid
      metrikte sınıflandırın; hangi metrikte izometri olduğunu \texttt{isometry}
      alanlarıyla karşılaştırın.
\end{enumerate}

\subsection*{Teori}

\begin{enumerate}[label=T\arabic*.]
\item İzometri $\Rightarrow$ genişlemez $\Rightarrow$ Lipschitz ($K=1$) zincirini ispatlayın.
\item Lipschitz $\Rightarrow$ üniform sürekli olduğunu $\delta = \varepsilon/K$ seçimi
      ile ispatlayın.
\item İzometrinin daima injektif olduğunu ispatlayın ($f(x)=f(y) \Rightarrow
      d_2(f(x),f(y))=0$); ardından düzlemde sabit bir fonksiyonun neden izometri
      olamadığını açıklayın.
\end{enumerate}


\backmatter

\chapter*{Dizin}
\addcontentsline{toc}{chapter}{Dizin}

\begin{description}
  \item[Açık top] Bkz.\ Bölüm 14.
  \item[Banach sabit-nokta teoremi] Bkz.\ Bölüm 15.
  \item[Bağlantılı uzay] Bkz.\ Bölüm 8.
  \item[Cauchy dizisi] Bkz.\ Bölüm 15.
  \item[Heine-Borel teoremi] Bkz.\ Bölümler 7, 15.
  \item[Homeomorfizma] Bkz.\ Bölümler 4, 16.
  \item[İzometri] Bkz.\ Bölüm 16.
  \item[Kompaktlık] Bkz.\ Bölüm 7.
  \item[Lipschitz sabit] Bkz.\ Bölüm 16.
  \item[Metrik aksiyomları] Bkz.\ Bölüm 14.
  \item[Sayılabilirlik aksiyomları] Bkz.\ Bölüm 9.
  \item[Tamlık] Bkz.\ Bölüm 15.
  \item[Topoloji aksiyomları] Bkz.\ Bölüm 4.
  \item[Urysohn fonksiyon teoremi] Bkz.\ Bölüm 6.
\end{description}

\end{document}
